% 本文件由 LaTeX 排版 Agent 根据有效 translation.json 自动生成。
% author=John Chrysostom; work_id=2306; title=弟茂德前书讲道集

\chapter{弟茂德前书讲道集}

% structure: p0001 source=h1 target=omitted reason=page-title-already-rendered-by-parent

序言\newline{}讲道 1\newline{}讲道 2\newline{}讲道 3\newline{}讲道 4\newline{}讲道 5\newline{}讲道 6\newline{}讲道 7\newline{}讲道 8\newline{}讲道 9\newline{}讲道 10\newline{}讲道 11\newline{}讲道 12\newline{}讲道 13\newline{}讲道 14\newline{}讲道 15\newline{}讲道 16\newline{}讲道 17\newline{}讲道 18\par

\SourceRecord{Translated by Philip Schaff.；Nicene and Post-Nicene Fathers, First Series；Vol. 13.；Edited by Philip Schaff.；Buffalo, NY: Christian Literature Publishing Co.,；1889.；<http://www.newadvent.org/fathers/2306.htm>.}

% structure: p0001 source=h1 target=section reason=page-title

\section{论\BibleBook{弟茂德}{前书}讲道}

1. 弟茂德也是宗徒保禄的门徒之一。路加为这位青年的非凡品质作证，告诉我们，他\BibleQuote{被吕斯特辣及依科尼雍的弟兄们所称赞}\BibleRef{宗 16:2} 他立即成了门徒和教师，并表现出独一无二的明智：他听到保禄宣讲时不坚持割损礼，并且明白保禄曾在这事上与伯多禄对抗，他不仅不反对宣讲，反而服从了这礼规。因为据说保禄\BibleQuote{带他去，给他行了割损礼}\BibleRef{宗 16:3}，尽管他已成年，却将自己的全部计划托付给了他。\par

保禄对他的爱足以证明他的品格。因为他在别处论到他说：\BibleQuote{你们知道他的贤能，他怎样如同儿子同父亲一样，与我共同为福音服务}\BibleRef{斐 2:22} 又在致格林多人书中写道：\BibleQuote{我打发弟茂德到你们那里去，他是我的爱子，在主内忠信的人}\BibleRef{格前 4:17} 又说：\BibleQuote{谁也不可轻视他，因为他做主的工，如同我一样}\BibleRef{格前 16:10} 在致希伯来人书中，他写道：\BibleQuote{你们要知道，我们的兄弟弟茂德已得释放}\BibleRef{希 13:23} 的确，他对他的爱处处可见，现今所行的奇迹仍然证明他的价值。\par

2. 如果问，为什么他只致书给弟铎和弟茂德，虽然息拉也被认可，路加也是如此——他写道：\BibleQuote{只有路加同我在一起}\BibleRef{弟后 4:11}，克肋孟也是他的同工，他说\BibleQuote{与克肋孟以及其他同工}\BibleRef{斐 4:3}——那么，为什么他只写信给弟铎和弟茂德呢？这是因为他已经把照顾教会的责任托付给了他们，某些明确的地方已分配给他们，而其他人则跟随他。因为弟茂德在德行上如此卓越，以致他的年轻并不妨碍他晋升；因此他写道：\BibleQuote{不要让人小看你年轻}\BibleRef{弟前 4:12}；又说：\Quote{待年轻的妇女如同姐妹。}因为哪里有了德行，其他一切都是多余的，就不会有任何阻碍。因此，当宗徒论及主教时，在他所要求的许多事项中，没有特别提及年龄。如果他说主教应是\BibleQuote{一个妻子的丈夫}，并\BibleQuote{使儿女顺服}\BibleRef{弟前 3:2}，这并非说必须要有妻子和儿女，而是说，如果有人碰巧从世俗生活被提升到那职位，他们应当知道如何管理自己的家庭和儿女，以及所有托付给他们的人。因为如果一个人既是世俗的，又在这些方面有欠缺，怎么能把教会的事务托付给他呢？\par

3. 但你会说，为什么他要致书给一位已被任命为教师职务的门徒呢？难道他不应该在受派遣之前就为职务做好准备吗？是的；但他所需要的教导不是适合门徒的，而是适合教师的。因此你会在整封信中看到他调整自己的教导以适应教师。这样，在开头他并没有说：\Quote{不要理会那些讲别样道理的人}，而是说：\BibleQuote{嘱咐他们不可教导别的教义}\BibleRef{弟前 1:3}\par

\SourceRecord{Translated by Philip Schaff.；Nicene and Post-Nicene Fathers, First Series；Vol. 13.；Edited by Philip Schaff.；Buffalo, NY: Christian Literature Publishing Co.,；1889.；<http://www.newadvent.org/fathers/230600.htm>.}

% structure: p0001 source=h1 target=section reason=page-title

\section{第一篇讲道：弟茂德前书}

\BibleRef{弟前 1:1-2}\par

\Emph{\BibleQuote{保禄，因天主我们的救主及基督耶稣我们的希望的命令，作耶稣基督的宗徒；致书给因信德作我真子的弟茂德：愿恩宠、仁慈与平安，由天主父及我们的主基督耶稣赐与你。} [修订标准本省略了\OriginalTerm{}{\GreekText{κυρίου}}，而译为\InnerQuote{基督耶稣我们的希望}，\OriginalTerm{}{\GreekText{τῆς} \GreekText{ἐλπίδος} \GreekText{ἡμῶν}}。]}\par

宗徒的尊位是伟大而可敬佩的，我们看见保禄不断地陈明其原因，似乎不是他自取自荣，而是受托付，且不得不如此。因为当他说自己\Quote{蒙召}，且是\Quote{因天主的旨意}，又在别处说\Quote{我却有不得已的责任}\BibleRef{格前 9:16}，并说\Quote{我为此被分别出来}，这样的说法除去了一切骄傲和野心的想法。因为那擅自进入不是天主所赐的职分者，应受最严厉的谴责；同样，那接受了职分而拒绝退缩的，也犯了另一种罪，即悖逆和不顺服。因此保禄在这书信的开端如此说：\Quote{保禄，因天主命令作耶稣基督的宗徒。}他在这里没有说\Quote{保禄蒙召}，而是说\Quote{因命令}。他如此开始，为免弟茂德以为保禄像对门徒一样对他说话，而感到人性的软弱。但这命令是在哪里下达的呢？我们在宗徒大事录中读到：\Quote{圣神说：你们为我选拔保禄和巴尔纳伯。}\BibleRef{宗 13:2} 保禄在著作中处处加上宗徒之名，为教训听众，不要以为他所传授的教训是从人来的。因为宗徒不能凭自己说什么，而称自己为宗徒，立即将听众引向那差遣他的。因此他在所有书信的开端都采用这称号，从而赋予他的话权威，如这里他说：\Quote{保禄，因天主我们的救主的命令，作耶稣基督的宗徒。}但似乎并未见到圣父在何处命令他。处处是基督对他说话。因此\Quote{祂对我说：\InnerQuote{离开，因为我要打发你到远方外邦人那里去。}}\BibleRef{宗 22:21}；又说：\Quote{你必须被带到凯撒面前。}\BibleRef{宗 27:24} 但圣子所命令的，他认为就是圣父的命令，正如圣神的命令就是圣子的命令。因为他是由圣神所派遣，由圣神所分别出来，而他说这是天主的命令。这是否有损圣子的权能呢？因为祂的宗徒是由圣父的命令差遣的？绝无此事。因为请看，他如何表明权能是两者共有的。因为说了\Quote{因天主我们的救主的命令}，他接着加上\Quote{及主耶稣基督，我们的希望}。请看，他如何恰当地应用这些名号。确实，圣咏作者将这名号归于圣父，说：\Quote{大地四极的希望。}\BibleRef{咏 63:5} 再者，有福的保禄在另一处写道：\Quote{因此我们劳苦，并且受辱，因为我们指望永生的天主。}师傅必须比门徒承受更多的危险。因为\Quote{我必打击牧人，（他说，）羊群就四散。}\BibleRef{匝 13:7} 因此魔鬼对师傅更猛烈地发怒，因为若他们被毁，羊群也被驱散。因为杀害羊，他只是减少了羊群；但当除去了牧人，他便毁灭了整个羊群，因此他更猛烈地攻击牧人，以较小的努力造成更大的祸害，并且在一个灵魂中造成全体的毁灭。为此，保禄在开始时提振并鼓励弟茂德的灵魂，说：我们有天主作我们的救主，基督作我们的希望。我们受许多苦，但我们的希望是大的；我们暴露于陷阱和危险中，但拯救我们的不是人，而是天主。我们的救主并不软弱，因为祂是天主，无论什么危险都不能胜过我们；我们的希望也不致蒙羞，因为它是基督。因为我们能从两方面承受危险：要么迅速脱离，要么在危险中因美好的希望而得支持。\par

但是保禄从未称自己是圣父的宗徒，而总是基督的宗徒。因为他使一切为两者所共有。他甚至称福音本身为\Quote{天主的福音}。并且他暗示，我们在此所受的一切，现在的事都算不得什么。\par

\BibleQuote{致书给因信德作我真子的弟茂德。}\par

这也令人鼓舞。因为若他表现出如此信德，以致被特别称为\Person{保禄}的\Quote{亲生}儿子，那么对于未来他也当满怀信心。因为信德的特质就是即使发生似乎与恩许相悖的事，也不沮丧或不安。但请注意，他说\Quote{我的儿子}，甚至\Quote{我亲生的儿子}，然而他们并非同一本体。但那又如何？难道他是无理性的种类吗？有人会说：\Quote{好吧，}有人会说，\Quote{他不是出于保禄，所以这不表示出于另一个人。}那么又如何？难道他是另一种本体吗？也不是这样，因为在说了\Quote{我亲生的儿子}之后，他又加上\Quote{在信德中}，以表明他确实是\Quote{他亲生的儿子}，真正出于他。没有区别。他与保禄的相似在于信德，如同人的生育在本体方面的相似。在人类中，儿子像父亲，但在天主方面，相近性更大。因为在这里，父亲和儿子虽同一本体，却在许多细节上不同，如肤色、形象、悟性、年龄、心性、灵魂和身体的禀赋等，在许多方面他们可能相似或不相似，但在神圣本质中没有这种不相似。\Quote{按命令}。这比\Quote{称作}更有力，正如我们从其他经文所知。正如他在这里称\Person{弟茂德}为\Quote{我亲生的儿子}，同样他对\Place{格林多}人说：\Quote{在基督耶稣内我生了你们}，即在信德内；但他加上\Quote{亲生}一词，以显示他与自己的特殊相似，以及他对他的爱和深情。再次注意\Quote{在}应用于信德。\Quote{我亲生的儿子}，他说，\Quote{在信德内}。看何等光荣的分别，他不仅称他为\Quote{儿子}，而且称他为\Quote{亲生的}儿子。\par

第二节。\BibleQuote{恩宠、怜悯与平安，由天主我们的父和耶稣基督我们的主赐予。}见：\BibleRef{弟前 1:2}\par

为何在其他书信中没有提及怜悯，而这里却提到了？这是他深情的进一步标志。他为自己儿子祈求更大的祝福，带着父母般的忧虑。因为他的忧虑如此之深，以致给了\Person{弟茂德}一些在别处未曾给过的指示，要他照顾自己的身体；在那里他说：\BibleQuote{因你的胃和屡次的不适，稍微用点酒吧。}见：\BibleRef{弟前 5:23} 教师们确实更需要怜悯。\par

他说：\Quote{由天主我们的父，}以及\Quote{和耶稣基督我们的主。}\par

这里也有安慰。因为如果天主是我们的父，祂就照看我们如同儿子，正如基督所说：\BibleQuote{你们中间有哪个人，儿子求饼，反而给他石头呢？}见：\BibleRef{玛 7:9}\par

第三节。\BibleQuote{当我往\Place{马其顿}去的时候，曾劝你仍留在\Place{厄弗所}。}见：\BibleRef{弟前 1:3}\par

请注意他措辞的温和，更像仆人而非主人。因为他不是说\Quote{我命令}，或\Quote{我吩咐}甚或\Quote{我劝告}，而是\Quote{我曾劝你}。但这种语气并非对所有人适用，只有温良而诚实的门徒才可如此对待。腐化而不诚实的人则须以另一种方式对待，正如\Person{保禄}自己在别处指示的：\BibleQuote{以全权责备他们。}见：\BibleRef{铎 2:15}；而在这里他说\Quote{命令}，不是\Quote{劝告}，而是\Quote{命令某些人不要传授别的道理}。这是什么意思？难道\Person{保禄}寄给他们的书信不够吗？不，是足够的；但人们有时会忽视书信，或者这封信可能是在这些书信写成之前。他本人曾在那城住过一段时间。有狄安娜神庙，并在那里遭受了巨大的苦难。因为当戏园里的集会解散后，他召来门徒劝勉了他们，然后发现必须乘船离开，虽然后来又返回他们那里。值得探究的是，他是否当时就派\Person{弟茂德}驻扎在那里。因为他说\Quote{你要命令某些人不要传授别的道理}；他没有点名，以免公开责备使他们更加无耻。在那城中，有些犹太人的假宗徒，他们想要迫使信徒遵守犹太法律，这是他在书信中到处指出的一个错误；他们这样做并非出于良心，而是出于虚荣心，想拥有门徒，出于对蒙福的\Person{保禄}的嫉妒和对抗精神。这就是所谓\Quote{传授别的道理}。\par

第四节。\BibleQuote{不可听从荒谬无凭的传说和无穷的族谱。}见：\BibleRef{弟前 1:4}\par

所谓\Quote{神话}，他指的不是法律；决不是；而是杜撰、伪造和假冒的道理。因为犹太人似乎把他们的全部谈论都浪费在这些无益之点上。他们数算他们的父祖，以便博得历史和考据的名声。他说：\Quote{好叫你吩咐某些人不要讲授别的道理，也不要听从神话和无穷的族谱，这些事只会引起辩论，而不是天主的救恩计划藉信德。}他为什么称它们为\Quote{无穷}？因为它们没有终点，或没有用处，或不易为我们理解。请注意他如何不赞成辩论。因为哪里有信德，就不需要辩论。哪里没有好奇的余地，问题就是多余的。辩论是信德的颠覆。因为寻求的人还没有找到。辩论的人无法相信。因此他的忠告是，我们不应忙于辩论，因为如果我们辩论，那便不是信德；信德使理性安息。但为何基督却说：\BibleQuote{你们寻求，就会找到；你们敲，就会给你们开门}\BibleRef{玛 7:7}；以及\BibleQuote{你们查考圣经，因为你们认为其中有永生}\BibleRef{若 5:39}？那里的寻求是指祈祷和热切渴望，他吩咐\Quote{查考圣经}，不是为了引入辩论的辛劳，而是为了结束它们，好使我们能够确定和稳固它们的真义，而不是一直辩论，而是了结它。他公正地说：\Quote{天主的救恩计划}。因为天主愿意施予的恩惠是伟大的；但它们的伟大不是通过理性所能领悟的。因此这必须是信德的工作，信德是我们灵魂的最佳良药。因此这种辩论与天主的救恩计划相悖。因为藉信德所施予的是什么？领受他的仁慈，成为更好的人；对任何事情都不怀疑、不争辩；而是安心信靠。因为那些\Quote{引起辩论}的东西取代了信德以及信德所成就和建立的。基督说过我们必须因信德得救；这些教师质疑甚至否认这一点。因为既然宣告已经临在，但其结果尚未到来，就需要信德。但他们被法律规条所占据，阻碍了信德。他在这里似乎也暗指希腊人，当他说\Quote{神话和族谱}时，因为他们列举他们的诸神。\par

\Emph{道德教训}。因此，我们不要听从辩论。因为我们被称为信友，正是为了毫不迟疑地相信所传授给我们的一切，而不存任何怀疑。因为如果所断言的是出于人，我们就应当审查它们；但既然出于天主，就只应敬畏和相信。如果我们不信，怎能确信有天主的存在？因为当你向他质问时，你怎能知道有天主？对天主的认识，最明显地表现为在没有证明和论证的情况下相信他。连希腊人都知道这一点；因为有人曾说，他们相信他们的神明告诉他们的，甚至没有证明；而且什么？——他们是神明的后代。但我何必谈及神明？至于那个人，一个骗子兼巫术师（我指的是\Person{毕达哥拉斯}），他们也照样行事，因为关于他说过：他说过。在他们的神庙上方，有一个沉默的形象，她的手指放在嘴上，压住嘴唇，意味深长地劝诫所有经过的人保持沉默。他们的教义如此神圣，而我们的反而不及吗？甚至还要被嘲笑？这真是极度的疯狂！希腊人的教条确实理应受到质疑。因为它们的性质如此，不过是争辩、理性冲突、疑惑和结论。但我们的教条远非如此。因为人的智慧发明了他们的，而我们的却是藉圣神的恩宠所教导的。他们的教义是疯狂和愚蠢，我们的却是真智慧。在他们那里，既没有老师，也没有学生，所有人都是辩论者。在这里，无论是老师还是学生，每个人都应当向应当学习的人学习，而不怀疑，只服从；不争辩，只相信。因为所有古人都藉着信德得了美好见证，而没有信德，一切都崩溃了。我何必谈论天上的事物？我们将会发现，地上的事物也同样依赖信德。因为没有信德，就没有贸易、契约或任何这类事情。如果在这虚假的事物中，信德如此必要，那么在那些真实的事物中，更当如何呢？\par

因此，让我们追求信德，持守信德，如此我们就能从灵魂中驱逐所有毁灭性的教义，例如那些关乎本命和命运的教义。如果你相信有复活和审判，你就能从心中驱逐所有那些虚假的意见。相信有一位公义的天主，你就不会相信有不公义的本命。相信有一位天主和天主眷顾，你就不会相信有一个本命，它主宰万事。相信有一个惩罚的地方和一个天国，你就不会相信一个剥夺我们自由意志、使我们受必然和强制支配的本命。那就不要播种，不要种植，不要作战，不要从事任何工作！因为无论你愿意与否，事情都会按照本命的进程发生！我们还需要祈祷做什么？如果存在这样的本命，你凭什么配作基督徒？因为那样你就无需负责。生活中的技艺从何而来？它们也来自本命吗？是的，你说，并且一个人注定要通过劳动变得聪明。但你能给我指出一个人没有劳动就学会了技艺吗？你不能。那么他的技能不是来自本命，而是来自劳动。\par

但为什么一个堕落邪恶的人，没有从父亲那里继承，却成了富人，而另一个人，历经无限辛劳，却仍然贫穷？因为他们总是提出这样的问题，总是讨论财富和贫穷，从不考虑恶习和美德。在这个问题上，不要谈论那个，而是向我指出一个人，他努力行善却变成了恶人；或者一个人，没有努力却变成了善人。因为如果命运有任何力量，它的力量应该表现在最重要的事情上：在恶习和美德上，而不是在贫穷和财富上。你又会问：为什么一个人多病，而另一个人健康？为什么一个人受尊敬，另一个人受羞辱？为什么这个人万事顺利，而那个人却只遇到失败和阻碍？放弃\OriginalTerm{出生星位}{nativity}的观念，你就会明白。坚定地相信有天主和\OriginalTerm{天主眷顾}{Providence}，这一切都会明朗。你说：\Quote{但我不能相信有\OriginalTerm{天主眷顾}{Providence}，当有这样的混乱时。我能相信善良的天主把财富赐给淫乱者、腐败和不诚实的人，而不赐给有德的人吗？我怎能相信这个？因为必须有事实来奠定信仰。}那么，这些情况来自一个正义的还是不正义的\OriginalTerm{出生星位}{nativity}？你说：\Quote{不正义的。}那么是谁造了它？你说：\Quote{不是天主，它是\OriginalTerm{无生的}{unbegotten}。}但\OriginalTerm{无生的}{unbegotten}怎么能产生这些？因为它们是矛盾的。\Quote{那么这些绝不是天主的工程。}那么我们要问谁造了地、海、天、季节？你回答：\Quote{出生星位。}那么出生星位在无生命的事物中产生了如此有序与和谐，但在我们这些为此而造的事物中却产生了如此多的混乱？就好像一个人在建造房屋时，小心地把它造得华丽，却不对他的家人稍加考虑。但谁保存了季节的交替？谁制定了自然的常规律法？谁决定了昼夜的行程？这些都超越任何这样的出生星位。但你说：\Quote{这些都是自己生成的。}然而，这样一个井然有序的系统怎么能自行产生呢？\par

你说：\Quote{但富者、健康者、有名望者从哪里来？有些人怎样因贪婪致富，有些人因继承，有些人因暴力？为什么天主允许恶人兴旺？}我们回答：因为按照各人应得的报应并不在此处发生，而是保留到将来。那么，向我指出任何这样的事发生！\Quote{好吧，}你说，\Quote{给我现世，我就不指望将来。}但正是因为你寻求现世，所以你不领受。因为如果当现世的享受不在你手中时，你如此热切地寻求眼前的事物，以至于把它们置于将来之上，那么如果你拥有纯粹的快乐，你会做什么呢？因此天主向你表明这些事是虚无的、无关紧要的；因为它们如果不是无关紧要的，祂就不会把它们赐给这样的人。你会承认，一个人是高是矮，是黑是白，是无关紧要的；同样，一个人是富是贫也是如此。因为告诉我，必要的事物不是平等地赐给所有人吗？例如德能、神恩的分配？如果你正确理解天主的仁慈，你就不会抱怨缺少世俗之物，而你却与别人平等地享有这些最好的恩赐；并且知道平等的分配，你就不会渴望在其他方面优越。就像一个仆人从主人的恩惠中，与同伴平等地享受食物、衣服、住宿以及一切必需品，却因有更长的指甲或更多的头发而自夸；同样，基督徒因那些只暂时享受的事物而自高自大。正是为此，天主把这些事物从我们这里撤去，以熄灭这种疯狂，并把我们的情感从它们转移到\OriginalTerm{天堂}{heaven}。然而我们仍不学智慧。就像一个孩子拥有一个玩具，宁愿它超过必需品，而他的父亲为了勉强引导他走向对他更好的事物，夺走他的玩具；同样，天主把这些事物从我们这里拿走，为引导我们到天堂。如果你问为什么祂允许恶人富有，那是因为他们在祂眼中并不重要。如果义人也富有，那与其说是祂使他们如此，不如说是祂允许这样。现在我们肤浅地说这些，如同对不熟悉圣经的人说的。但如果你们相信并留意天主圣言，我们的讲论就是不必要的，因为圣言会教导你们一切。并且为使你们明白财富、健康、荣耀都算不了什么，我可以向你们指出许多人，当他们可以获取财富时，却不寻求财富；当他们可以享受健康时，却克苦自己的身体；当他们可以上升到荣耀时，却以被轻视为目标。但没有一个善人曾努力要成为恶人。因此让我们停止寻求下面的事物，让我们寻求天上的事物；这样我们就能获得它们，并享受永恒的喜乐，靠我们主耶稣基督的恩宠和慈爱。愿光荣、权能和尊崇归于祂，及圣父和圣神，从今时直到永远。阿们。\par

\SourceRecord{Translated by Philip Schaff.；Nicene and Post-Nicene Fathers, First Series；Vol. 13.；Edited by Philip Schaff.；Buffalo, NY: Christian Literature Publishing Co.,；1889.；<http://www.newadvent.org/fathers/230601.htm>.}

% structure: p0001 source=h1 target=section reason=page-title

\section{论\WorkTitle{弟茂德前书} 第二篇讲道}

\Emph{\BibleRef{弟前 1:5-7}}\par

\BibleQuote{命令的目的就是爱德，这爱德出于纯洁的心、良好的良心和无伪的信德：有些人偏离了这些，便转向虚浮的言论；想要作律法的教师，却不明白自己所说的话，也不肯定自己所主张的。}\par

对人而言，没有什么比轻视友谊、不竭尽全力培养友谊更有害的了；反之，也没有什么比尽力追求友谊更有益的了。基督在以下的话中表明了这一点：\BibleQuote{如果你们中两个人在地上同心合意，无论他们求什么，我在天之父必要成全他们}\newline{}\BibleRef{玛 18:19}；又说：\BibleQuote{由于罪恶的增多，爱德必要冷却。}\newline{}\BibleRef{玛 24:12}。这就是一切异端产生的原因。因为人们不爱自己的弟兄，便嫉妒那些享有盛誉的人，由嫉妒而贪图权势，由贪图权势而引入异端。为此，保禄说了\Quote{好叫你能吩咐某些人不可讲授别的道理}之后，现在指出实现这一点的方法是爱德。因此，正如当他称\BibleQuote{基督是法律的终结}\newline{}\BibleRef{罗 10:4}，即法律的圆满，这与此处相关，同样，这条命令也隐含在爱德内。正如医药的目的是健康，一旦健康，便无需多费周折；同样，哪里有了爱德，哪里便无需许多命令。但他说的是哪一种爱德呢？是真诚的，不仅在于言语，而且出自心性、情感与同情。他说\Quote{出于纯洁的心}，或指正直的生活，或指真诚的情感。因为不纯洁的生活也会造成分裂。\BibleQuote{凡作恶的，都憎恨光明}\newline{}\BibleRef{若 3:20}。的确，甚至恶人之间也有友谊。强盗和杀人犯可以彼此相爱，但这不是\Quote{出于良好的良心}，不是出于\Quote{纯洁的心}，而是出于不纯洁的心，不是出于\Quote{无伪的信德}，而是出于虚伪和假冒的信德。因为信德指明真理，真诚的信德产生爱德，真正信天主的人不能容忍放弃爱德。\par

第6节。\BibleQuote{有些人偏离了这些，便转向虚浮的言论。}\newline{}\BibleRef{弟前 1:6}\par

他说\Quote{偏离}，这很恰当，因为射得正直而不偏离目标，需要技巧，需要有圣神的引导。因为有许多事情会让我们偏离正道，而我们只应注视一个目标。\par

第7节。\BibleQuote{想要作律法的教师。}\newline{}\BibleRef{弟前 1:7}\par

这里我们看到邪恶的另一个原因：贪图权势。因此基督说：\BibleQuote{不要称为拉比}\newline{}\BibleRef{玛 23:8}；宗徒又说：\BibleQuote{因为他们不遵守法律，只愿因你们的肉体而夸耀}\newline{}\BibleRef{迦 6:13}。他的意思是，他们贪图优越，因而忽视真理。\par

\BibleQuote{不明白自己所说的话，也不肯定自己所主张的。}\par

这里他谴责他们，因为他们不知道法律的目的与目标，也不知道它应当具有权威的时期。但若出于无知，为何称为罪呢？因为无知不仅起源于他们想要作律法的教师，而且起源于他们没有持守爱德。不仅如此，他们的无知就来自这些原因。因为当灵魂沉溺于属肉的事物时，其视力的清晰度就会变得模糊；远离爱德，就会陷入好争，而且心灵之眼被蒙蔽。因为被任何现世欲望所占有的人，如同沉醉于情欲，不能成为真理的公正裁判。\par

\BibleQuote{不知道自己所肯定的。}\par

很可能他们谈论法律，并详论其净化仪式和其他肉体的礼规。于是宗徒避免谴责这些，要么因为它们毫无价值，要么至多是属灵事物的影子与预象，而是以更吸引人的方式赞美法律，此处将十诫称为法律，并借此摒弃其余。因为如果连这些诫命也惩罚违犯者，并对我们成为无用，那么其他的就更不用说了。\par

第8、9节。\BibleQuote{但我们知道法律是好的，只要人合法地使用它；也知道法律不是为正义的人设立的。}\newline{}\BibleRef{弟前 1:8-9}\par

法律，他似乎在说，是好的，但又是，不好的。那么是怎么回事呢？如果有人不合法地使用它，它就不是好的吗？不，即使那样它也是好的。但他的意思是：若有人在行动中遵守它；因为这就是此处所指的\Quote{合法地使用它}。但若他们用言语解释它，却在行为上忽视它，那就是不合法地使用它。因为这样的人使用它，却无益于自己。此外还可以提到另一种方式。是什么呢？那就是，法律，如果你正确使用它，它就把你引向基督。因为既然它的目的是使人成义，却未能实现这一点，它就转而将我们交给那能够做到的那一位。还有一种合法地使用法律的方式，就是当我们遵守它，但把它当作一件多余的事。怎么是多余的事呢？正如马勒恰当地使用，不是被那跳跃的马嚼着，而是被那仅为装样子而佩戴它的马所用；同样，那自我约束，却不是因法律的条文所迫的人，就是合法地使用法律。合法地使用法律的人是意识到自己不需要它的人，因为那已经如此有德，以至于遵守法律不是出于对它的恐惧，而是出于德性的原则，他就是合法地、平安地使用它：也就是说，若一个人如此使用它，不是出于畏惧它，而是更将良心的谴责而非将来的惩罚放在眼前。而且，他称那已获得德行的人为义人。因此，不需要受法律教导的人就是合法地使用法律。正如字行是放在儿童面前的；但若有人无需其帮助，凭其他知识就按它们所指示的去做，就显出更高的技巧，是更好的读者；同样，那在法律之上的人，就不受法律的管束。因为他在更高的程度上遵守法律，因为他遵守不是出于恐惧，而是出于德性的倾向；那畏惧惩罚的人遵守法律的方式与那追求赏报的人不同。在法律之下的人做事不像在法律之上的人。因为过一种在法律之上的生活就是合法地使用法律。他合法地使用法律并遵守它，因为他成就了超越法律的事，且不需要它的教导。因为法律大部分是禁止邪恶；但这本身并不使人成为义人，还需要行善事。因此，那些像奴隶一样避免邪恶的人，并不达到法律的标准。因为法律是为处罚违犯而设立的。这样的人确实使用它，但那是为了畏惧它的惩罚。经上说：\Quote{你愿意不怕掌权的吗？你只管行善。}\BibleRef{罗 13:3}，这暗示法律只威胁恶人。但法律对那行为配得冠冕的人有什么用呢？正如医生只对那有伤的人有用，对健康无恙的人无用。\Quote{但为不法及不顺服的人，为不虔诚及犯罪的人。}他也称犹太人为\Quote{不法及不顺服的人}。\Quote{法律（他说）激起忿怒，}就是说，对作恶的人。但对那配得赏报的人呢？\Quote{藉着法律，人得知罪。}\BibleRef{罗 3:20}那么，对义人又如何呢？\Quote{法律不是为义人设立的，}他说。为什么呢？因为他免于法律的刑罚，且无须等待法律教导他当尽的本分，因为内有圣神的恩宠指引他。因为法律赐下，是要人因畏惧其恐吓而受管教。但驯良的马无需马勒，那不需教导的人无需师傅。\par

\Quote{但为不法及不顺服的人，为不虔诚及犯罪的人，为不圣洁及亵渎的人，为弑父及弑母的人。}他不仅停止于泛论罪过，也不仅列举这些，而是逐一列举各种罪，似乎是为使人羞于受法律的管束；并在具体列出一些后，又加上对未提及之罪的指涉，尽管他所列举的已足以使人退避。那么他说这些是论及谁呢？论及犹太人，因为他们曾是\Quote{弑父及弑母的人}；他们是\Quote{亵渎及不圣洁的人}，因为他说的\Quote{不虔诚及犯罪的人}也包括这些；而既然是这样，法律就必须赐给他们。因为他们岂非屡次崇拜偶像？岂非用石头砸死梅瑟？他们的手岂非沾满亲人的血？众先知岂非不断为此控告他们？但对于那些受天上哲学教导的人，这些诫命是多余的。\Quote{为弑父及弑母的人，为杀人犯，为奸淫者，为同性恋者，为拐骗者，为说谎者，为发虚誓者，以及一切违反健全道理的事}；因为他所提及的一切都是腐败灵魂的情欲，因此与健全的道理相悖。\par

第11节。\Quote{这是照着那可称颂之天主的荣耀福音，这福音交托了我。}\BibleRef{弟前 1:11}\par

因此，法律对于证实福音仍是必要的，但对于那些服从它的人却是不必要的。他称福音为\Quote{荣耀的福音}。有些人因福音所受的迫害以及基督的苦难而感到羞耻，因此，为了这些人，也为了其他人，他称之为\Quote{荣耀的福音}，以此表明基督的苦难是我们的荣耀。或许他也展望未来。因为如果我们现在的状态暴露于羞辱和责备，将来却不会如此；而福音属于将来，不属于现在。那么，天使为何说：\BibleQuote{看，我给你们报告一个大喜讯，为你们诞生了一位救主}（\BibleRef{路 2:10}）？因为祂生来就是他们的救主，尽管祂的奇迹并非从祂诞生时就开始。他说：\Quote{照著蒙福的天主的福音}。他所指的荣耀，要么是事奉天主的荣耀，要么是说，即使现在的事物充满了祂的荣耀，但将来的事物更会如此；那时\BibleQuote{祂的仇敌要被放在祂的脚下}（\BibleRef{格前 15:25}），那时再也没有反对，那时义人要看见所有那些\BibleQuote{眼所未见，耳所未闻，人心所未曾想到的}福事。（\BibleRef{格前 2:9}）因为我们的救主说：\BibleQuote{我愿意他们也在那里，同我在一起，为使他们享见祢所赐给我的光荣}（\BibleRef{若 17:24}）。\par

\Emph{伦理方面}。让我们学习这些人是怎样的人，并称他们有福，思考他们那时将享受何等福乐，将参与何等光明与荣耀。这世界的荣耀毫无价值，不能持久；即使它持久，也只是到死为止，之后便完全熄灭。因为经上说：\BibleQuote{他的荣耀不能随他下降}（\BibleRef{咏 48:17}）。对于许多人来说，它甚至不能持续到生命结束。但在那荣耀中，绝无此类之事；它持久不衰，永无止境。因为天主的事物就是如此：持久，超越一切变化或终结。因为那境界的荣耀不是来自外在，而是来自内在。我的意思是，它不在于众多仆役、车辆或昂贵的衣裳；不依赖这些，人本身就被荣耀所包裹。在此世，没有这些东西，人就显得赤身露体。在浴场中，我们看见显赫者、平凡者与卑贱者都一般赤裸。大人物在公共场合常因某些时刻被仆人遗弃而陷入危险。但在那世界，人们随身携带自己的荣耀，圣人们如同天使，无论出现在哪里，都自身带着荣耀。更确切地说，正如太阳无需衣裳，不需外援，无论出现在哪里，它的荣耀立刻照耀；那时也是如此。\par

那么，让我们追求那荣耀，没有什么比它更可敬重；并抛弃世界的荣耀，因为它完全一文不值。\BibleQuote{不要因你的衣服和服饰而夸耀。} \BibleRef{德 11:4} 这是古时给简单人的忠告。的确，舞者、娼妓、戏子所穿戴的袍子比你更华丽、更昂贵。此外，这种夸耀所凭借的东西，若被蛾子咬，就能夺去你对它的享受。你看到这是多么不稳定的事物，即现世生命的荣耀？你以昆虫制造和毁坏的东西为傲。因为据说印度的昆虫吐丝，制成你袍子的细线。但倒要寻求由上而来的衣袍，一种可赞叹的、光亮的衣服，真正的金衣——不是由恶人的手从矿山中挖出的金子，而是品德的产物。让我们穿上不是由穷人或奴隶制作，而是由主亲自编织的袍子。但你说，你的衣服镶有金子！那与你何干？制作它的人，而非穿它的人，才是被赞叹的对象，因为赞叹本当归给制作者。那件衣服在洗衣匠的架子上被撑开，但穿它的人却被赞叹？但木头架子穿着它，它被捆在上面。正如那架子穿着它，却非为使用，同样有些妇女也是如此，为了保养衣服，她们说，让它透风，以免被虫蛀！那么，对一件如此无价值的东西操心，做任何事，为它冒丧失救恩的危险，嘲笑地狱，蔑视天主，忽视饥饿的基督，岂不是极端的愚昧？不要再谈论印度、阿拉伯、波斯的珍贵香料，湿的和干的，香水和香膏，如此昂贵却又无用。女人啊，你为什么把香水倒在充满内在不洁的身体上？为什么浪费在令人厌恶的东西上，如同把香水倒在泥土上，或把香膏滴在砖头上？你若愿意，有一种珍贵的香膏和芬芳，你可以用来膏抹你的灵魂；它不是来自阿拉伯、埃塞俄比亚或波斯，而是来自天上本身；不是用黄金购得，而是靠德性的意志和不虚伪的信德。买这香膏吧，它的香气能充满世界。宗徒们所散发的就是这种香气。\BibleQuote{因为我们是（他说）一种馨香，在某些人身上是致死的，在某些人身上是致生的。} \BibleRef{格后 2:15} 这是什么意思？正如人们所说，猪被香水窒息！但这属灵的芬芳不仅熏陶了宗徒们的身体，也熏陶了他们的衣服；保禄的衣服如此浸透了它，以致能驱逐魔鬼。有什么香叶、肉桂、没药如此甜美或如此有效，如同这香气？因为若它能驱逐魔鬼，还有什么它不能成就？让我们用这香膏来装备自己。圣神的恩宠将透过施舍提供它。我们将散发这种香气，当我们进入另一个世界时。正如在这里，身上洒了甜香的人吸引所有人的注意，无论是在浴室、集会还是任何人群聚集之处，众人都跟随他、注视他；同样，在那个世界，当灵魂带着这种属灵的馨香而来时，众人都起立让路。甚至在这里，魔鬼和一切恶习都害怕接近它，不能忍受它，因为它使它们窒息。那么，让我们不要带着那标志着柔弱的香水，而要带着这标志着刚毅、真正可赞叹、充满神圣信心的香水。这是一种香料，不是地上的产物，而是从品德中生长出来的，不会枯萎，永远绽放。这使拥有它的人获得尊荣。我们在受洗时以此被傅油，那时我们散发着它的芬芳；但我们必须靠日后的照料来保留这香气。古时司祭们被膏油傅油，作为德性的象征，司祭应当在自己周围散发这德性的芬芳。\par

没有什么比罪的恶臭更令人作呕，这使圣咏者说：\BibleQuote{我的创伤溃烂流脓。} \BibleRef{咏 37:5} 因为罪比腐烂本身还要污秽。例如，有什么比奸淫更令人作呕？即使犯罪时不易察觉，但犯罪之后，其恶劣本性、所沾染的不洁、咒诅和可憎之处便会被察觉。一切罪皆如此。在犯罪之前，它带有某种快乐；但犯罪之后，快乐停止并消失，痛苦和羞愧接踵而来。但正义则相反：开始时伴随辛劳，但结束时却是快乐和安息。但即使在今世，正如在一种情况下，罪恶的快乐并非快乐，因为对耻辱和惩罚的期待；在另一种情况下，辛劳因赏报的盼望而不觉为辛劳。酗酒的快乐是什么？是饮酒的可怜满足，甚至几乎不是。因为当不省人事随之而来时，人看不见眼前的一切，比疯子更糟，还有什么享受可言？甚至可以说，奸淫本身也没有快乐。因为当情欲剥夺了灵魂的判断力时，还能有真正的快乐吗？我们也可以说瘙痒是快乐！我称那真正的快乐，是灵魂不受情欲影响、不激动、不被身体压倒时的快乐。因为咬牙切齿、歪眼斜视、激动兴奋到有失体面，有什么快乐可言？它远非愉快，反而人们急于逃脱它，结束后还感到痛苦。但如果它是快乐，他们就会希望不逃脱它，而是继续它。因此，它只有快乐之名。\par

但我们所享受的快乐并非如此；它们真正令人愉快，不使人激动，也不使人燃烧。它们使灵魂自由，并令其欢喜和扩展。保禄的快乐就是如此，当他说：\Quote{我为此欢喜，并且将继续欢喜}；又说：\Quote{你们要常常在主内喜乐。} \BibleRef{斐 1:18} 因为罪恶的快乐伴随着羞耻和定罪；它是在暗中放纵，并伴随着无限的不安。但真正的快乐免除了这一切。那么，让我们追求这快乐，好使我们获得未来的美好事物，因我们的主耶稣基督的恩宠和怜悯，愿光荣归于祂，等等。\par

\SourceRecord{Translated by Philip Schaff.；Nicene and Post-Nicene Fathers, First Series；Vol. 13.；Edited by Philip Schaff.；Buffalo, NY: Christian Literature Publishing Co.,；1889.；<http://www.newadvent.org/fathers/230602.htm>.}

% structure: p0001 source=h1 target=section reason=page-title

\section{第三篇论\BibleBook{弟茂德}{前书}讲道辞}

\Emph{\BibleRef{弟前 1:12-14}}\par

\Quote{ \Emph{我感谢我们的主基督耶稣，祂赐给了我能力，因为祂看我是忠信的，把我置于服务中[修订版：服务于祂]，\OriginalTerm{服务于祂}{\GreekText{εἰς} \GreekText{διακονίαν}}；我从前是亵渎者、迫害者、凌辱者；但我蒙了怜悯，因为我是在不信中无知而行的。我们主的恩宠格外丰盈，连同在基督耶稣内的信德和爱德。} }\par

谦逊的好处是公认的，然而它却不易遇到。有足够的谦逊言论，但谦逊的心却无处可寻。有福的保禄如此培养这一品质，以至于他总在寻找谦逊的诱因。那些自知有伟大功绩的人，若要谦逊，必须与自己激烈斗争。而他也同样容易受到强烈诱惑，他良好的良心像肿块一样使他膨胀。请观察他这里的方法。\Quote{我受托，} 他曾说，\Quote{天主的荣耀福音，那些仍坚持法律的人无权分受；因为它现在与福音对立，它们的差异如此之大，以至于受法律驱使的人还不配分受福音；就如我们说，那些需要惩罚和锁链的人无权被接纳进入哲学家的行列。} 因此，他充满了高傲的思想，使用了宏大的言辞后，立刻贬低自己，并促使他人也这样做。在说了\Quote{福音托付给了他} 之后，为免这话显得出自骄傲，他立刻控制自己，以修正的方式补充道：\Quote{我感谢我们的主基督耶稣，祂赐给了我能力，因为祂看我是忠信的，把我置于服务中。} 我们看到，他处处隐藏自己的功德，将一切归于天主，但只限于不取消自由意志。因为不信者可能会说：如果一切都是出于天主，而我们自己毫无贡献，祂像对待木头石头一样将我们从邪恶转向爱智慧，那么祂为何使保禄成为如此，而不使犹达斯这样？为消除这一异议，注意他措词的谨慎：\Quote{它被托付，} 他说，\Quote{在我的信任中。} 这是他自己的优点和功德，但不完全是自己的；因为他说：\Quote{我感谢基督耶稣，祂赐给了我能力。} 这是天主的部份：然后又是他自己的：\Quote{因为祂看我是忠信的。} 的确是因为他愿意自己成为有用的。\par

第13节。\BibleQuote{我从前是亵渎者、迫害者、凌辱者；但我蒙了怜悯，因为我是在不信中无知而行的。}\BibleRef{弟前 1:13}\par

因此我们看到，他承认自己的部份和天主的部份，并将较大的部份归于天主的眷顾时，他减轻了自己的部份，但只限于如我们之前所说，与自由意志一致。\Quote{祂赐给了我能力} 是什么意思？我告诉你。他肩负的重担如此沉重，需要来自高处的许多帮助。试想，每日暴露在侮辱、嘲笑、陷阱、危险、嘲讽、责难和死亡中，却不疲倦、不滑倒、不后退，反而虽每日被无数箭矢攻击，却勇敢承受，坚定不可动摇。这绝非人力所能致，但亦非单靠神能，而是他的决心所致。因为基督预知他将成为怎样的人，从而选择了他，这在祂开始宣讲前为他所作的见证中是明显的。\Quote{他是为我所选用的器皿，把我的名字带到外邦人和君王面前。} \BibleRef{宗 9:15} 因为正如那些在战争中高举王旗的人需要力量和技巧，以免旗帜落入敌人手中；同样，那些在战争与和平中持守基督之名的人，需要强大的力量，以保护它免受指控者的侵害。承受基督之名，好好持守并背起十字架，确实需要极大的力量。因为凡在行为、言语或思想上做出任何与基督不相称之事的人，并未持守祂的名，也没有基督住在他内。因为持守这名的人，不是在人丛中凯旋，而是穿过诸天，众天使肃然起敬，侍立在他周围，钦佩他。\par

\Quote{我感谢主，祂赐给了我能力。} 请注意他如何甚至为他自己的部份感谢天主。因为他承认自己成为\Quote{所选用的器皿} 是出于祂的恩惠。\Quote{因为天主不偏待人。} 但我感谢祂，因祂\Quote{认为我配得上这一服务。} 因为这是祂看我是忠信的证明。家中的管家不仅因被信任而感谢主人，而且视之为主人认为他比他人更忠信的标志：这里也是如此。然后，请注意他描述自己的先前生活时如何放大天主的仁慈和关爱：\Quote{我从前，} 他说，\Quote{是亵渎者、迫害者、凌辱者。} 而当他说到仍不信的犹太人时，他反而减轻了他们的罪责。\Quote{因为我为他们作证：他们对天主有热心，但不是按着真知识。} \BibleRef{罗 10:2} 但论到自己，他说：\Quote{我从前是亵渎者和迫害者。} 请注意他如何贬低自己！他是如此远离自爱，如此充满谦逊，以至于不满足于称自己为迫害者和亵渎者，反而加重自己的罪责，表明他的罪不止于自身，他不仅是亵渎者，而且在亵渎的疯狂中，迫害那些愿意敬虔的人。\par

\BibleQuote{但我蒙了怜悯，因为我是在不信中无知而行的。}\BibleRef{弟前 1:13}\par

那么，为何其他犹太人没有得到慈悲？因为他们所做的，并非无知，而是明知故犯，清楚知道自己在做什么。为此我们有福音作者的作证：\BibleQuote{犹太人中有许多人信了祂，但因法利塞人而不敢公开承认，因为他们喜爱人的称赞胜过天主的称赞。}\BibleRef{若 12:42} 基督又对他们说：\BibleQuote{你们既然彼此寻求光荣，而不寻求出自唯一天主的光荣，怎能相信我呢？}\BibleRef{若 5:44} 那瞎子的父母\BibleQuote{这样说，是因为怕犹太人，因为犹太人已经议定，若有人认祂为默西亚，就把他逐出会堂。}\BibleRef{若 9:22} 连犹太人也说：\BibleQuote{你们看不到我们毫无进展吗？看，世界都跟祂去了。}\BibleRef{若 12:19} 可见他们对权力的贪恋处处妨碍他们。当他们承认除了天主无人能赦罪时，基督随即行了他们所承认的神性标志，这就不可能出于无知。但那时保禄在哪里？或许有人会说，他正坐在加玛里耳的脚前，并没有参与合谋反对耶稣的人群；因为加玛里耳似乎并非一个野心勃勃的人。那么后来保禄如何与众人同谋呢？他看到这教义正在壮大，即将得势，并被普遍接受。因为在基督在世时，门徒与祂交往，后来与他们的老师交往，但当他们完全分离时，保禄并不像其他犹太人那样出于权力欲行事，而是出于热心。因为他前往大马士革的动机是什么？他以为这教义是有害的，害怕它的宣讲会四处蔓延。但犹太人影响他们行动的不是对大众的关心，而是权力欲。因此他们说：\BibleQuote{罗马人必要来，夺去我们的圣殿和国家。}\BibleRef{若 11:48} 这种让他们不安的恐惧，不就是对人的恐惧吗？但值得探究的是，像保禄这样精通法律的人，怎能无知呢？因为正是他说：\BibleQuote{这福音是天主先前藉着祂的圣人先知在圣经所预许的。}\BibleRef{罗 4:2} 那么，你这对祖先法律热心、在加玛里耳脚前长大的人，怎么会不知道呢？然而那些在湖边、河边度日的人，甚至税吏，都接受了福音，而你这位研究法律的人却在迫害它！正是为此，他谴责自己说：\BibleQuote{我原是宗徒中最小的一个。}\BibleRef{格前 9:9} 正是为此，他承认自己的无知，这无知是由不信产生的。因此他说，他得到了\Quote{慈悲}。那么，他说\Quote{祂以为我忠实}是什么意思？他绝不放弃自己主人的任何权利；他把自己的一切都归于主，不为自己索取什么，也不把属于天主的荣耀归于自己。因此，在另一处我们听到他高呼：\BibleQuote{诸位，你们为什么这样行事？我们也是和你们有同样性情的人。}\BibleRef{宗 14:15} 同样，\Quote{祂以为我忠实}。又说：\BibleQuote{我比他们众人更劳碌，但不是我，而是同我一起的天主的恩宠。}\BibleRef{格前 15:10} 又说：\BibleQuote{是天主在我们内运行，使我们愿意并实行。}\BibleRef{斐 2:13} 因此，他承认自己\Quote{得了慈悲}，就是承认自己应受惩罚，因为慈悲是为了这样的人。而在另一处，他论到犹太人说：\BibleQuote{以色列的一部分硬心。}\BibleRef{罗 11:25}\par

第14节。\Quote{并且我们的主的恩宠极其丰盛，连同在基督耶稣内的信德与爱德。}\BibleRef{弟前 1:14}\par

这里加上这句话，以免我们听到他得了慈悲，就仅仅理解为：他本是应受惩罚的、迫害者和亵渎者，却未被惩罚。但慈悲并不限于此，即未施惩罚；它还隐含着许多其他大恩。因为天主不仅免除了我们即将面临的惩罚，还使我们成为\Quote{义人}、\Quote{儿子}、\Quote{弟兄}、\Quote{继承人}、\Quote{同继承人}。因此，他说\Quote{恩宠极其丰盛}。因为所赐的恩惠超出了慈悲，它们不仅是慈悲的结果，更是亲情与超乎寻常的爱的结果。在如此阐明天主的爱之后——这爱不满足于向亵渎者和迫害者显示慈悲，反而将其他福泽丰富地赐予他——他加上\Quote{连同在基督耶稣内的信德与爱德}，以防止不信者剥夺自由意志的错误。他说，我们只贡献了这一点：我们相信祂能拯救我们。\par

\Emph{伦理教训}。那么，让我们藉着基督爱天主。\Quote{藉着基督}是什么意思？就是说，是祂，而不是法律，使我们能够这样做。注意我们欠基督什么，欠法律什么。他不仅说恩宠丰盛，而是说\Quote{极其丰盛}，因为它立刻将那些应受无限惩罚的人带到了义子的地位。\par

再请注意，\Quote{in}是用以表示\Quote{through}的。因为不仅信德是必需的，而且爱德也是。因为仍有许多人相信基督是天主，却并不爱祂，行为也不像爱祂的人。因为当他们将一切置于祂之前：金钱、出生、命运、占卜、预兆，那又是怎么回事呢？当我们生活悖逆祂时，请问，我们的爱德何在？谁若有一个温暖而深情的朋友？愿他也同样爱基督。那么，即使不能再多，也愿他爱那为我们这些祂的仇敌、我们自己毫无功德可言的、赐下祂的圣子的那一位。我说功德吗？我们犯了无数的罪，竟敢冒犯祂到极点，且毫无理由！然而，祂在无数次显示良善与眷顾之后，甚至那时也没有抛弃我们。就在我们极大地得罪祂的时刻，祂却为我们赐下祂的圣子。而我们，在接受了如此大的恩惠之后，被提升为祂的朋友，并因祂而被视为配得一切福分，却仍没有像爱朋友那样爱祂！那么，我们还能有什么希望呢？你或许因这话而战栗，但我倒愿你因事实而战栗！什么？你说，怎能显明我们甚至没有像爱朋友那样爱天主？我将努力向你们说明——但愿我的话是毫无根据、毫无意义的！但我恐怕事实正是如此。请想想：真正的朋友，常常为他们所爱的人承受亏损。但为了基督，没有人愿意承受亏损，甚至不以自己的现状为足。为了朋友，我们乐意忍受侮辱，挑起争斗；但为了基督，没有人能忍受敌意；俗话说：\Quote{宁可无缘无故被爱，也不要无缘无故被恨。}\par

我们当中没有一个人会不接济一位饥饿的朋友；然而基督每天来到我们这里，所求的并不是大事，仅仅只是饼，我们却连看也不看他一眼，尽管我们自己饱胀欲呕，因贪食而肚腹鼓胀：我们的气息泄漏出昨日的酒气，我们过着奢侈的生活，将财产浪费在娼妓、食客和谄媚者身上，甚至浪费在怪物、白痴和侏儒身上；因为人将这类天生的缺陷转化为取乐的材料。再者，对于真正的朋友，我们不嫉妒，也不因他们的成功而苦恼；然而对基督（的仆人）我们却怀有这种情绪，并且我们对人的友谊显得比敬畏天主更强大，因为嫉妒和不真诚的人显然敬畏人胜过敬畏天主。这是怎么回事呢？天主洞察人心，人却在天主面前毫不收敛他的欺诈；然而若这同一个人被人发现欺诈，他便认为自己完了，羞愧得脸红。何必提这个呢？如果一位朋友身处困境，我们去探望他，倘若我们稍有拖延，便怕受责备。但我们不去探望基督，尽管他在监狱里一而再、再而三地死去；不仅如此，如果我们在信友中有朋友，我们去探望他们，不是因为他们基督徒，而是因为他们是我们的朋友。这样，我们做事没有一件出于对天主的敬畏或爱，而是有些出于友谊，有些出于习惯。当我们看见一位朋友远行时，我们哭泣、忧心；如果我们看见他离世，我们哀悼他，尽管我们知道我们不会分离太久，他将在复活时归还给我们。但基督离开我们，或者不如说我们天天拒绝他，我们却不悲伤，也不觉得伤害他、冒犯他、激怒他、做他不喜悦的事是奇怪的；可怕的事不是我们不把他当朋友对待；因为我要证明我们甚至把他当仇敌对待。你问：怎么？因为\BibleQuote{随从肉性的思念，是与天主为仇}，正如\Person{保禄}所说，而我们将这思念常带在身边。当基督向我们走来，来到我们的门口时，我们迫害他。因为邪恶的行为实际上就是如此，我们每一天因我们的贪婪和强夺而使他受侮辱。难道有人因宣讲他的圣言、造福他的教会而获得好名声吗？那么他成了嫉妒的对象，因为他做天主的工作。我们以为我们嫉妒他，但我们的嫉妒传到基督身上。我们假装希望这益处不来自别人，而来自我们自己。但这不能是为了基督，而是为了我们自己；否则，好处由别人或我们自己来做都应是无所谓的。如果一个医生发现自己无法医治他那快要失明的儿子，他会拒绝另一位能够完成治疗的医生的帮助吗？远非如此！\Quote{让我的儿子康复}，他几乎会对他说，\Quote{无论是由你还是由我。}为什么呢？因为他不会考虑自己，而是考虑什么对儿子有益。同样，如果我们顾及\Quote{基督}，就会引导我们说：\Quote{让好事成就，无论是由我们自己还是由任何其他人。}正如\Person{保禄}所说：\BibleQuote{或是假意，或是诚心，无论如何，基督是被传扬了。}\BibleRef{斐 1:18}。同样，\Person{梅瑟}回答时，当有些人想激怒他对\Person{厄耳达得}和\Person{摩丹}发怒，因为他们说预言：\BibleQuote{你为我的缘故嫉妒人吗？但愿上主的人民都能成为先知！}\BibleRef{户 11:29}。这些嫉妒的情感出于虚荣；这不正是反对者和仇敌的情感吗？有人说你的坏话？爱他！你说不可能。不，只要你愿意，是可能的。因为你若只爱那说你好的，你有什么可夸的呢？那不是为了主的缘故，而是为了那人好话的缘故。有人伤害了你？向他行善！因为报答那对你有恩的人，功德极小。你受了极大的冤屈和损失？设法用相反的来补偿。是的，我恳求你。让我们这样尽我们的本分。让我们停止仇恨和伤害我们的仇敌。他命令我们\BibleQuote{爱你们的仇敌}\BibleRef{玛 5:44}：但我们却迫害那爱我们的他。愿天主禁止！我们都在口头上说，但行为上却不如此。我们的心智被罪恶如此蒙蔽，以致我们在行为上容忍那些口头上认为不可容忍的事。那么，让我们停止那些损害和毁灭我们救恩的事，好使我们获得那些作为他朋友所能获得的祝福。因为基督说：\BibleQuote{父啊！我所在的地方，愿我的门徒也同我在那里，好使他们看见我的荣耀}\BibleRef{若 17:24}，愿我们众人借着耶稣基督的恩宠与爱，都能达到这境地。\par

\SourceRecord{Translated by Philip Schaff.；Nicene and Post-Nicene Fathers, First Series；Vol. 13.；Edited by Philip Schaff.；Buffalo, NY: Christian Literature Publishing Co.,；1889.；<http://www.newadvent.org/fathers/230603.htm>.}

% structure: p0001 source=h1 target=section reason=page-title

\section{弟茂德前书讲道集第四篇}

\BibleRef{弟前 1:15-16}\par

\Emph{\BibleQuote{这话是确实的，值得完全接纳：耶稣基督来到世上，为拯救罪人；其中我是罪魁。然而我蒙了怜悯，是为要叫我这个罪魁在耶稣基督身上显出祂一切的忍耐，给后来信祂得永生的人作榜样。}}\par

天主的恩惠远超人的希望和期待，以致往往人不相信。因为天主赐予我们人从未期望、从未想到的事。为此，宗徒们费了许多口舌，务使我们相信天主赐予我们的恩赐。因为人领受了很大的好处时，会问自己是否在做梦，不敢相信；天主的恩赐也是如此。那么，什么事被认为难以置信呢？就是那些曾是仇敌、罪人，既未因法律成义，也未因行为成义的人，竟可以立即因着信德而获得最高的恩宠。为此，保禄在\WorkTitle{罗马书}中已详细论述，在此又详细论述。他说：\BibleQuote{这话是确实的，}他说，\BibleQuote{值得完全接纳：耶稣基督来到世上，为拯救罪人。}\par

由于犹太人特别受这一点吸引，他便说服他们不要听从法律，因为他们若没有信德就不能获得救恩。这正是他争辩的；因为在他们看来，一个从前一生在虚妄和邪恶的行为中度过的人，竟能因着他的信德得救，这是难以置信的。为此他说：\BibleQuote{这是值得相信的话。}但有些人不仅不信，甚至还反对，正如现在希腊人所做的：\Quote{那么，我们作恶以成善吧。}这是他们从他的话\BibleQuote{罪恶在哪里越多，恩宠在那里也格外丰富}中戏谑得出推论。\BibleRef{罗 3:8}所以当我们向他们谈论地狱时，他们说：这怎配得上天主呢？当人发现自己的仆人犯罪，他宽恕他，认为他配得宽恕，难道天主却永恒地惩罚吗？当我们谈及圣洗及藉之赦罪时，他们也说：这也不配天主，因为一个犯下无数罪行的人竟能得赦免。这是何等的心态乖戾，何等的好辩精神！确实，如果宽恕是恶，那么惩罚就是善；但如果惩罚是恶，那么宽恕它便是善。我是按着他们的观念说的，因为按着我们的观念，二者都是善。这我将在另一时间说明，因为目前时间不够处理如此深奥且需要详细辩论的主题。我必须在适当的时候向你们的爱德讲述。目前我们继续我们既定的主题。他说：\BibleQuote{这话是确实的。}但为何它值得相信呢？\par

这从前文和后文都可以看出。请看他是如何预备我们接受这个断言，然后又如何详细阐述。因为他先前已宣告，祂对我这个\BibleQuote{亵渎者和迫害者}动了怜悯；这便是一种预备。并且，祂不但动了怜悯，还\BibleQuote{算我是忠信的}。他的意思是，我们远不应怀疑祂动了怜悯。因为如果一个人看见一个囚犯被带进王宫，他不会怀疑那人得了怜悯。而保禄的处境正是如此，因为他以自己为榜样。他也不羞于称自己为罪人，反而以此为乐，因为这样最能显明天主对他的眷顾的神迹，以及天主认为他配得如此特别的恩惠。\par

然而，他在这里称自己为罪人，甚至罪魁，而在别处他却断言自己\BibleQuote{就法律所要求的义而言，是无瑕可指的}。\BibleRef{斐 3:6}这是因为，相对于天主所成就的义，即那真正追求的成义，就连在法律上为义的人也是罪人，\BibleQuote{因为所有的人都犯了罪，缺少了天主的光荣}。\BibleRef{罗 3:23}因此，他并未简单地说\Quote{义}，而是说\Quote{法律所要求的义}。好比一个人获得了财富，在自己看来是富有的，但与王室的宝藏相比，却非常贫穷，而且是穷人中的最贫穷者；此处的道理也是如此。与天使相比，即使义人也是罪人；如果保禄，他履行了法律所要求的义，尚且是罪魁，那么还有什么人可以称为义人呢？因为他这样说并非要谴责自己的人生不洁，切勿这样想；而是将自己法律上的义与天主的义相比较，显出其毫无价值，不仅如此，他还证明那些拥有法律之义的人也是罪人。\par

第16节。\BibleQuote{然而我蒙了怜悯，是为要叫我这个罪魁在耶稣基督身上显出祂一切的忍耐，给后来信祂得永生的人作榜样。}\BibleRef{弟前 1:16}\par

请注意他如何进一步谦卑自贬，又举出一个新的、不那么光彩的理由。因为他因无知而蒙受怜悯，这并不那么意味着那蒙怜悯者是个罪人，或深陷定罪；但要说他蒙受怜悯，是为了使今后没有一个罪人因绝望而不敢寻求怜悯，而是每人确信必得同样的恩惠，这实在是一种极度的谦卑，乃至他曾自称为罪人之首，\BibleQuote{一个亵渎者、迫害者，不配称为宗徒}，也未曾说过类似的话。这可以通过一个例子来说明。假设一座城，居民全是恶人，有的更恶，有的稍轻，但都应受惩罚；其中有一个人比其余所有人更该受罚，犯有各种恶行。若宣布国王愿意赦免所有人，就不如亲眼看见这个最恶的罪犯真的被赦免那样容易相信。那时就不再有任何怀疑了。这就是保禄所说的：天主愿意给人充分保证，祂赦免他们的一切过犯，就拣选了那比任何人更罪恶的人作为祂怜悯的对象；因为当我蒙怜悯时，他论证说，别人就没有怀疑的余地了：就如平常说话时我们会说，\Quote{如果天主赦免了这样的人，祂就永远不会惩罚任何人了}；如此他表明自己虽不配得赦免，但为了别人的得救，首先获得了那赦免。因此，他说，既然我得了救恩，就谁也不要怀疑救恩了。且看这位蒙福者的谦卑；他没有说，\BibleQuote{为在我身上彰显}祂的\BibleQuote{宽忍}，而是\BibleQuote{一切宽忍}；好像他曾说，在任何情形下，祂都不能显出比在我身上更大的宽忍了，也找不到一个罪人如此需要祂的一切赦免、一切宽忍；不是仅一部分，就像那些部分犯罪的人那样，而是\BibleQuote{一切}宽忍。\par

\BibleQuote{为给那些将来相信的人作榜样。}这话是为安慰、为鼓励。但因他曾高声赞扬圣子，赞扬祂所显示的伟大爱情，以免有人以为他把圣父排除在外，他也把这光荣归于圣父。\par

第17节。\BibleQuote{愿尊贵和荣耀归于永世的君王，那不朽坏、不可见、独一智慧的天主，直到永远。阿们。}（\BibleRef{弟前 1:17}）\par

因此，我们不仅光荣圣子，也光荣圣父。让我们在这里与异端辩论。论到圣父，他说：\BibleQuote{归于独一智慧的天主。}那么圣子不是天主吗？\BibleQuote{归于独一的不朽者。}那么圣子不是不朽的吗？或者祂自己并不拥有那以后将赐予我们的不朽？他们回答说，圣子是天主，是不朽的，但不与圣父相同。那么，是什么？祂的本质较低，因此不朽也较低吗？那么什么是较大或较小的不朽？不朽不是别的，就是不受朽坏。因为光荣有大有小，但不朽没有大小之分，正如健康没有大小之分。要么是可朽坏的，要么是完全不可朽坏的。那么我们也像祂一样不朽吗？千万不要！当然不是！为什么？因为祂按本性拥有，而我们则是外来的。那么你为什么区别？他们说，圣父不是由任何别的成为这样的，而圣子是由圣父而成这样的。这我们也承认，不否认圣子是从圣父不可朽坏地受生的。他（保禄）的意思是，我们光荣圣父，因为祂生了这样的圣子。这样，当圣子做了大事，圣父便最受光荣。因为圣子的光荣又归于祂。既然祂生了全能而自足的圣子，那么圣子自足、自立、无软弱，这不仅是圣子的光荣，更是圣父的光荣。论到圣子，经上说：\BibleQuote{藉着祂造了世界。}（\BibleRef{希 1:2}）在我们中间，创造与制作有区别。因为一个制作、劳苦、执行，另一个统治；为什么？因为执行的是较低的。但那里并非如此；主权不在一个，制作在另一个。因为我们听见\BibleQuote{藉着祂造了世界}时，并不把圣父排除在创造之外。当我们说\Quote{归于不朽的君王}时，也不否认圣子的统治。因为这些是两者共同的，每项都属于二者。圣父创造，即生了创造的圣子；圣子统治，作为万有的主。因为祂不是为雇佣而工作，也不是服从别人，像我们中的工人那样，而是出于自己的良善和对人类的爱。圣子曾被看见过吗？无人能断言。那么，\BibleQuote{归于不朽、不可见、独一智慧的天主}是什么意思？或者当说：\BibleQuote{除祂以外，没有别的名字使我们得救}（\BibleRef{宗 4:12}）：又说：\BibleQuote{在别的里面没有救恩}（\BibleRef{宗 4:12}）？\par

\BibleQuote{愿尊贵和荣耀归于祂，直到永远。阿们。}\par

如今，尊贵和荣耀并非空话；既然祂不仅用言语，也用祂为我们所做的尊荣了我们，我们也当用行为和事实尊荣祂。但这尊荣临到我们，却达不到祂，因为祂不需要来自我们的尊荣，我们却需要来自祂的尊荣。\par

因此，尊荣祂，就是尊荣自己。睁开眼凝视太阳之光的人，自己获得快乐，欣赏那星辰之美，却并未给那发光体带去恩惠，也没有增加它的光辉，因为它仍是原来的样子；对于天主更是如此。那赞美、尊荣天主的人，是为自己的救恩和最大的益处而行；何以如此？因为他追随德行，且被祂尊荣。因为祂说：\BibleQuote{尊重我的，我必尊重他。}（\BibleRef{撒上 4:30}）然而，如果祂从我们的尊荣中得不到益处，祂又如何被尊荣呢？就如祂被说成饥渴一样。祂承担我们的一切，为能以各种方式吸引我们归向祂。祂被说成接受尊荣，甚至侮辱，好使我们畏惧。尽管如此，我们仍不归向祂！\par

\Emph{伦理}. 那么，让我们\Quote{光荣天主}，并在\Quote{我们的身体和我们的灵魂}上承载天主。\BibleRef{格前 6:20} 如何以身体光荣祂？又如何以灵魂光荣祂？这里灵魂被称为神，以区别于身体。但我们如何能在身体和灵魂上光荣祂？那在身体上光荣祂的人，不犯奸淫或淫乱，避免贪食和醉酒，不追求外表的炫耀，只为自己作健康所需的准备；同样，那妇女不涂抹香料或脂粉，满足于天主所造的模样，不添加自己的巧技。你为什么在你所造的工作上添加你自己的装饰？天主的工艺还不够你满足吗？或者你试图为它增添优雅，仿佛你是更好的艺术家？你这样装饰自己的身体，并非为了自己，而是为了吸引大批的爱慕者，并且侮辱了你的造物主。不要说：\Quote{我能做什么？这并非我的意愿，但为了丈夫我必须这样做。除非我同意，否则我无法赢得他的爱。}天主使你美丽，是为了让人在美丽中钦佩祂，而不是遭受侮辱。因此，不要如此忘恩负义，而要以谦逊和贞洁来报答祂。天主使你美丽，是为了增加你谦逊的考验。因为一个可爱的人保持谦逊，比一个没有吸引力的人更难。圣经为什么告诉我们，\Quote{若瑟体态秀雅，容貌俊美}\BibleRef{创 39:6}？岂非为了让我们更钦佩他美貌中的谦逊？天主使你美丽，你为什么使自己不美丽？因为如同有人用泥巴涂抹金像，那些使用脂粉的妇女也是如此。你用红白泥土涂抹自己！但你说，相貌平平的人有理由这样做。为什么？为了隐藏丑陋？这是徒劳的尝试。因为自然的外貌何时能被矫揉造作所改善？你为何因不美而苦恼，既然这不是耻辱？因为听智者的话：\Quote{不要因人的美貌称赞他，也不要因人的外貌而厌恶他。}\BibleRef{德 11:2} 让天主，这位最卓越的工匠，更受钦佩，而非人——人没有任何功劳使他成为现在的样子。告诉我，美貌有什么好处？没有。它使拥有者遭受更大的考验、不幸、危险和猜疑。没有美貌的人免于猜疑；拥有美貌的人，除非她极其谨慎，否则会招致恶名，而比一切更糟的是，丈夫的猜疑——他看见她的美貌时所得的快乐，不如嫉妒所带来的痛苦。她的美貌因熟悉而在眼中褪色，而她的品格却因软弱、放纵和放荡的指控而受损，她的灵魂也变得卑下，充满傲慢。美貌面临这些邪恶。但没有这种吸引力的人则安然无恙。狗不会攻击她；她如同一只羔羊，安息在安全的牧场，没有狼闯入骚扰她，因为牧人在旁保护。\par

真正的优越，不在于一个人美丽而另一个人平庸，而在于一种优越：一个人即使不美丽，却不贞洁；另一个人却不邪恶。告诉我，眼睛的完美在于什么？是柔软、转动、圆润、黝黑，还是清晰和敏锐？灯的完美在于形状优雅、制作精良，还是放射光芒、照亮整个房屋？我们不能说不是这样，因为前者无关紧要，后者才是真正的目的。因此，我们常对负责灯的女仆说：\Quote{你把灯做坏了。}所以灯的作用完全是发光。同样，眼睛的外观无关紧要，只要它完全履行职责。我们把昏暗或混乱、睁开却看不见的眼睛称为坏眼睛。因为坏的正是不能履行其本职——这正是眼睛的缺陷。至于鼻子，告诉我，什么时候它是好的？是直挺、两侧光滑、比例匀称？还是能迅速接收气味并传达到大脑？任何人都能回答这个问题。\par

现在，让我们用一个例子来说明——比如夹子，我指的是那种称为夹子的工具；我们说那些能够夹起并握住东西的夹子是好的，而不仅仅是形状美观优雅的。同样，好的牙齿是那些适合分割和咀嚼食物的，而不是那些排列美丽的。将同样的道理应用到身体的其他部分，我们称那些健康、正确履行本职的肢体为美丽。因此，我们认为任何工具、植物或动物是好的，不是因其形状或颜色，而是因其达成目的。一个仆人被认为是好的，是因为他有用于我们的服务、随时待命，而不是因为相貌俊美却放荡。我相信你现在明白，你如何能够变得美丽。\par

既然最大、最重要的恩惠是所有人平等享有的，我们就没有任何不利。无论美丽与否，我们同样看见这宇宙、太阳、月亮和星辰；我们呼吸同样的空气，同样分享水和地上的果实。如果我们说些奇怪的话，平庸的人比美丽的人更健康。因为美丽的人为了保持美貌，不参与劳动，沉溺于怠惰和精致的生活，由此损害了身体的精力；而其他人没有这种顾虑，将所有注意力单纯而完全地投入积极的活动。\par

我们该当\BibleQuote{光荣天主，并在我们的身体上承担和背负祂}（\BibleRef{格前 6:20}）。我们不要追求美丽的容貌，那种关心是虚妄而无益的。我们不要教导我们的丈夫只欣赏纯粹外在的形态；因为如果这就是你的装扮，他习惯看你的脸庞，就容易让一个娼妓迷住。但如果你教导他喜爱善行和谦逊，他就不会轻易游移，因为他看到在娼妓身上没有这些品质，反而相反。也不可教导他被笑谈或放荡的衣饰所吸引，免得你为自己预备毒药。让他习惯于以谦逊为乐，如果你衣着端庄，就能做到这点。但如果你有招摇的神态、不稳重的举止，你怎能以严肃的语调对他说话呢？谁不会藐视和嘲笑你呢？\par

但是，我们如何能在我们的精神上光荣天主呢？通过实践德行，装饰灵魂。因为这种装饰并不禁止。因此，当我们各方面都好时，我们光荣天主，而在那个大日子，祂将更加光荣我们。因为\BibleQuote{我断定，现今的苦楚，与那将要显示于我们的光荣，是不能比较的}（\BibleRef{罗 8:18}）。愿我们大家都能分享这光荣，这是天主因我们的主耶稣基督的恩宠和慈爱所赐予的。\par

\SourceRecord{Translated by Philip Schaff.；Nicene and Post-Nicene Fathers, First Series；Vol. 13.；Edited by Philip Schaff.；Buffalo, NY: Christian Literature Publishing Co.,；1889.；<http://www.newadvent.org/fathers/230604.htm>.}

% structure: p0001 source=h1 target=section reason=page-title

\section{《论弟茂德前书》第五篇讲道词}

弟茂德前书 1:18, 19\par

\BibleQuote{\Emph{我儿}（\OriginalTerm{孩子}{\GreekText{τέκνον}}）\Emph{弟茂德，我依据以前指着你所说的预言，把这训令托付给你，好叫你凭那些预言，能打一场好仗；保持信德和良心纯洁；有些人抛弃了良心，就在信德上翻了船。}}\par

教师和司铎的职务具有极大的尊严，而推举一位称职者需要天主的拣选。古时如此，现今亦然，当我们不带人情、不因任何现世考虑、不为友情或敌意所动摇而作选择时，也是如此。因为我们虽不像古圣那样领受那么大的圣神恩赐，但一个良好的意向足以吸引天主的拣选临于我们。宗徒们在拣选玛弟亚时，尚未领受圣神，但他们将这事交托于祈祷，便拣选他加入宗徒之列。因为他们不看重人间的友情。如今我们也该如此。但我们已陷入极度的疏忽，甚至连明显的事都放过。若是我们忽略明显的事，天主又怎会将隐密的事启示给我们呢？正如经上所说：\BibleQuote{你们在小事上不忠信，谁还把大事和真事托付给你们呢？}\BibleRef{路 16:11} 但在那时，当没有任何人为因素时，司铎的任命也是藉着预言进行的。所谓\Quote{藉着预言}是什么意思？就是藉着圣神。预言不仅指预告未来，也包括启示现在。藉着预言，撒乌耳被发现\BibleQuote{藏在行李中}\BibleRef{撒上 10:22}。因为天主向义人启示。同样，也藉着预言说：\BibleQuote{你们给我选拔出巴尔纳伯和扫禄来}\BibleRef{宗 13:2}。弟茂德也是这样被选上的，保禄关于他用了复数的\Quote{预言}，也许是指保禄\Quote{带他行了割礼}的时候，以及在他按立弟茂德时，正如保禄在他给弟茂德的信中所说：\BibleQuote{不要忽略你身上的神恩}\BibleRef{弟前 4:14}。因此，为激励他，使他保持清醒警醒，保禄提醒他是谁拣选并按立了他，就好像说：\Quote{天主拣选了你。祂把使命交托给你，你不是由人的投票产生的。所以，不要滥用或羞辱天主的任命。} 当保禄说到\Quote{训令}时，暗示了某种负担，他加上：\Quote{我儿弟茂德，我把这训令托付给你。} 他像对待儿子，自己的儿子一样训令他，不是凭专横或独断的权柄，而是如同父亲，说：\Quote{我儿弟茂德。} 但\Quote{托付}一词表示要谨慎持守，且不是我们自己的东西。因为这不是我们为自己取得的，而是天主赐予我们的；不仅这职务，还有\Quote{信德和良心纯洁}。我们所以当持守祂所赐予的。因为若祂不来，信德便不可寻，那由教育而来的纯洁生活也不可得。如同保禄说：\Quote{不是我在命令你，而是那拣选你的那一位}，而这正是\BibleQuote{依据以前指着你所说的预言}的意思。听从它们，服从它们。\par

又说：你命令什么？\BibleQuote{好叫你凭那些预言，能打一场好仗。} 他们拣选了你，你就该为此而\Quote{好仗}。他称\Quote{好仗}，因为还有不好的仗，就如他所说：\BibleQuote{你们曾将你们的肢体献给不洁和不义，作奴隶}\BibleRef{罗 6:19}。那些人在暴君手下服役，而你却在君王手下服役。为什么他称它为\Quote{仗}？为表明所有人，尤其是教师，必须进行何等大的搏斗；我们需要强壮的臂膀、清醒、警醒和持续警惕；我们必须为流血和冲突作准备，要列阵而战，毫不松懈。他说：\Quote{你在其中打仗}。因为在军队中，并非所有人都担任相同的职务，而是各有岗位；同样在教会中，有的担任教师，有的担任门徒，有的担任平信徒。而你是教师。因为这一点还不够，他又说：\par

19节。\BibleQuote{保持信德和良心纯洁。}\BibleRef{弟前 1:19}\par

因为想要成为教师的人，必须先教导自己。正如从未作过好兵的人，绝不会成为将军，教师也是如此；所以他在别处说：\BibleQuote{免得我给别人传了福音，自己反而被淘汰}\BibleRef{格前 9:27}。他说：\BibleQuote{保持信德和良心纯洁}，这样你才能管理他人。我们听了这话，不要轻视上级的劝勉，即使我们是教师。因为若是弟茂德——我们所有人加在一起也不配与他相比——尚且领受命令并受指教，而且他本人处于教师的职位，那么我们更该如此。\BibleQuote{有些人抛弃了良心，就在信德上翻了船}。这是自然的。因为当生活败坏时，就会生出与之相称的道理，由此许多人陷入邪恶的深渊，转向外教。因为他们为了不受未来恐惧的折磨，试图说服自己的灵魂，说我们所宣讲的是假的。有些人因完全凭理性推究而偏离信德；因为推究招致翻船，而信德则如安全的船。\par

那些偏离信德的人必要遭遇翻船，保禄用一个例子表明这事。\par

20节。\BibleQuote{其中有依默纳约和亚历山大}\BibleRef{弟前 1:20}\par

他要藉此教导我们。你看，从那时起已有引诱人的教师，好奇的探究者，以及远离信德、凭自己理性探究的人。正如船难者赤身露体、一无所有，那背离信德的人也同样无助，不知该立足何处或安身何处，也无法从善生中获得任何益处。因为头若错乱，身体其余部分有何用？倘若没有善生的信德无济于事，那么相反的情况更是如此。如果天主为了我们而轻视他自己的，那么我们更该为了他而轻视自己的。因为事实如此：任何人若背离信德，便毫无稳定，他左右漂浮，直到最终沉入深渊。\par

\Quote{我已把他们交给撒殚，使他们受训不再亵渎！}可见，凭自己理性探究神圣之事便是亵渎。人的理性与神圣之事有何共通之处？然而撒殚如何训导他们不再亵渎？他尚未教导自己，仍是亵渎者，怎能教导别人？并非\Quote{他应教导}，而是\Emph{他们应受教导}。并非撒殚成就此事，尽管结果是如此。正如他在别处论及那行淫者时所说：\Quote{把这样的人交给撒殚，败坏他的肉体。}不是为拯救肉体，而是为\Quote{使灵魂得救。}\BibleRef{格前 5:5} 因此，这是无人称的说法。这如何实现呢？正如刽子手虽自身背负无数罪行，却成为他人的纠正者；恶神也是如此。但为何你不亲自惩罚他们，如你对待那个巴尔耶稣，或如伯多禄对待阿纳尼雅，反而把他们交给撒殚？不是为惩罚他们，而是为使他们受训。因为他有此权能，从其他经文可见：\Quote{你们愿意怎样？要我带着棍棒到你们那里去吗？}\BibleRef{格前 4:21} 又说：\Quote{免得我照主所赐给我的权能过于严厉，这权能是为建立，而不是为破坏。}\BibleRef{格后 13:10} 那么，为何他要呼求撒殚惩罚他们？为使羞辱更大，因为严厉和惩罚更为显著。或者，他们亲自惩罚那些尚未相信的人，但那些背离的人，他们则交给撒殚。为何伯多禄惩罚阿纳尼雅？因为当他在试探圣神时，仍是不信者。为使不信者得知他们无法逃脱，他们亲自施加惩罚；但那些已得知此事却后来背离的人，他们交给撒殚；表明他们并非靠自己的能力维持，而是靠对他们的关心；凡骄傲自大的人都被交给他。如同国王亲手杀死敌人，却将臣民交给刽子手惩罚，此处亦然。这些事为显示赐予宗徒们的权柄。这权柄不小：能够如此使魔鬼服从他们的命令。这表明魔鬼即便不情愿也事奉服从他们，这是恩宠权能的明证。请听如何交付：\Quote{当你们聚会时，我的灵魂，并借我们主耶稣基督的权能，把这样的人交给撒殚。}\BibleRef{格前 5:4} 他立即被逐出公共集会，与羊群分离，变得孤独无助；他被交给狼。因为如同云柱标示希伯来人的营帐，同样圣神区分了教会。因此，任何人若在外，便被吞灭；因宗徒的判决，他被逐出羊栈。主也这样把犹达斯交给撒殚。因为\Quote{吃过那饼后，撒殚就进入了他里面。}\BibleRef{若 13:27} 或者可以说：那些他们希望改正的人，他们并不亲自惩罚，而把惩罚留给怙恶不悛的人。或者，因交付他人而更令人畏惧。约伯也被交给撒殚，但不是因为他的罪，而是为更充分地证明他的价值。\par

这样的事例至今仍多有发生。因为司铎不能知道谁是罪人、谁是配不上地领受圣奥迹的人，天主便往往如此将他们交付给撒殚。疾病、攻击、忧苦、灾祸等之所以临到他们，正是为此。保禄指明了这点：\Quote{因此，在你们中间有许多软弱患病的人，死的也不少。}（\BibleRef{格前 11:30}）但有人说：我们不过一年领受一次！这正是祸患所在：你们衡量自己领受的相称与否，不是按心灵的纯洁，而是按时间的间隔。你们以为不常领圣体便是谨慎，却不知即便只一次，不相称地领受，便受了烙伤；而相称地领受，即使频繁，也是有益的。频繁领受并非冒昧，但整个人生中仅一次不相称地领受，便是冒昧。但我们愚昧可悲：一年中犯了无数过犯，却不急于从这些过犯中获得赦免，反而满足于自己不常鲁莽无耻地侮辱基督的身体，却忘了那些钉死基督的人，也只钉了他一次。那么，因为只犯了一次，过犯就减轻了吗？犹达斯只出卖了主一次。难道这使他免受惩罚吗？在这事上，为什么还要考虑时间呢？我们领受的时候，应当是良心纯洁的时候。复活节的奥迹并不比现在庆祝的奥迹更有功效。两者是同一的。圣神的恩宠相同；它常是逾越节。你们这些已领受入门圣事的人知道：在预备日、安息日、主日、殉道者纪念日，所举行的是同一祭献。保禄说：\Quote{你们每次吃这饼，喝这杯，就是宣告主的死亡。}（\BibleRef{格前 11:26}）这祭献的举行并没有限定时间；那么为何那时被称为逾越节？因为那时基督为我们受苦。所以，不要让时间影响你领受的态度。在任何时候，都有同样的能力、同样的尊威、同样的恩宠、同一个身体；没有一次庆祝比另一次更圣洁或更不圣洁。你们这些在这场合中只看到世间帷幔和更华丽的随从的人，也知道这一点。这些日子所唯一多的，是从它们开始，当基督被祭献时，我们的救恩之日就开始了。但就这些奥迹而言，那些日子并无更高之处。\par

当你走近领取肉体食物时，你会洗手和口；但当你走近这神粮时，你却没有洁净灵魂，反而满是不洁地走近。但你说：四十天的斋戒难道不足以洗清我们罪过的巨堆吗？告诉我，这有什么用？如果你想收藏某种珍贵的香膏，你会把容器弄干净来盛放，但放进去不久后又往里面倒垃圾，那香气岂不消失了吗？我们身上也发生这样的事：我们尽己所能使自己配得亲近，然后又再次玷污自己！那这又有什么好处呢？这话我们甚至是对那些能在四十天内洗净自己的人说的。\par

所以，我恳求你们，不要忽略我们的救恩，免得我们的劳作归于徒然。因为那离弃自己罪恶，却又去犯同样罪的人，就像\Quote{狗转过来吃自己所吐的。}（\BibleRef{箴 26:11}）但如果我们按所应当的行事，谨慎自己的道路，我们必将配得那些高贵的赏报。愿天主因我们的主耶稣基督的恩宠和慈爱，使我们众人获得这赏报。\par

\SourceRecord{Translated by Philip Schaff.；Nicene and Post-Nicene Fathers, First Series；Vol. 13.；Edited by Philip Schaff.；Buffalo, NY: Christian Literature Publishing Co.,；1889.；<http://www.newadvent.org/fathers/230605.htm>.}

% structure: p0001 source=h1 target=section reason=page-title

\section{第六篇讲道论\WorkTitle{弟茂德前书}}

\Emph{\BibleRef{弟前 2:1-4}}\par

\Emph{\BibleQuote{所以我劝勉，首先要为人人祈求、祈祷、转求和谢恩；为君王及一切有权位者，好使我们能敬虔端庄地度宁静平安的生活。这在我们的救主天主面前是美好可悦纳的，祂愿意人人得救，并认识真理。} 【R.V.: 祂愿意人人得救，等等。】}\par

司铎是全世界的共同父亲，正如他事奉的天主一样。因此他理应关心所有人。为此他说：\BibleQuote{所以我劝勉，首先要为人人祈求、祈祷、转求和谢恩。}由此产生两项益处。首先，消除对局外人的仇恨；因为没有人会恨他所祈祷的人；而他们也因着为他们所作的祈祷而变得更好，并失去对我们的凶暴态度。因为没有什么比爱与被爱更能吸引人服从教导。想想看，那些迫害、鞭打、放逐、屠杀基督徒的人，听到那些被他们如此野蛮对待的人为他们向天主献上热忱的祈祷，这会是怎样的情形。请注意他如何希望基督徒超越一切恶待。正如一个怀抱婴儿的父亲被小儿打在脸上，也不会减少对它的慈爱；同样，我们即使被外人打击，也不应减少对他们的善意。\Quote{首先}是什么意思？它指在日常礼仪中；而受过入门教导的人知道，这如何每天在早晚实行，我们如何为全世界、为君王和一切有权位者献上祈祷。但也许有人会说，他是指不是为所有人，而是为所有信徒。那么他为何提到君王？因为那时君王并非天主的敬拜者，有一长串不虔敬的统治者。而且为了不显得奉承他们，他先说\Quote{为人人}，然后\Quote{为君王}；因为他若只提君王，可能会引起怀疑。而且，既然某些基督徒的心灵听到这个可能会迟疑，并拒绝劝告，如果在举行神圣奥迹时必须为外教君王祈祷，他就向他们展示其益处，至少以这种方式使他们接受劝告，\BibleQuote{好使我们能度宁静平安的生活}；意思是说，他们的安全是我们的保障；正如他在致罗马人书中，劝他们服从统治者，\BibleQuote{不是因愤怒，而是因良心}。\BibleRef{罗 13:5} 因为天主为了公共福祉设立了政府。因此，当他们为此而作战，并守护我们的安全时，我们不为他们战争和危险中的安全而祈祷，难道不合理吗？这不是奉承，而是符合正义的规则。因为如果他们得不到保存，在战争中不顺利，我们的事务必定陷入混乱和麻烦；如果他们被消灭，我们要么自己服务，要么像逃亡者一样四处流散。因为他们是我们面前筑起的壁垒，被围在其中的人得享和平与安全。\par

他说：\BibleQuote{祈求、祈祷、转求和谢恩。} 因为我们必须为他人所得的恩惠感谢天主，就如祂使太阳既照恶人也照善人，降雨给义人也给不义之人。请注意他如何不仅藉祈祷，还藉谢恩将我们联合并束缚在一起。因为被迫为邻人的恩惠感谢天主的人，也必然爱他，并对他心存善意。如果我们必须为邻人的恩惠感恩，那么更要为自己所遇的事、为未知的事、甚至为违背我们意愿的事和看似痛苦的事感恩，因为天主安排万事皆为我们有益。\par

\Emph{道德教训}。让我们每一次祈祷都伴随谢恩。如果我们受命为近人祈祷，不仅为信徒，也为不信者，那么想想，祈祷反对你的弟兄是何等错误。什么？祂命令你为仇敌祈祷，而你却祈祷反对你的弟兄？但你的祈祷不是反对他，而是反对你自己。因为你用这些不敬的话激怒天主：\Quote{同样待他！} \Quote{如此待他！} \Quote{击打他！} \Quote{报复他！} 这样的话远非基督门徒所应为，他们应谦和温良。从那曾蒙恩领受如此神圣奥迹的口，不可说出任何苦毒之言。那曾接触主圣体的舌头，不可说出冒犯的话，务要纯洁，不可承载咒骂。因为如果\BibleQuote{辱骂者必不能继承天主的国}\BibleRef{格前 6:10}，那么咒骂者更将不能。因为咒骂者必是伤害人者；伤害与祈祷彼此对立，咒骂与祈祷相距甚远，控告与祈祷更是天壤之别。你用祈祷和息天主，却又发出诅咒？如果你不宽恕，你也不会被宽恕。\BibleRef{玛 6:15} 但你不宽恕，反倒恳求天主不宽恕，这是多么极端的邪恶！如果不宽恕者得不到宽恕，那么祈求主不宽恕的人，又如何得宽恕？伤害的是你自己，不是他。因为即使你的祈祷即将为你自己蒙允，在这种情况下，因从污秽的口中献上，也绝不会被接纳。因为那咒骂的口确实被所有冒犯和不洁所玷污。\par

当你本应为自己的罪战栗，恳切地寻求赦免时，你却来激动天主反对你的弟兄——难道你不惧怕，也不想想自己的事吗？你难道看不见自己在做什么吗？甚至效法学校里孩童的行为。如果他们看见同班的人在背诵功课，都因懒惰而挨打，一个接一个地被严厉盘问并受棍棒责打，他们就会吓得要死；即使他们的同伴打了他们，而且打得很重，他们也无暇生气，也不向老师告状；他们的灵魂充满了恐惧。他们只关注一件事：如何能平安进出而不受鞭打，他们的心思都在那个时刻。当他们出来时，无论是否挨打，从玩伴那里受到的打击从不因为他们喜悦而进入他们的脑海。而你，当你为自己的罪焦虑地站立时，你怎能不战栗地提及别人的过失？你怎能恳求天主的赦免？因为你以诅咒别人的方式使自己的处境更糟，并且你禁止他对你的过失宽容。他岂不会说：\Quote{如果你要我如此严厉地追讨别人对你的冒犯，你怎能期望我赦免你对我所犯的冒犯？} 让我们终于学会做基督徒吧！如果我们不知道如何祈祷——这是一件非常简单容易的事——我们还能知道什么呢？让我们学着像基督徒一样祈祷。那些是外邦人的祈祷，犹太人的恳求。基督徒的祈祷正相反，为了赦免和忘记别人对我们的冒犯。经上说：\Quote{被人辱骂，我们就祝福；被人迫害，我们就忍受；被人毁谤，我们就劝慰。} \BibleRef{格前 4:12} 听听斯德望说：\Quote{主，不要将这罪归于他们。} \BibleRef{宗 6:60} 他不是祈祷反对他们，而是为他们祈祷。你呢，不是为他们祈祷，而是对他们发出诅咒。那么，你在邪恶的程度上与他卓越的程度相当。请问我们钦佩谁？是那些被他祈祷的人，还是为他祈祷的人？当然是后者！如果我们这样，那么天主更是如此。你希望你的敌人被击倒吗？为他祈祷：然而不要带着这样的意图，不是要击倒他。那确实会是结果，但不要把它当作你的目标。那位蒙福的殉道者无辜地承受了一切，却仍为他们祈祷；我们从敌人那里遭受许多事，都是罪有应得。如果那位无辜受苦的人尚且不敢不为他的敌人祈祷，那么我们应该受什么惩罚呢？我们遭受的是应得的惩罚，却不为他们祈祷，反而祈祷反对他们？你确实以为你是在打击别人，但实际上你是把剑刺向自己。你以这种激起天主对别人愤怒的方式，不允许审判者对你自己的过错宽容。因为 \Quote{你们用什么尺度量，也要用什么尺度量给你们；你们用什么审判审判人，也要受什么审判。} \BibleRef{玛 7:2} 因此，让我们愿意宽恕，好使天主也愿意宽恕我们。\par

我愿你们不仅听闻这些事，更要遵行。因为此刻记忆只存留了言词，也许连这些也存留不住。我们彼此分离后，若有人未曾与会，问你们我们讲了什么，有人说不出来；有人只知道我们谈论的主题，回答说曾有一篇讲道，论及宽恕伤害和为仇人祈祷，却将所说的一切完全遗漏，因为记不住；另一些人记得一点，但终究有限。因此，如果你们听了毫无所得，我恳求你们甚至不要来听讲。因为有何益处呢？受的谴责更大，惩罚更重，如果经过如此多的劝诫，我们仍依然故我。为此，天主赐给我们一篇确定的祈祷文，好使我们不祈求任何属人、任何属世的事物。你们这些信友知道你们应当祈求什么，知道整篇祈祷文是共通的。但有人说：\Quote{那里并没有命令要为不信者祈祷。}倘若你们理解那篇祈祷文的效力、深邃和隐藏的宝藏，你们就不会这样说。只要展开它，你们就会发现这一点也包含其中。因为当人在祈祷中说\BibleQuote{愿祢的旨意奉行在人间，如同在天上}时，其中便隐含了这一点。如今，因为天上没有不信者，也没有冒犯者；因此，如果这祈祷文只为信友，那说法就毫无意义了。倘若信友遵行天主的旨意，而不信者不遵行，祂的旨意便没有奉行在人间，如同在天上。但它的意思是：天上没有恶人，同样愿地上也没有；但要吸引众人敬畏祢，使众人成为天使，连那些恨我们、与我们为敌的人也是如此。你们岂不见天主每日被信友和不信者在言语和行为上亵渎和讥嘲吗？那么怎样？祂为此曾熄灭太阳吗？或停止月亮运行吗？祂曾粉碎天穹、拔起大地吗？祂曾使海干涸吗？曾封闭水源吗？或搅乱空气吗？不，相反，祂使祂的太阳升起，降雨降下，按时赐下地上的果实，如此向亵渎者、麻木者、污秽者、迫害者供给年年的养料，不是一两天，而是他们的一生。那么你们要效法祂，尽人性能力所及去仿效祂。你们不能叫太阳升起吗？那就禁止恶言。你们不能降雨吗？那就止住咒骂。你们不能赐食物吗？那就戒除傲慢。你们这样的奉献已足够。天主对祂仇人的良善由祂的作为显明。你至少用言语如此做：为你的仇人祈祷，这样你便像你在天之父了。我们讲过多少次这个主题了！我们也不会停止讲论；只愿有些效果。并非我们困倦或厌倦讲话；只求你们听的人不要厌烦。如今，当人不愿遵行所说的话时，似乎就厌烦。因为实行的人，常常爱听同样的事，并不因此厌烦；因为这正是对他的称赞。但厌烦纯粹来自不遵行所吩咐的。因此说话者变得令人讨厌。如果一个人实践施舍，听见别人谈论施舍，他并不厌倦，反而喜悦，因为他听见自己的善行被推荐和宣扬。因此，当我们听有关宽恕伤害的讲道而不悦时，那是因为我们对容忍不感兴趣，我们没有实践它；因为如果我们拥有这实际，就不因它被提起而痛苦。因此，如果你们不愿我们令人疲惫或厌烦，就按照我们所宣讲的去实行，在你们的行动中展示我们讲论的主题。因为我们永不会停止讲论这些事，直到你们的行为与它们相符。我们这样做，更是出于对你们的关心和爱。因为号手必须吹号，即使无人出战；他必须尽自己的本分。我们这样做，并非希望给你们带来更重的谴责，而是为从我们自己身上避免它。此外，对你们的爱催迫我们，因为如果那事（愿天主禁止！）临到你们，它将撕裂并折磨我们的心。我们所推荐的并非什么昂贵的过程：它不需要糟蹋财物，也不需要漫长艰辛的旅程。只需意愿。它是一个字，一个心中的决意。让我们只在舌上设守卫，在嘴唇安门和闩，好使我们不说出任何冒犯天主的话。这样做是为我们自己的益处，并非为那些我们为他们祈祷的人。因为让我们永远思考：那祝福自己仇人的，是祝福自己；那咒骂自己仇人的，是咒骂自己；那为自己仇人祈祷的，不是为他祈祷，而是为自己祈祷。如果我们如此行，我们将能够实行这卓越的德行，并如此获得所应许的福分，赖我们的主耶稣基督的恩宠和慈爱。\par

\SourceRecord{Translated by Philip Schaff.；Nicene and Post-Nicene Fathers, First Series；Vol. 13.；Edited by Philip Schaff.；Buffalo, NY: Christian Literature Publishing Co.,；1889.；<http://www.newadvent.org/fathers/230606.htm>.}

% structure: p0001 source=h1 target=section reason=page-title

\section{弟茂德前书讲道集 第七篇}

\BibleRef{弟前 2:2-4}\par

\BibleQuote{愿我们能平安宁静地度日，在一切虔敬和端庄中生活。因为这在我们的救主天主面前是美好且可悦纳的；祂愿意所有人都得救，并认识真理。}\par

如果为了结束公开的战争、骚乱和战斗，司铎被劝诫为君王和官长献上祈祷，那么个人更应当这样做。因为有三类极其严重的战争：第一类是公开的，当我们的士兵遭受外军攻击时；第二类是即使在和平时期，我们彼此相争；第三类是个人与自己相争，这是最糟糕的。因为对外战争不能对我们造成很大伤害。请问，即使它屠杀我们、将我们剪除，又能怎样？它不伤害灵魂。第二类也不能违背我们意愿地伤害我们，因为即便别人与我们相争，我们自己仍可保持和平。先知如此说：\BibleQuote{他们因我的爱而与我为敌，但我却专务祈祷} \BibleRef{咏 108:4}；又说：\BibleQuote{我与那些恨恶和平的人和平相处} \BibleRef{咏 119:6}；以及：\BibleQuote{我愿和平，但我一发言，他们便备战。}但第三类，我们无法无危险地逃脱。因为当身体与灵魂相争，激起邪恶的欲望，用感官享乐或愤怒、嫉妒的不良情欲来武装自己反对灵魂时，我们便无法获得所应许的福分，直到这场战争结束；凡不平息这骚动的人，必被击倒，所受的创伤将带来地狱中的死亡。因此，我们每天都需要谨慎和极大的焦虑，使这战争不在我们内挑起，或者，如果挑起了，也不要持续，而是被平息并进入沉睡。因为如果整个世界享有深远的和平，而你与自己相争，那有什么益处呢？因此，这乃是我们应当保持的和平。如果我们拥有它，任何外在的事物都不能伤害我们。为此，公共和平贡献不小：因此经上说：\BibleQuote{愿我们能平安宁静地度日。}但如果有人在安静时心烦意乱，他是一个可怜的人。你看见祂所说的和平就是我称为第三类的吗？因此当祂说\BibleQuote{愿我们能平安宁静地度日}时，祂并不止于此，而是加上\BibleQuote{在一切虔敬和端庄中。}但我们若没有那和平，就不能生活在虔敬和端庄中。因为当好奇的推理扰乱我们的信德时，还有什么和平？或者当不洁的灵时，还有什么和平？\par

为了使我们不以为祂说的是所有人都过的那种生活，当祂说\BibleQuote{愿我们能平安宁静地度日}时，祂又加上\BibleQuote{在一切虔敬和端庄中}，因为外邦人、放荡者以及骄奢淫逸之人都可以过着平安宁静的生活。显然，祂的意思并非如此，从祂所加上的\BibleQuote{在一切虔敬和端庄中}就可以明白。这样的生活充满了陷阱和冲突，灵魂每日被自己思想的骚动所伤。但祂真正所指的生活从下文显而易见，也由于祂并非简单地说虔敬，而是加上\BibleQuote{一切虔敬}。因为说这话时，祂似乎坚持的虔敬不仅在于教义，也在于生活所支撑的虔敬，因为两方面都必须要求虔敬。因为教义虔敬而生活不虔敬有什么益处呢？生活不虔敬是极有可能的，请听那位有福的宗徒在别处所说：\BibleQuote{他们声称认识天主，但行为上却否认祂。} \BibleRef{铎 1:16} 又说：\BibleQuote{他已背弃信德，比不信者更坏。} \BibleRef{弟前 5:8} 以及：\BibleQuote{若有称为弟兄的人，是行淫的、或贪婪的、或拜偶像的} \BibleRef{格前 5:11} ，这样的人不尊重天主。还有：\BibleQuote{凡恨自己弟兄的，就不认识天主。} \BibleRef{若一 2:9} 不虔敬的方式各式各样。因此祂说：\BibleQuote{一切虔敬和规范。}因为不仅行淫者不端庄，贪婪的人也可称为无秩序和不节制。因为贪财是一种欲望，不亚于身体的食欲，凡不克制它的人就被称为放荡。因为人被称为放荡，是由于不约束自己的欲望，所以易怒的、嫉妒的、贪婪的、诡诈的，以及一切生活在罪中的人，都可称为放荡、无秩序、放纵。\par

第三节。\BibleQuote{因为这在我们的救主天主面前是美好且可悦纳的。} \BibleRef{弟前 2:3}\par

所谓\Quote{可悦纳的}是指什么？是为所有人祈祷。这是天主所接纳的，这是祂所愿意的。\par

第四节。\BibleQuote{祂愿意所有人都得救，并认识真理。} \BibleRef{弟前 2:4}\par

效法天主！如果祂愿意所有人都得救，那么就有理由为所有人祈祷；如果祂愿意所有人都得救，你也要愿意；如果你愿意，就为此祈祷，因为愿望引向祈祷。注意祂从各方面催促灵魂为外邦人祈祷，表明从中生出多么大的益处；\BibleQuote{使我们能平安宁静地度日}；而比这更重要的，是它中悦天主，如此人便相似祂，因为他们愿意祂所愿意的。这足以使野兽羞愧。因此，不要害怕为外邦人祈祷，因为天主自己愿意如此；但只应害怕为任何人作反对他们的祈祷，因为祂不愿意。如果你为外邦人祈祷，当然也应该为异端者祈祷，因为我们当为所有人祈祷，而不是迫害。这还有另一个理由，因为我们共有同一本性，天主命令并悦纳彼此之间的慈爱和亲情。\par

但你说：若主自己愿意赐予，何需我的祈祷？这祈祷对他人与自己都大有裨益。它吸引人走向爱德，又倾向你于仁爱。它有能力吸引他人信奉；因为许多人因彼此争竞而远离天主。这也就是他如今所称的天主的救恩：\BibleQuote{祂愿意所有的人得救}；没有这个，其他一切都无甚重要，只是名义上的救恩，空谈而已。\BibleQuote{并得以认识真理。}真理是什么真理？信德在于祂。其实他先前已说过：\BibleQuote{嘱咐某些人不要传别的道理。}但为使无人视他们为仇敌，因而挑起纷争反对他们，他说：\BibleQuote{祂愿意所有的人得救，并得以认识真理}；说了这话后，他又加上了：\par

第五节。\BibleQuote{因为只有一位天主，也只有一个中保，在天主与人之间。}\BibleRef{弟前 2:5}\par

他先前说：\BibleQuote{得以认识真理}，暗示世界不在真理之中。如今他说：\BibleQuote{只有一位天主}，意即并非像有些人所说的有许多神，并说祂派遣了祂的圣子为中保，以此证明祂愿意所有的人得救。但圣子岂非天主？的确祂是；那么为何他说\BibleQuote{一位天主}？这是相对于偶像而言，并非相对于圣子。因为他谈论的是真理与谬误。中保应当与两方都有共融，即他所中保的双方。因为中保的特性就是与他所中保的每一方紧密共融。若他只与一方相连而与另一方分离，他就不再是中保了。因此，若祂不分享圣父的本性，祂就不是中保，而是分离的。正如祂因来到人间而分享人的本性，同样祂因出自天主而分享天主的本性。因为祂要中保于两种本性之间，祂必须接近这两种本性；正如处于两地之间的地方与每地相连，同样处于本性之间的也必须与每种本性相连。因此，祂成为人，也同时是天主。人不能成为中保，因为他还需向天主恳求。天主也不能成为中保，因为那些祂应从中保的人不能接纳祂。正如他在别处说：\BibleQuote{只有一个天主圣父……也只有一个主耶稣基督}\BibleRef{格前 8:6}；同样在这里也说\BibleQuote{一位}天主和\BibleQuote{一位}中保；他没有说两位；因为他不想那个数目被歪曲为他所谈论的多神论。所以他写下了\BibleQuote{一位}和\BibleQuote{一位}。你看圣经的措辞多么准确！因为虽然一加一是二，但我们不该这么说，尽管理性提示如此。而在这里，你虽说并非一加一是二，但你说的却是理性所不提示的。\Quote{如果祂生了，祂也受了苦。}\BibleQuote{因为只有一位天主，}他说，\BibleQuote{也只有一个中保，在天主与人之间，就是那人基督耶稣。}\par

第六节。\BibleQuote{祂舍了自己作万民的赎价，这证实了在适当的时期。}\BibleRef{弟前 2:6}\par

那么基督也是外邦人的赎价吗？毫无疑问，基督甚至为外邦人死了；而你却不愿为他们祈祷。那么你问，为什么他们不信？因为他们不愿意；但祂的那部分已经完成。祂的受苦是一个\BibleQuote{证实}，他说；因为祂来，本为\BibleQuote{为真理作证}，即天父的真理，并被杀害。如此，不仅天父为祂作证，祂也为天父作证。\BibleQuote{因为我奉我父的名而来，}祂说\BibleRef{若 5:43}。又说：\BibleQuote{从来没有人见过天主。}\BibleRef{若 1:18}又说：\BibleQuote{使他们认识你，唯一的真天主。}\BibleRef{若 17:3}还有：\BibleQuote{天主是神，}\BibleRef{若 4:24}祂甚至作证至死。但这\BibleQuote{在适当的时期}意即在最适宜的时期。\par

第七节。\BibleQuote{我为此被立为宣讲者及宗徒——我在基督内说实话，并不说谎——外邦人的教师，在信德和真理上。}\BibleRef{弟前 2:7}\par

既然基督为外邦人受苦，且我被分别为\BibleQuote{外邦人的教师}，你为何拒绝为他们祈祷？他充分表明自己的可信性，说他是为此\BibleQuote{被立}\BibleRef{宗 13:2}，即被分别出来，当时其他宗徒在教导外邦人一事上退后；他加上\BibleQuote{在信德和真理上}，以表明那信德中毫无诡诈。可见恩宠的扩展。因为犹太人从不为外邦人祈祷；但如今恩宠延及他们；当他说他被分别为外邦人的教师时，他暗示恩宠现已遍及世界每一角落。\par

\Quote{祂把自己当作赎价，}祂说，那么祂如何被圣父交出呢？因为是出于祂的良善。\Quote{赎价}是什么意思？天主本要惩罚他们，但祂没有这样做。他们将要灭亡，但祂代替他们赐下了自己的圣子，并派遣我们作使者去宣讲十字架。这些事足以吸引众人，并显示基督的爱。\Emph{道德训诲}。确实，天主赐予我们的恩惠是如此真实、如此不可言喻的伟大。祂为仇视、憎恨、拒绝祂的敌人牺牲了自己。没有人会为朋友、为兄弟、为儿女做的事，主却为祂的仆人做了；这位主自己并非像祂的仆人那样，而是天主为人，为不配的人。因为如果他们是配得的，如果他们遵行了祂的旨意，那就没有那么奇妙了；但祂为如此忘恩负义、如此顽梗的受造物而死，这使每个人的心灵都感到惊异。因为人不会为自己的同类做的事，天主却为我们做了！然而，在祂向我们显示这样的爱之后，我们却退缩，没有认真地去爱基督。祂为我们牺牲了自己；我们却没有为祂做出任何牺牲。当祂需要必需的饮食时，我们忽略祂；当祂患病赤身露体时，我们不探望祂。我们不该受什么？什么愤怒、什么惩罚、什么地狱？即使没有其他诱因，这一点也足以说服每一个人：祂屈尊把人类的苦难当作自己的，说：\Quote{我饥饿，我口渴。}\par

财富的暴政啊！或者更确切地说，是那些甘愿做它奴隶的人的邪恶！因为财富本身没有大能，而是因着我们的软弱和奴性：是我们卑鄙、卑躬屈膝、属肉体和没有悟性。因为金钱有什么能力呢？它是哑的、无感觉的。如果魔鬼，那个邪恶的灵，那个狡猾的万物搅乱者，没有能力，那么金钱有什么能力呢？当你看到银子时，想象它是锡！你不能吗？那就把它当作它本来的样子；因为它本是泥土。但如果你不能这样推理，就想想我们也要死去，许多拥有它的人几乎没有从中得到任何益处，成千上万以它为荣的人现在已是尘土和灰烬。他们正在遭受极重的惩罚，远比如今用玻璃和陶器进食的人更贫穷；那些曾经躺在象牙床榻上的人，现在比躺在粪堆上的人更穷。但它悦目吗！有多少其他东西更悦目！花朵、纯净的天空、穹苍、明亮的太阳，远更令眼睛愉悦。因为它有许多锈迹，因此有人认为它是黑色的，这从变黑的雕像可以看出。但太阳、天穹、星辰中没有黑暗。这些灿烂之物远比它的颜色令人愉悦。因此，使它讨人喜欢的不是它的光彩，而是贪婪和邪恶；这些，而不是金钱，给人快乐。把这些从你的灵魂中驱逐出去，那么看起来如此珍贵的东西对你而言将比泥土更无价值。那些发烧的人渴望见到泥土，好像它是泉水；但健康的人很少甚至不想喝水。摆脱这种病态的渴望，你就会看到事物的本来面目。为了证明我没有说谎，要知道，我可以指出许多这样做过的人。扑灭这火焰，你就会看到这些东西比花朵更不值得。\par

黄金是好的吗？是的，它好用于施舍、救济穷人；它好，但不是为了无益的使用，不是为囤积或埋在地下，不是为戴在手上、脚上或头上。它被发现的目的是为了用它解救俘虏，而不是把它做成链子戴在天主的肖像上。为这个目的使用你的金子：解救被捆绑的人，而不是捆绑那自由的人。告诉我，你为什么把毫无价值的东西看得比一切更重？难道因为它是金的，它就不是链子了吗？材料有什么不同？无论是金是铁，它仍然是链子；不，金的更重。那么是什么使它轻省呢？无非是虚荣，以及被看见戴着链子的乐趣，你反该为此羞愧。为了显明这一点，把它系上，把佩戴者放在旷野或无人看见的地方，链子立刻就会觉得沉重，并被看作是负担。\par

亲爱的，让我们害怕，免得我们被判听到那些可怕的话：\BibleQuote{捆起他的手脚。} \BibleRef{玛 22:13} 女人啊，为什么你现在对自己这样做？没有囚犯手脚都被捆绑。你为什么还要捆绑你的头？因为你不满足于手脚，还用许多链子捆绑你的头和颈。我略过这些东西带来的挂虑、恐惧、惊慌、以及因它们与丈夫发生的争吵，如果丈夫需要它们，失去它们中的任何一件对人是何等的死亡。你能称之为快乐吗？为了悦他人的眼目，你使自己受捆绑、挂虑、危险、不安和每日的争吵？这该受一切责备和谴责。不，我恳求你，我们不要这样做，让我们挣脱每一种\BibleQuote{不义的束缚} \BibleRef{宗 8:23}；让我们把饼分给饥饿的人，并做所有其他能确保我们在天主面前有自信的事，使我们能获得借着我们的主耶稣基督所应许的福分。\par

\SourceRecord{Translated by Philip Schaff.；Nicene and Post-Nicene Fathers, First Series；Vol. 13.；Edited by Philip Schaff.；Buffalo, NY: Christian Literature Publishing Co.,；1889.；<http://www.newadvent.org/fathers/230607.htm>.}

% structure: p0001 source=h1 target=section reason=page-title

\section{第八篇讲道论弟茂德前书}

\BibleRef{弟前 2:8-10}\par

\BibleQuote{我愿男人在各地举起圣洁的手祈祷，不发怒，不怀疑。同样，也愿女人以端庄的服饰，以羞赧和自律装饰自己；不以编发、黄金、珍珠或昂贵的衣服为装饰；但要以善行——这原合乎那自承虔敬的女人。}\par

基督说：\BibleQuote{你们祈祷时，不要像假善人一样；因为他们爱在会堂里和十字路口站着祈祷，为显给人看。我实在告诉你们，他们已获得了他们的赏报。至于你，当你祈祷时，要进入你的内室，关上门，向你在暗中的父祈祷；你的父在暗中看见，必要报答你。} \BibleRef{玛 6:5} 那么保禄说什么？\BibleQuote{我愿男人在各地举起圣洁的手祈祷，不发怒，不怀疑。} 这并不与那相矛盾，断乎不是，而是完全协调。但是如何，以何种方式？我们必须先考虑\BibleQuote{进入你的内室}是什么意思，以及基督为何命令这事，如果我们可以在任何地方祈祷？或是我们不可在教堂祈祷，也不可在房屋的任何其他部分，而只可在内室？那么那句话是什么意思？基督是在推荐我们避免炫耀，当他命令我们不仅私下祈祷，而且秘密地祈祷。因为当他说\BibleQuote{不要叫左手知道右手所做的} \BibleRef{玛 6:3}，他考虑的不是手，而是命令他们极力防备炫耀；他在这里也是如此；他没有把祈祷限制在一个地方，而只要求一件事，即没有虚荣。保禄的目的在于区分基督徒与犹太人的祈祷，因此请注意他所说的：\BibleQuote{在各地举起圣洁的手}，这是犹太人所不允许的，因为他们不得在其他地方接近天主、献祭和举行他们的礼仪，而是从世界各地聚集在一个地方，他们必须在圣殿内举行全部敬拜。与此相反，他提出了他的训诫，并使他们脱离这种必要，他实际上说，我们的方式不像犹太人的；因为正如基督命令我们为众人祈祷，因为他为众人死了，我也为众人宣讲这些事，所以\BibleQuote{在各地祈祷}是好的。从此，所考虑的不是地方，而是祈祷的方式；\BibleQuote{在各地祈祷}，但\BibleQuote{在各地举起圣洁的手}。那才是所需要的。什么是\BibleQuote{圣洁的}？纯洁的。什么是纯洁的？不是用水洗的，而是脱离贪婪、凶杀、抢夺、暴力，\BibleQuote{不发怒，不怀疑}。这是什么意思？谁祈祷时会发怒？意思是不怀恶意。让祈祷者的心思纯洁，脱离一切激情。不要让任何人怀着敌意、或不可爱的性情、或\BibleQuote{怀疑}来接近天主。什么是\BibleQuote{不怀疑}？我们来听。它暗示我们不应有任何怀疑，以为我们不会被垂听。因为经上说：\BibleQuote{你们求，只要相信，就必得到。} \BibleRef{玛 21:22} 又说：\BibleQuote{当你们站着祈祷时，若你们对任何人有什么怨愤，要宽恕。} \BibleRef{谷 11:25} 这就是不发怒、不怀疑地祈祷。但我怎能相信我会得到所求呢？通过不求任何与他愿意赐予相悖的东西，不求任何与伟大君王不相称的东西，不求任何世俗的东西，而只求一切属灵的祝福；如果你\BibleQuote{不发怒}地接近他，有纯洁的手，\BibleQuote{圣洁的手}：用于施舍的手是圣洁的。这样接近他，你必定得到所求。因为\BibleQuote{你们虽然不好，尚且知道把好东西给你们的儿女，何况你们在天之父，岂不更把好东西给求他的人吗？} \BibleRef{玛 7:11} 通过怀疑，他指的是疑惑。同样，他说：我愿意女人不发怒、不怀疑地接近天主，举起圣洁的手：她们不应随从自己的欲望，也不应贪婪或抢夺。因为一个女人若不亲自抢劫或偷窃，而通过她的丈夫来做，那又如何？然而保禄对女人有更多要求，她们要装饰自己\BibleQuote{以端庄的服饰，以羞赧和自律；不以编发、黄金、珍珠或昂贵的衣服为装饰；但要以善行——这原合乎那自承虔敬的女人。} 但是什么是\BibleQuote{端庄的服饰}？那样的服装完全遮盖她们，端庄得体，没有多余的装饰，因为前者是合宜的，后者则不然。\par

\Emph{道德}。什么？你接近天主祈祷，却带着编发和金饰？你是来跳舞吗？来结婚吗？来参加喜庆游行吗？在那里，这样的编发、这样的贵重衣服是合时的，这里却一样也不需要。你来祈祷，来求赦罪，来为过犯恳求，祈求主，希望使他对你仁慈。为什么你要打扮自己？这不是恳求者的服装。你怎能呻吟？你怎能哭泣？怎能热切祈祷，当你如此装扮？你若哭泣，你的眼泪将成为旁观者的笑柄。哭泣的妇女不应佩戴黄金。那只是做戏和虚伪。因为不是做戏吗？从这样一个被奢侈和野心吞没的灵魂中流出眼泪？除去这样的虚伪！天主是轻慢不得的！这是演员和舞女的服装，她们靠舞台生活。这类东西一点也不适合端庄的女子，她们应以\BibleQuote{羞赧和自律}来装饰自己。\par

不要效法娼妓。因为她们用这样的服饰引诱许多情人；许多人因此蒙受可耻的嫌疑，非但没有从装饰中得到任何益处，反而因这种形象伤害了许多人。因为淫妇即使有端庄的名声，也无法从中获益，到了那日子，那审判人隐秘的主必使一切显明；同样，端庄的妇女若用这种服饰装作淫妇，就会失去贞洁的益处。因为许多人因这种看法而受害。\Quote{我能做什么，} 你说，\Quote{若别人怀疑我？} 但你却因你的服饰、你的眼神、你的举止而给了机会。因此\Person{保禄}论及服饰和端庄很多。如果他甚至要除去那些只是财富象征的东西，如金子、珍珠和华服；那么更要除去那些表示刻意装饰的东西，如涂脂、画眼、扭捏的步态、造作的声音、慵懒而淫荡的眼神；对披风、紧身衣、精美的腰带、紧贴的鞋子精心打扮？因为他说到\Quote{端庄的服饰}和\Quote{羞怯}时就暗示了这一切。因为这类事是可耻和不端。\par

请容忍我，我恳求你们，因为我的目的不是用这坦率的责备来伤害或痛苦你们，而是要移除我羊群中一切不相称的事。但如果这些禁令是针对那些有丈夫、富有、生活奢侈的妇女，那么更应针对那些宣发童贞的人。但你说，哪个贞女戴金子或编发？然而在简单的服装中也可能有刻意的精致，远超这些。你可以在普通服装上比穿金戴银的人更注重外表。因为当一件深色长袍用腰带紧束在胸前（如同舞台上的舞者所穿），精致到既不宽也不窄，恰到好处；胸部以许多褶皱衬托，这岂不是比任何丝袍更诱惑吗？当鞋子因黑色而闪亮，收成尖头，模仿绘画的雅致，几乎看不见鞋底的宽度——或者，虽然你确实不画脸，却花许多时间和精力洗脸，并在额前披上比脸还白的面纱——上面再戴一顶兜帽，其黑色衬托出白色的对比——这一切岂不是服饰的虚荣？对于眼珠不断转动，又能说什么？对于穿上胸衣，巧妙地时而隐藏、时而显露系扣？因为她们有时也暴露这些，以显示腰带的精美，将兜帽完全绕在头上。然后像演员一样，戴上紧贴的手套，仿佛长在手上；我们还可以谈论她们的步态和其他比金饰更迷人的技巧。亲爱的，让我们害怕，免得我们也听到先知对热衷外在装饰的希伯来妇女所说的话：\BibleQuote{以绳索代替腰带，以麻绳代替发辫，以秃头代替卷发。}\BibleRef{依 3:24}这些和许多其他只为被看见和吸引观看者而发明的东西，比金饰更迷人。这些不是小过，而是令天主不悦，足以破坏童贞的一切克己。\par

贞女，你以基督为你的新郎，为何要吸引世间的爱人？他要审判你如同淫妇。你为何不佩戴那使他喜悦的装饰，就是端庄、贞洁、有序和朴素的服饰？这是俗艳可耻的。我们再也无法分辨娼妓与贞女，她们已经到了如此无耻的地步。贞女的服饰不应刻意，而应朴素不费心；但如今她们有许多技巧使服饰显眼。妇人，停止这愚昧吧。将这份关心转移到你的灵魂上，转移到内在的装饰上。因为包裹你的外在装饰不容内在变得美丽。关心外表的人轻视内在，正如不关心外表的人专注内在。不要说：\Quote{唉！我穿着破旧的衣服、简陋的鞋、无价值的面纱；这些有什么装饰呢？}不要欺骗自己。正如我说过，用这些比用昂贵的衣服更能讲究外表，尤其是当它们紧贴身体、做成不端庄的样子、闪闪发亮时。你向我辩解，但你对天主能说什么，他知道你做事的心思和意念？\Quote{这不是为了邪淫！}也许不是，而是为了羡慕；你难道不因被羡慕这些事而羞耻吗？但你说：\Quote{我这样穿只是偶然，并无此类动机。}天主知道你对我说什么：难道你要向我交代吗？不，是要向那鉴察你行为、将来必审问你的那一位，万有在他面前都是赤裸敞开的。正因如此，我们现在劝诫这些事，免得你们受那些严厉的审判。因此让我们害怕，免得他借先知对犹太妇女所说的话责备你：\BibleQuote{她们到我这里来，行走时卖弄风情、扭捏作态，脚步叮当。}\BibleRef{依 3:16}\par

你承担了一场伟大的竞赛，其中需要的是摔跤，而非装饰；等待你的是战斗，而非怠惰与安逸。请看竞技场上的斗士与摔跤手。他们关心自己的步伐或衣着吗？不，他们藐视这一切，披上一件浸透油脂的外衣，只关注一件事：击伤对方，而不被击中。魔鬼咬牙切齿，伺机从各方面毁灭你，而你却毫不在意，或只在意这魔鬼的装饰。关于声音我暂且不提，虽然在这方面也多有矫饰；也不提香水及其他类似的奢侈。世俗妇女正是因这些事讥笑我们。贞洁的尊荣已丧失。没有人按照应有的方式尊敬贞女。她们给自己的不名誉提供了缘由。难道她们在天主的教会里不应被视若自天而降的女性吗？但现在她们被轻视，且是罪有应得，不过那些明智的人不在此列。但是，当一个有丈夫、子女且主持家务的女子，看见你——本该被钉死于世俗的人——比她自己更迷恋世俗时，她岂不要讥笑并蔑视你吗？看那操心！看那劳苦！你穿着朴素，却比那穿戴最昂贵饰品的人更甚；你比那浑身金饰的女子更刻意于外表。你不追求适合你的事，却追求那不适合你的，而你的本分原是致力于善工。因此，贞女不如世俗妇女受尊敬。因为她们没有行与贞洁誓愿相称的工作。这话不是对所有人说的；或者更确切地说，是对所有人说的：对有过失的人，好叫他们学习谦逊；对无过的人，好叫他们以谦逊教导他人。但要当心，这责备不要在实际行为中应验。因为我说这些，不是为叫你们忧愁，而是为纠正你们，好使我们为你们夸耀。但愿我们都能做天主所喜悦的事，并为光荣祂而生活，好获得我们的主耶稣基督的恩宠与慈爱所应许的福分，偕同祂，等等。\par

\SourceRecord{Translated by Philip Schaff.；Nicene and Post-Nicene Fathers, First Series；Vol. 13.；Edited by Philip Schaff.；Buffalo, NY: Christian Literature Publishing Co.,；1889.；<http://www.newadvent.org/fathers/230608.htm>.}

% structure: p0001 source=h1 target=section reason=page-title

\section{弟茂德前书第九篇讲道}

\BibleRef{弟前 2:11-15}\par

\BibleQuote{让她们在完全顺服中安静地学习。但我不允许女人教导，也不允许她管辖男人，只要沉静。因为先造的是亚当，后造的是夏娃；且不是亚当受了欺骗，而是女人受了欺骗，陷于过犯。然而她必在生产[借着生产]上得救，若她们常存信德、爱德、圣德，并持守节制。}\par

蒙福的保禄要求妇女极大的谦逊和得体的举止，不仅限于衣着和外貌，他甚至要规范她们的言语。他说：\BibleQuote{女人要在安静中学习}；即：她们在教堂里完全不要说话。他在致格林多人的书信中也给出了同样的规则，他说：\BibleQuote{妇女在教堂里说话是可耻的}\BibleRef{格前 14:35}，理由是法律使她们服从男人。又在别处说：\BibleQuote{若她们愿意学习什么，让她们在家问自己的丈夫}\BibleRef{格前 14:35}。那时，妇女因这样的教导而保持安静；但现在她们中间常有大声喧哗，许多吵闹和谈话，而且没有比这个地方更甚的。在这里，人们看到她们全都在说话，比在市场或浴池更甚。因为她们好像来这里消遣，全都忙于谈论无益的话题。这样一切混乱，她们似乎不明白，除非安静，她们无法学到任何有益的东西。因为当我们的讲道与谈话相抵触，没有人注意所说的话，对她们有什么好处呢？妇女应当安静到这种程度：在教堂里，不仅不允许谈论世俗之事，甚至不允许谈论属灵之事。这是秩序，这是谦逊，这比任何衣服更能装饰她。这样穿戴，她才能以最相称的方式献上祈祷。\par

\BibleQuote{但我不允许女人教导。}他说：\BibleQuote{我不容许}。这个命令在这里有什么地位？最相宜的。他正在谈论安静、得体、谦逊，所以说了他愿意她们在教堂里不说话之后，为断绝一切谈话的机会，他说：让她们不教导，而是占据学习者的位置。因为这样她们就会以沉默表示服从。因为女性生来有些爱说话；因此他各方面约束她们。他说：\BibleQuote{因为先造的是亚当，后造的是夏娃；而且不是亚当受了欺骗，而是女人受了欺骗，陷于过犯。}\par

如果问：这与现今的妇女有什么关系？这表明男性享受了更高的尊荣。男人是先造的；他在别处也表明他们的优越性。\BibleQuote{而且男人不是为女人创造的，女人是为男人创造的。}\BibleRef{格前 11:9}他为什么说这个？他希望男人在各方面居首位；他既因上述理由，也因之后发生的事，意指让男人居首位。因为女人曾一度教导了男人，使他成为不顺服的共犯，导致了我们的毁灭。因此，因为她滥用了她对男人、或不如说是与男人平等的权力，天主使她服从于她的丈夫。\BibleQuote{你要依附你的丈夫？}\BibleRef{创 3:16}这话以前并没有对她说过。\par

但亚当怎么没有被欺骗？如果他没有被欺骗，那么他也没有违犯？请仔细注意。女人说：\BibleQuote{那蛇欺骗了我}。但男人并没有说：\BibleQuote{女人欺骗了我}，而是说：\BibleQuote{她把树上的果子给我，我就吃了。}被一个同类、同一种类的受造物欺骗，与被一个低等和从属的动物欺骗，是不同的。这才是真正的被欺骗。因此，与女人相比，他被说成是\Quote{没有被欺骗}。因为她被一个低等和从属者迷惑，而他被一个平等者迷惑。再者，关于男人，并没有说他\BibleQuote{见那树好作食物}，而是关于女人，说她\BibleQuote{吃了，并给她丈夫}；所以他违犯不是被食欲所俘虏，而仅仅是出于妻子的劝说。女人教导了一次，就毁了一切。因此他说，让她不要教导。但这件事与其他女人何干？当然有关，因为女性是软弱和易变的，他是在谈论整个性别。因为他说\Quote{女人}而不是\Quote{厄娃}，这是整个性别的通称，不是她的专名。那么，因为她的过错，整个性别都包括在过犯中吗？正如他论及亚当时所说的：\BibleQuote{效法亚当违犯的样式，他是将来那位的预像}\BibleRef{罗 5:14}，同样，女性性别违犯了，男性却没有。那么，妇女不能得救吗？能，借着生产。因为他说\BibleQuote{若她们常存信德、爱德、圣德，并持守节制}，并不是指厄娃。什么信德？什么爱德？什么带着节制的圣德？意思是：\Quote{女人们，不要灰心，因为你们的性别招致了责备。天主给了你们另一个得救的机会，即通过教养子女，所以你们不仅靠自己，也靠他人得救。}看他提及了多少问题。他说：\BibleQuote{女人受了欺骗，陷于过犯}。哪个女人？厄娃。那么她必因生产得救？他不是这个意思，而是整个女性种族将得救。那么它没有陷于过犯吗？有，但厄娃违犯了，然而整个性别仍将\BibleQuote{因生产}而得救。为什么不是因她们自己的德行？难道他把其他人排除在这救恩之外吗？那么童贞女、不育的、在生育前丧夫的寡妇怎么办？她们会灭亡吗？她们没有希望吗？然而童贞女是最受尊敬的。他到底要说什么？\par

有些人这样解释他的意思：正如发生在第一个女人身上的事，导致了整个女性性别的从属地位（因为他说，既然厄娃是第二个受造的，并且被置于从属地位，就让其余的女性也处于从属地位），同样，因为她犯了罪，其余的女性也处于罪中。但这不是合理的推理；因为在创造时，一切都是天主的恩赐，但在这里，情况是女人犯罪的结果。然而，这就是他所说的话的含义。正如众人因一人而死，因为那人犯了罪，同样，整个女性种族都犯了罪，因为女人在罪中。但她不必悲伤。天主给了她不小的安慰，即生育。如果说这是出于本性，那么那件事也是出于本性；因为不仅那出于本性的事被赐予了，而且养育子女的事也被赐予了。\Quote{如果他们持守信德、爱德和清醒审慎的圣德}；也就是说，如果在生育之后，她们使子女持守爱德和纯洁。藉此，她们将因子女而获得不小的赏报，因为她们为基督的服役培养了战士。所谓圣德，他是指善生、端庄和清醒审慎。\par

第三章第一节。\BibleQuote{这话是确实的。}\BibleRef{弟前 3:1}\par

这话与当前的主题有关，而不是与下文有关，即关于主教的职务。因为既然有人怀疑，他便确认这话是真实的，即父亲们可以从子女的德行中获益，母亲们亦然，当她们好好抚养子女的时候。但如果她本人沉溺于邪恶和恶习呢？那时她还会因养育子女而获益吗？难道她不会很可能把子女教养得和她一样吗？因此，这话不是针对任何女人，而是针对有德行的妇人，说她也因此获得巨大的赏报。\par

\Emph{道德教训}。你们作父母的，请听这话，你们的养育子女不会失去赏报。他在下文也这样说：\BibleQuote{有好名声，因善行而受称赞；如果她养育了子女}\BibleRef{弟前 5:10}。在其他赞美中，他举出这一条，因为将天主所赐的子女奉献给天主，绝非轻省的赞美。因为如果他们所奠定的根基是好的，他们的赏报就大；如果他们疏忽了，他们的惩罚也同样大。厄里就是因他的子女而丧亡的。因为他本该告诫他们，事实上他也告诫了他们，却没有按应当的方式去做；但是因他不愿让他们痛苦，便毁灭了自己和他们。请听这话，你们作父亲的，要悉心养育子女，\BibleQuote{在主的教育和规劝之中}\BibleRef{弗 6:4}。青年是野性的，需要许多训导者、教师、指导者、陪伴者和监护人；即使有了所有这些，若能约束住也是幸运的。因为青年就像一匹未驯服的马，或一头未驯服的野兽。但是如果我们从一开始，从最早的年龄，就用良好的规条来固定它，日后就不再需要太多辛劳；因为养成的良好习惯对他们来说就像法律。我们不要让他们做任何虽有快感却有害的事；也不要纵容他们，以为他们不过是孩子。特别要在贞洁上训练他们，因为这是青年最大的祸害。为此，需要许多奋斗和极大的关注。让我们及早为他们娶妻，好让他们的新娘能接受纯洁无玷的身体，这样他们的爱情会更炽热。婚前贞洁的，婚后会更加贞洁；婚前行淫的，婚后也会行淫。\BibleQuote{一切饼，}经上记着，\BibleQuote{对犯奸淫者都是甘甜的。}\BibleRef{德 23:17}。新郎的头上通常戴花环，作为胜利的象征，表示他们未受快乐征服地进入婚姻之床。但如果他被快乐所俘，把自己交给了娼妓，那么他为何还戴花环呢？因为他已经受了征服。\par

让我们以这些事劝诫他们。让我们时而采用劝告，时而警告，时而威胁。子女是我们受托的重大责任。让我们对他们倾注极大的关心，并尽一切努力，使恶者不能从我们手中夺走他们。但现在我们的做法恰恰相反。我们确实尽一切努力使我们的农场秩序井然，并把它交给忠实的管家，我们为它寻找驴夫、骡夫、监工和精明的会计。但我们不去寻找更加重要得多的事，即一个我们可以把我们的儿子托付给他作为其道德监护的人，虽然这比所有其他财产都更有价值。我们确实是为儿子才如此关注我们的产业。我们为子女照管财产，却丝毫不关心子女本身。这是何等荒谬！正确塑造你儿子的灵魂，其余的一切都会加给你们。如果灵魂不好，财富对他毫无益处；如果灵魂被塑造得良善，贫穷也不会伤害他。你愿意他富有吗？教他良善：因为这样他就能致富；即使不能，他也不会比那些拥有财富的人更糟。但如果他是邪恶的，即使你留给他无限的财富，你也没有留下人来照管它，并且你使他比那些陷入极度贫困的人更糟。因为对于那些禀性不善的子女来说，贫穷比财富更好。因为贫穷能在一定程度上将他们约束在美德中，即使违反他们的意愿。而金钱却不允许那些愿意清醒的人继续保持清醒，它引导他们离开，毁灭他们，使他们陷入无尽的危险。\par

母亲们，要特别谨慎地好好管教你们的女儿；因为管理她们是容易的。要警醒照顾她们，使她们能做居家的人。最重要的是，教导她们虔诚、端庄、轻视财富、不在乎装饰。这样来安排她们的婚姻。因为如果你们这样塑造她们，你们不仅拯救她们，也拯救将要与她们结婚的丈夫，不仅丈夫，还有孩子，不仅孩子，还有孙子孙女。因为根若好，好枝条就会长出来，并且变得更好，而所有这一切你们都将得到赏报。因此，让我们做一切事，如同不仅有益于一个灵魂，而是通过那一个利益许多灵魂。因为她们本当从父家走向婚姻，如同从训练场出来的战士，具备所有必需的知识，并成为能将整团面转化为自身美德的酵母。要让你们的儿子如此端庄，以他们的稳重和节制著称，以便他们从天主和人那里都得到极大的称赞。让他们学习控制自己的欲望，避免奢侈浪费，成为好的管家，有爱心并服从规则。因为这样他们就能为父母确保好的赏报，如此一切事都将为光荣天主而做，并为了我们的救恩，借着我们的主耶稣基督，等等。\par

\SourceRecord{Translated by Philip Schaff.；Nicene and Post-Nicene Fathers, First Series；Vol. 13.；Edited by Philip Schaff.；Buffalo, NY: Christian Literature Publishing Co.,；1889.；<http://www.newadvent.org/fathers/230609.htm>.}

% structure: p0001 source=h1 target=section reason=page-title

\section{\BibleBook{弟茂德}{前书} 第十篇讲道}

\BibleRef{弟前 3:1-4}\par

\BibleQuote{如果有人渴望主教的职分，他渴望的是美好的工作。所以主教必须无可指摘，只作一个妻子的丈夫，有节制，清醒，端庄，好客，善于教导；不酗酒，不打架，不贪不义之财；但要温和，不争论，不贪婪；善于管理自己的家庭，使儿女凡事端庄顺服。}\par

既然现在要讨论主教的职务，他就首先指明主教应当是怎样的人。这里他并非在劝勉弟茂德的过程中谈论此事，而是向众人说话，并透过弟茂德教导其他人。他说了什么？\Quote{如果有人渴望主教的职分}，我并不责备他，因为这是保护的工作。若有人有此渴望，并非贪图管辖和权柄，而是愿意保护教会，我就不责备他。\Quote{因为他渴望的是美好的工作。} 甚至梅瑟也渴望这个职分，尽管不是权柄，而他的渴望使他遭受了那嘲讽：\BibleQuote{谁立了你作我们的首领和审判者？}\BibleRef{宗 7:27} 所以，若有人如此渴望，就让他渴望吧。因为 \OriginalTerm{主教职}{Episcopate} 的得名，就是源于监管一切。\par

\Quote{那么，主教}，他说，\Quote{必须无可指摘，只作一个妻子的丈夫。} 他立下这规矩，并非意味着主教不应没有妻子，而是禁止他拥有超过一个妻子。因为连犹太人也允许再婚，甚至同时有两个妻子。因为\BibleQuote{婚姻应受尊重} \BibleRef{希 13:4}。但有些人说，这话的意思是主教应当只作一个妻子的丈夫。\Quote{无可指摘。} 这个词内含一切德行；所以若有人自知有罪，就不应渴望被自己的行为所玷污的职务。因为这样的人应当被管理，而不是管理别人。因为作首领的人应当比任何明灯更明亮；他的生活应当毫无瑕疵，使众人都仰望他，并以他的生活为榜样。但他在使用这一劝言时，并非没有共同的目的。因为他自己也要任命主教（他也劝勉弟铎在写给弟铎的书信中这样做），而且很可能有许多人渴望那职务，因此他敦促这些劝诫。\Quote{有节制}，他说，就是谨慎，有千百只眼睛在身边，目光敏锐，不让心智的眼睛模糊。因为许多事情发生，使人不能清楚看见，看见事物的真相。因为忧虑、烦恼和来自四面八方的重担都压在他身上。所以他必须警醒，不仅为自己的事，也为别人的事。他必须保持清醒，精神热切，如同燃烧的火；他必须日夜劳苦、尽职，甚至超过将军对军队的关心；他必须关心和挂虑所有人。\Quote{端庄，好客。} 因为属下的大多数人都具备这些品质（在这些方面他们应当与统治他们的人相等），他为了指明主教的特质，就加上：\Quote{善于教导。} 因为这对于被管理的人并非要求，但对于受托管理的人却是最必要的。\par

\Quote{不酗酒}：这里他并非主要指醉酒不节，而是指傲慢和无耻。\Quote{不好斗}：这也不是指动手打人。那么 \Quote{不好斗} 是什么意思呢？因为有些人时不宜地打击弟兄的良心，似乎这话是针对他们说的。\Quote{不贪不义之财，但要温和，不好争论，不贪婪；善于管理自己的家庭，使儿女凡事端庄顺服。} 如果这样，那么 \BibleQuote{结了婚的人挂虑世俗的事}\BibleRef{格前 7:33}，而主教不应挂虑世俗的事，为什么他说只作一个妻子的丈夫呢？有些人认为他是针对不结婚的人说的。但若非如此，有妻子的人可以好像没有妻子。\BibleRef{格前 7:29} 因为那时这种自由是恰当允许的，适合当时环境的本性。而且，人若愿意，是可能这样规范行为的。因为正如财富使人难进天国，然而富人常常得以进入，婚姻也是如此。但为什么当他说主教应当 \Quote{不酗酒、好客} 时，他本应提到更大的事？为什么他没有说主教应当是天使般、不受人类情绪影响的人？那些基督所说的伟大品质在哪里，即使是那些在他治理下的人也应当具有？就是死于世界，时刻准备舍命，正如基督所说：\BibleQuote{善牧为羊舍命}\BibleRef{若 10:11}；又说：\BibleQuote{不背起自己的十字架跟随我的，不配作我的门徒}\BibleRef{玛 10:38}。但他说 \Quote{不酗酒}；真是好前景，如果这就是要劝诫主教的事情！为什么他没有说主教应当已经超脱世界？但你要求主教比世上的人还低吗？因为对这些世上的信徒，他说：\BibleQuote{治死你们在地上的肢体}\BibleRef{哥 3:5}，以及 \BibleQuote{死了的人已经脱离了罪}\BibleRef{罗 6:7}。\BibleQuote{凡属于基督的人已经把肉体钉在十字架上}；而基督又说：\BibleQuote{凡不撇下一切所有的，不配作我的门徒}\BibleRef{路 15:33}。\par

但因为教会将受到攻击，他所要求的不是那种卓越高超的德行，而是中等程度的德行；因为清醒、品行端正、节制，是许多人都具备的品质。\BibleQuote{让自己的儿女凡事端庄顺服。}这是必要的，好使在他的家中能树立榜样。因为谁肯相信一个不能使自己的儿子顺服的人，能够管理一个陌生人呢？\BibleQuote{善于管理自己的家。}连外人也说：善于管理家庭的人，也将是好的政治家。因为教会好比一个小家庭，就像在家里有儿女、妻子和仆役，男人管理他们所有的人；同样，在教会里有妇女、儿童和仆人。如果管理教会的人有同僚分享他的权力，那么男人也有一个同僚，就是他的妻子。教会难道不该供养她的寡妇和童贞女吗？同样，在家庭中也有仆人和女儿需要供养。事实上，管理家庭更容易；因此他问：\BibleQuote{人若不知道管理自己的家，怎能照管天主的教会呢？}\par

第六节。\BibleQuote{不可为初奉教者。}\BibleRef{弟前3:6}他不是说，不可为年轻人，而是不可为新奉教者。因为他曾说：\BibleQuote{我栽种了，阿颇罗浇灌了，但使之生长的是天主。}\BibleRef{格前3:6}他希望他们指出这样的人，就用了这个词。否则，有什么阻止他说\Quote{不可为年轻人}呢？因为如果只有年轻是个障碍，为什么他自己任命了年轻的弟茂德呢？（他藉着对他说\BibleQuote{不可让人小看你年轻}来证明这一点。）\BibleRef{弟前4:12}因为他知道弟茂德伟大的德行和严格的生活。知道这点的他写道：\BibleQuote{你自幼便学习了圣经。}\BibleRef{弟后3:15}而他实行严格禁食的证据是这句话：\BibleQuote{为了你屡次的病症，稍微用点酒。}他写给他这些以及其他事情；如果他不知道弟茂德这些善工，就不会写信给他，也不会给他的门徒这样的吩咐。但是当时有许多人从外教人皈依并受洗，他说：\Quote{不要立刻就把一个初奉教者，即这些新归化者之一，提拔到尊严的职位上。}因为，如果他在还没有好好做门徒之前，就立刻被立为教师，他就会变得傲慢。如果他在学会服从之前就被任命为统治者之一，他就会自高自大：因此他加上：\BibleQuote{免得他自高自大，就落在魔鬼所受的审判里，}即落入那因骄傲而临到撒殚的同一种审判中。\par

第七节。\BibleQuote{并且他必须在那些外人中也有好名声，免得他落在毁谤与魔鬼的罗网里。}\BibleRef{弟前3:7}\par

这话说得对，因为他一定会被他们指责；也许出于同样的原因，他说了\BibleQuote{一个妻子的丈夫}，尽管在别处他说：\BibleQuote{我愿众人像我一样！}\BibleRef{格前7:7}，即保持独身。因此，为了不把他们限制得太窄，不要求过份严厉的生活，他满足于规定中等的德行。因为需要在每个城市任命一位管理者，正如他写给弟铎的：\BibleQuote{你又应在各城设立长老，像我吩咐你的那样。}\BibleRef{铎1:5}但是，如果他有好名声和好的声誉，却不配得，那该怎么办呢？首先，这不容易发生。好人在敌人中获得好名声是很难的。但事实上，他并没有让这一点单独成立；他说：\Quote{也}要有好名声，即除此之外的其他品质。那么，如果他们出于嫉妒无理地毁谤他，特别是他们是外教人，该怎么办呢？这是不可预期的。因为连他们也尊重一个无可指摘的人。那么，为什么论到自己，他说：\BibleQuote{经历毁谤与美名}？\BibleRef{格后6:6}因为受到攻击的不是他的生活，而是他的宣讲。因此他说\BibleQuote{经历毁谤}。他们因宣讲被诽谤为欺骗者和骗子，这是因为他们无法攻击他们的道德品格和生活。为什么没有人说宗徒们是淫乱者、不洁者或贪婪者，而只说他们是欺骗者，这仅限于他们的宣讲？难道不是因为他们生活无可指摘吗？这是明显的。\par

因此，我们也当如此生活，这样就没有敌人，没有不信者能说我们的坏话。因为生活有德的人，连他们也尊敬。因为真理能堵住甚至敌人的口。\par

但他怎么\BibleQuote{落在罗网里}呢？是因屡次陷于同样的罪，如同那些外人一样。因为如果他是这样的人，那恶者很快就为他设下另一个网罗，并迅速使他毁灭。但如果他从敌人那里得到好名声，那他更会从朋友那里得到。因为一个生活无可指摘的人不太可能被毁谤，我们可以从基督的话推论：\BibleQuote{你们的光也当这样照在人前，叫他们看见你们的善行，光荣你们在天之父。}\EditorialNote{440 v. 16}但如果一个人被诬告，因特殊情况被诽谤呢？嗯，这是可能的情况；但即使这样的人也不应被提拔。因为结果非常可怕。因此经上说，他应当有\BibleQuote{好名声}，因为你们的善行是要发光的。所以，就像没有人会说太阳是黑暗的，甚至瞎子也不会（因为他会羞于反对所有人的意见），同样，没有谁会责备一个卓越良善的人。虽然外教人常因他的教义而诽谤他，但他们不会攻击他有德行的生活，反而会同其他人一起钦佩和尊敬它。\par

\Emph{道德}。那么，让我们这样生活，使天主的名不致受亵渎。一方面，我们不要着眼于人的声誉；另一方面，也不要使自己遭受恶名，而要在两方面都持守中庸，就如他所说：\Quote{你们在他们中，好像世上的光体。}\BibleRef{斐 2:15} 因为祂留我们在此世，正是为使我们成为光体，使我们被委任为他人之师，成为酵母，使我们如天使般在人间交往，如成人对待孩童，如属神者对待属血气者，使他们因我们得益，使我们成为种子，结出许多果实。如果我们能在生活中如此发光，就不需要言语了；如果我们能彰显善行，就不需要教师了。如果我们真是所应当的基督徒，就不会有外教人了。如果我们遵守基督的诫命，忍受冤屈，容许别人占我们的便宜，被辱骂却祝福，被虐待却善待他人 \BibleRef{格前 4:12}；如果这在我们的生活中成为常态，就没有人如此残忍而不皈依虔诚。为证明这一点：\Person{保禄}不过一人，却带领了多少人归向祂？如果我们都像他那样，我们能吸引多少个世界？看哪，基督徒的人数已超过外教人。在其他技艺中，一人可以同时教导一百个孩童；但在这里，教师比学习者多得多，却无人被领过来。因为受教者看重教师的德行；当他们看见我们表现同样的欲望，追求同样的目标——权力与荣誉——他们怎能钦佩基督教？他们看见我们的生活可指责，我们的灵魂属世。我们与同样人一样崇拜财富，甚至更甚。我们同样害怕死亡，同样畏惧贫穷，同样不耐疾病，同样贪求荣耀与统治。我们因贪财而折磨自己直至死亡，并服侍时代。那么，他们怎能相信？靠奇迹吗？但奇迹已不再发生。靠我们的言谈吗？它已腐化。靠爱德吗？已无处可寻。因此，我们不但要为自己的罪负责，也要为因这些罪对他人造成的伤害负责。\par

让我们于是恢复清醒，警醒度日，在地上彰显属天的生活。让我们说：\Quote{我们的家乡在天上}\BibleRef{斐 3:20}，并在地上坚持这场争斗。也许有人说，我们中间曾有伟大的人物，但希腊人会说：\Quote{怎样呢？}他说，\Quote{我怎能相信？因为我在你们的品行中看不到任何类似之事。如果你们这样自夸，我们也有自己的哲学家，他们的人生令人钦佩。}\Quote{但请展示给我另一位\Person{保禄}，或一位\Person{若望}：你们做不到。}他岂不因此嘲笑我们，因为我们的智慧只在言语上，不在行为上？如今，为了一枚小钱，你们就准备杀人或被杀！为了一小块土地，你们就接连起诉！为了一个孩子的死亡，你们就颠倒一切：我且不提那些令人哭泣的事——你们的占卜、预兆、迷信、占星术、符咒、占卜、咒语、巫术。这些都是呼喊的罪，足以激起天主的愤怒；祂派遣了自己的圣子之后，你们还敢做这样的事。\par

那么，除了哭泣，我们还能做什么？因为世上只有一小部分人在走向救恩的路，而那些丧亡的人听见了，就高兴自己并非独受惩罚，而是有许多同伴。但这有什么值得高兴的呢？这高兴本身将使他们受罚。不要以为在那里如同在这里，有同伴受苦能带来安慰。何以见得？让我说清楚。假设一人被命令烧死，他看见自己的儿子与他一同被烧，烧焦的肉味扑鼻而来；这岂不本身就是死吗？无疑。我还要告诉你们如何：即使那些没有受苦的人，看见这些事也会麻木、惊恐，何况那些自身受苦的人呢？不要对此感到惊奇。听一位智者的话：\Quote{你也变得软弱像我们吗？你也变得像我们一样吗？}\BibleRef{依 14:10} 因为人性倾向于同情，别人的苦难会触动我们的怜恤。那么，一个父亲看见儿子同受惩罚，或丈夫看见妻子，或人看见同胞，岂会得到安慰，而不是加重痛苦？我们岂不更加崩溃？但你说，在那里的情感不同。我知道不同；但那里有更悲惨的情感。因为那里有不能安慰的哀哭，所有人都目睹彼此的折磨。那些在痛苦中的人，因他人同受苦而得安慰吗？看见儿子、父亲、妻子或孙辈同受惩罚，断不是安慰！如果看见朋友如此，是安慰吗？不！不！这反而加重了我们自己的痛苦！此外，有些灾难因其严重性，不可能因共有而减轻。如果两人同时被扔进火里，他们会彼此安慰吗？告诉我：如果我们曾患剧烈的热病，岂不发现一切安慰都无济于事？因为有些灾难如此压倒一切，以致灵魂中没有安慰的余地。一个妻子失去丈夫，数算有多少人遭受同样损失，能减轻她的悲伤吗？因此，我们不要依靠任何这样的希望，反而要在悔改自己的罪中找到唯一的安慰，走上通往天堂的善路，好能获得天国，依靠我们的主耶稣基督的恩宠和慈爱。愿光荣归于祂，等等。\par

\SourceRecord{Translated by Philip Schaff.；Nicene and Post-Nicene Fathers, First Series；Vol. 13.；Edited by Philip Schaff.；Buffalo, NY: Christian Literature Publishing Co.,；1889.；<http://www.newadvent.org/fathers/230610.htm>.}

% structure: p0001 source=h1 target=section reason=page-title

\section{弟茂德前书第十一篇讲道}

\Emph{\BibleRef{弟前 3:8-10}}\par

\Emph{\BibleQuote{同样，执事也应庄重，不一口两舌，不贪杯，不贪不义之财；以纯洁的良心持守信德的奥迹。这些人也当先受考验，若无可指摘，然后才可作执事。}}\par

\BibleQuote{同样，执事。}也就是说，他们应具备与主教相同的品质。这些品质是什么呢？无可指摘、节制、好客、忍耐、不争斗、不贪财。当他说\BibleQuote{同样}时，指的是这一点，从他随后说的\BibleQuote{庄重，不一口两舌}可以明显看出，即不虚伪或欺骗。因为没有什么比欺骗更使人卑贱，在教会中，没有什么比不真诚更有害。\BibleQuote{不贪杯，不贪不义之财；以纯洁的良心持守信德的奥迹。}这样他解释了\BibleQuote{无可指摘}的意思。此处他以另一种方式要求执事\BibleQuote{不是初入教的}，当他说\BibleQuote{这些人也当先受考验}时，这里的\BibleQuote{也}是连接词，与先前论主教的论述相连，因为中间没有插入其他内容。\BibleQuote{不是初入教的}在此处的原因也相同。因为如果一个刚买来的奴隶，在未经过长期考验证明其品性之前，家里都不敢将任何事托付给他；那么，一个人从外教状态进入天主的教会，立刻被置于显要职位，岂不是荒谬的吗？\par

11节。\BibleRef{弟前 3:11} \BibleQuote{同样，女执事也应庄重，不毁谤，有节制，凡事忠心。}\par

有些人以为这是泛论妇女，其实不然；为何他要插入与主题无关的妇女问题呢？他指的是担任女执事职位的人。\par

12节。\BibleRef{弟前 3:12} \BibleQuote{执事应当只作一个妻子的丈夫。}\par

因此，这必须理解为论女执事。因为这一职务在教会中是必要、有用且尊贵的。请注意，他要求执事具备与主教同样的品德；虽然他们等级不同，但必须同样无可指摘，同样纯洁。\par

\BibleQuote{好好管理儿女和自己的家。}\par

13节。\BibleRef{弟前 3:13} \BibleQuote{因为善作执事的，自己就得到美好的地步，并且在基督耶稣里的信德上大有胆量。}\par

各处都要求他们好好管理儿女，以免因他们的不良行为使他人跌倒。\par

\BibleQuote{善作执事的，自己就得到美好的地步，}即晋升，\BibleQuote{并且在基督耶稣里的信德上大有胆量。}好像说，在较低级别上保持警惕的人，不久便会升至更高级别。\par

14至15节。\BibleRef{弟前 3:14-15} \BibleQuote{我指望快到你那里去，所以先将这些事写给你。倘若我耽延日久，你也可以知道在天主的家中当怎样行；这家就是永生天主的教会，真理的柱石和根基。}\par

为使弟茂德不致因接到这些指示而沮丧，他说：我这样写，并非我不来了，我确实要来；但若我耽搁了，你也不要忧愁。他这样写，是为了防止弟茂德灰心；对其他人而言，则是激励他们更加恳切。因为他即使只是许诺要临在，也会产生巨大影响。也不要觉得奇怪，虽然他藉着圣神预知万物，却对这事未知，只说\BibleQuote{我指望}要来，\BibleQuote{倘若我耽延}，表不确定。因为他受圣神引领，不凭自己意愿行事，所以对此事自然不确定。\par

\BibleQuote{你也可以知道在天主的家中当怎样行；这家就是永生天主的教会，真理的柱石和根基。}不像那犹太人的家。因为正是教会维护信德和圣言的宣讲。真理是教会的柱石和根基。\par

16节。\BibleRef{弟前 3:16} \BibleQuote{大哉！敬虔的奥迹，无人不以为然：天主（祂）在肉身显现，在圣神内被证明为义。}\par

此处他论及为我们的救恩计划。不要给我提铃铛、至圣所、大司祭。教会是世界的柱石。试想这奥迹，你必肃然起敬：它确实是\BibleQuote{大哉的奥迹}，\BibleQuote{敬虔的奥迹}，\BibleQuote{无人不以为然}，即毫无疑问。因为在指示司祭时，他并未要求任何像《肋未纪》中所记的事物，而是将一切引向另一位，说：\BibleQuote{天主在肉身显现。}造物主以肉身显现。\BibleQuote{在圣神内被证明为义。}正如经上说：\BibleQuote{智慧由其子女证明为义}，或因为祂没有行诡诈，如先知说：\BibleQuote{因他未行强暴，口中也没有诡诈。} \BibleRef{依 53:9} \BibleQuote{被天使看见。}因此天使与我们一同看见了天主子，之前他们并未见过祂。这奥迹真是伟大，真是伟大！\BibleQuote{被传于外邦，被世人信服。}祂被普世听闻和相信，正如先知预言的：\BibleQuote{他们的声音传遍天下。} \BibleRef{咏 18:4} 不要以为这些话只是空言，并非如此，它们充满隐藏的实义。\BibleQuote{被接在荣耀里。}祂升上云端。经上说：\BibleQuote{这离开你们被接升天的耶稣，你们见他怎样往天上去，他还要怎样来。} \BibleRef{宗 1:11}\par

真福保禄的谨慎值得注意。当他劝勉执事避免饮酒过量时，并未说：\Quote{不可醉酒}，而是说连\Quote{不能沉湎于酒}也不可。这是恰当的告诫；因为如果那些在圣殿供职的人根本不尝酒，那么这些人更不应如此。因为酒使人头脑混乱，即使不导致醉酒，也会削弱灵魂的能力，松弛其坚定。\par

他称为我们而行的救恩计划为\Quote{奥迹}，这称呼恰如其分，因为它并非显明于众人，甚至不显明于天使，因为它\BibleQuote{藉着教会显明}\BibleRef{弗 3:10}。因此他说：\Quote{毫无疑问，这奥迹是伟大的}。确实伟大。因为天主成了人，人也成了天主。一个人被看见而无罪！一个人被接升天，被传扬于世界！天使与我们一同看见了他。这确实是奥迹！因此，我们不要暴露这奥迹，不要到处宣扬它，而要以相称于这奥迹的方式生活。那些受托付保管奥迹的人是伟大的。若国王向我们透露一个秘密，我们视为恩惠。但天主已将他的奥迹托付给我们，而我们却忘恩负义，仿佛没有领受最大的恩惠。我们对此恩典的麻木不仁应令我们战栗。何以人人都知道的还是奥迹？首先，并非人人都知道；而且此前他们也不知道，但现在它被显明了。\par

\Emph{道德}。因此，在持守这奥迹上，让我们忠于所托。他托付给我们如此伟大的奥迹，我们却连自己的钱财也不信托他，尽管他吩咐我们把财富存在他那里，那里无人能夺去，既没有虫子也没有贼能毁坏。他应许偿还百倍，我们却不听从。然而在这里，若我们托付别人一笔存款，我们不会额外收回什么，反而若存款被归还就已感恩。若贼偷了它，他说：把账算在我身上；我不对你说，贼偷了，或被虫子蛀了。他在这里偿还百倍，再加上永生，但我们却不把财宝存在那里！\Quote{但是}，你说，\Quote{他还还得慢}。这正证明他恩赐的伟大，因为他不在今世偿还；或者说，他甚至在这里就偿还了百倍。保禄不是在这里留下了他的工具，伯多禄留下了他的渔竿和鱼钩，玛窦留下了他的税关吗？全世界不是为他们比为君王更敞开吗？万有不都伏在他们脚下吗？他们不是被立为统治者、为主人吗？人们不是将生命交在他们手中，完全信赖他们的谋略，并为他们服务吗？甚至现在，我们不也看到许多类似的事吗？许多贫穷卑微的人，仅操持着铁锹，几乎没有足够的必需食物，仅因隐修士的名声，就超越众人，被君王所尊敬。\par

这些事不算什么吗？那么，要想到这些只是附加的，本金存留在来生。轻看财富，你若想拥有财富；你若想真正富有，就变贫穷。因为天主的道理就是这样似非而是。他不愿你自己操心致富，而是藉他的恩宠。他说：把这些事交给我；你要关心属灵的事，好认识我的大能。逃避财富加给人的那奴役的轭。只要你紧抓财富，你就是贫穷的；当你轻看它时，你便双倍富有，因为这样的事要从各方涌向你，而且你不会缺少众人所缺少的任何东西。因为真正富有不是拥有许多，而是需要很少。君王只要有所欠缺，就与穷人无异。因为贫穷就是需要别人；按此论证，君王是贫穷的，因为他需要臣民。但那被钉死于世界的人却非如此；他什么也不缺，因为他的双手足以维持生计，如保禄所说：\BibleQuote{这些手供应了我的需要，以及与我同工的人的需要}\BibleRef{宗 20:34}。这位说\BibleQuote{似乎一无所有，却拥有一切}\BibleRef{格后 6:20}的就是他。他就是被吕斯特拉居民视为神的那位。你若想获得世俗之物，就追求天上之事；你若想享受世间之物，就轻看它们。因为他说：\BibleQuote{你们先该寻求天主的国和它的义德，这一切自会加给你们}\BibleRef{玛 6:33}。\par

你为何崇拜这些琐事？为何渴求无价值之物？人穷困多久？乞讨多久？举目望天，思念那里的财富，嘲笑黄金；思考它多么无用；它的享乐只持续今生，而与永恒相比，今生如同沙粒或水滴之于无边的海洋。这财富不是产业，不是财产，而是供使用的借贷。因为当你死时，无论情愿与否，你所拥有的一切都归别人，他们又转给别人，别人再转给别人。因为我们都是旅居者；房屋的承租人也许更真正是它的主人，因为主人死了，承租人活着，仍然享用房屋。若后者承租，前者也可说是承租：因为他建造了它，劳苦经营，装修好了。事实上，财产只是一个词：实际上我们所有人都是他人之物的拥有者。只有那些我们已经先行送到另一个世界的事物才真正属于我们。此世的财物不属于我们；我们只有终身的权益，或者说它们在我们有生之年就消逝了。只有灵魂的德行真正属于我们，如施舍和爱德。世俗之物，即使外教人也称之为外在之物，因为它们在我们之外。但让我们把它们变成内在的。因为我们离开这里时不能带走财富，却可以带走施舍。然而，让我们宁可把它们先送上去，好叫它们为我们预备永恒的居所\BibleRef{路 16:9}。\par

财物因使用而得名，而非因所有权；它们并非我们的，产业也非财产而是借贷。因为每份产业曾有过多少主人，还将有多少主人！有一句明智的谚语（民间谚语若含有智慧，也不应轻视）：\Quote{田地啊，你曾属于多少人，又将属于多少人？}这也是我们对房屋和一切财物应当说的。只有德行能随我们离去，并陪伴我们前往上界。那么，让我们舍弃并熄灭对财富的贪爱，好能在我们内燃起对天上事物的热爱。这两种情感不能占有同一个灵魂。因为经上说：\BibleQuote{他必要恨这一个而爱那一个，或是依附这一个而轻忽那一个。}\BibleRef{玛 6:24}你看见一个人带着长长的随从队伍，在街上开道，身穿丝绸衣服，高高骑在马上，硬挺着脖子吗？不要畏惧，而要微笑。如同我们看见孩子扮演国王时发笑，也当嘲笑他的排场，因为这不比孩子的游戏更好，甚至也不那么快意，因为缺乏孩童的那份纯真和单纯。在孩子那是欢笑和愉悦，在这里则是一个人变得可笑和可鄙。\par

光荣天主，祂使你免于这种戏剧性的炫耀。因为，如果你愿意，尽管你的地位卑微，你也能比那个坐在战车上高高在上的人更高贵。为什么呢？因为，虽然他的身体略略离开了地面，他的灵魂却紧贴地面，因为他说：\BibleQuote{我的力量，}他说，\BibleQuote{粘在我的肉上。}\BibleRef{咏 101:6}而你却以你的精神行走在天上。即使他有众多随从开道，又如何呢？他因此得的荣耀比他的马更多吗？驱赶人在前面为一只野兽清道是多么荒谬！那么骑在马上又是什么荣耀？那是连他的奴隶也分享的荣耀。然而有些人对此如此虚荣，即使不需要，也要让人牵着马跟在后面。还有比这更愚蠢的吗？想靠马匹、昂贵的衣服、随从来显耀自己！还有什么比由马匹和仆从构成的荣耀更可鄙的呢？你有德行吗？就不要用这些区分。让你自身有装饰。不要因他人的在场而欠下荣耀。这种荣誉最邪恶、腐败、卑劣的人也能得到；实际上所有富人都能。演员和舞者也可以骑马，并有仆人在前奔跑，但他们仍然是演员和舞者。他们的马匹和随从并未为他们赢得尊敬。因为当灵魂的美善缺乏时，这些外物的添加是多余且虚浮的。正如一面墙脆弱或身体失调时，无论你加上什么，它仍然不健全而衰败；同样，在这种情况下，灵魂保持不变，从外在事物中得不到任何益处，即使那人佩戴上千件金饰。因此，我们不要为这些事忧虑。让我们从现世事物中抽身，追求更大的、属灵的殊荣，这将使我们真正成为可敬的对象，使我们也能获得来世的福分，藉着我们的主耶稣基督的恩宠和慈爱，愿光荣归于祂，等等。\par

\SourceRecord{Translated by Philip Schaff.；Nicene and Post-Nicene Fathers, First Series；Vol. 13.；Edited by Philip Schaff.；Buffalo, NY: Christian Literature Publishing Co.,；1889.；<http://www.newadvent.org/fathers/230611.htm>.}

% structure: p0001 source=h1 target=section reason=page-title

\section{\WorkTitle{弟茂德前书第十二篇讲道词}}

\BibleRef{弟前 4:1-3}\par

\BibleQuote{现今圣神明明地说：在末期，有些人将离弃信德，听从诱惑人的鬼神和魔鬼的道理；他们以伪善讲谎话，良心如同被烙铁烙过；禁止嫁娶，又禁戒食物，就是天主所造、让相信并认识真理的人以感恩的心领受的。}\par

凡坚守信德的人，如同系在安全的锚上；凡离弃信德的人，则无处可安息；他们四处漂泊，最终被卷入灭亡的深渊。保禄先前已指出这一点，说有些人\Quote{已经于信德遭遇了船难}；现在他又说：\Quote{现今圣神明明地说，在末期，有些人将离弃信德，听从诱惑人的鬼神。}这是指摩尼派（\OriginalTerm{摩尼派}{Manichæans}）、克己派（\OriginalTerm{克己派}{Encratites}）、马西翁派（\OriginalTerm{马西翁派}{Marcionites}）以及他们的整个宗派，他们将离弃信德。你们看见吗？离弃信德是随后一切邪恶的原因！\par

但什么是\Quote{明明地}呢？就是清楚、明确、毫无疑问。他说：不要惊奇，有些人离弃信德却仍拘守犹太教。将来有一个时候，就连那些已分享信德的人也会陷入更严重的错误，不仅限于食物，而且涉及婚姻及其他类似的事，引入极其有害的观念。这里不是指犹太人（因为\Quote{末期}和\Quote{离弃信德}不适用于他们），而是指摩尼派（\OriginalTerm{摩尼派}{Manichees}）和这些教派的创始人。他非常恰当地称他们为\Quote{诱惑人的鬼神}，因为正是这些鬼神驱使他们说出这样的话。\Quote{以伪善讲谎话}，这说明他们不是由于无知而说不真实的话，而是像演戏一样，明知真理，却\Quote{良心如同被烙铁烙过}，也就是说，他们过着邪恶的生活。\par

但为什么他只提这些异端者？基督先前曾说：\Quote{绊倒人的事是免不了的}\BibleRef{玛 18:7}，他在撒种和稗子生长的比喻中也预言了同样的事。但此处请钦佩保禄的先知之恩，他在这些异端者出现之前的时代，就指明了具体的时间。就如他曾说：不要惊奇，在初期信德建立时，有些人试图引入这些有害的道理；因为当信德已确立一段长时间后，许多人将离弃信德。\Quote{禁止嫁娶，又禁戒食物}。那么他为何没有提及其他异端？虽然没有具体列出，但已包含在\Quote{诱惑人的鬼神和魔鬼的道理}这些说法中。但他不愿提前将这些事灌输给人；至于已经开始的关于食物的问题，他指明了。\Quote{就是天主所造、让相信并认识真理的人以感恩的心领受的}。为什么他没有说\Quote{也被不信的人领受}？不信的人如何领受？他们用自己的规则将自己排除在外。但难道奢侈是被禁止的吗？当然是的。但为什么？如果好的事物是为被领受而受造的，为什么？因为祂造了饼，但过量是被禁止的；祂也造了酒，但过量是被禁止的；我们并非被命令要避免美味佳肴，好像它们本身不洁，而是因为它们因过量而腐蚀灵魂。\par

第4节. \BibleQuote{因为天主所造的一切都是好的，若以感恩的心领受，没有一样是可弃绝的。}\BibleRef{弟前 4:4}\par

如果它是天主所造的，它就是好的。因为经上说：\BibleQuote{都很好}\BibleRef{创 1:31}。他如此谈论可食之物，是预先驳斥那些主张未受造的物质、并声称这些食物出自那物质之人的异端。但如果它是好的，为何要\Quote{因天主的话和祈祷}而成圣呢？因为它本是不洁的，所以需要被圣化？不是的，他这里是对那些认为其中某些食物是不洁的人说的；因此他提出两个要点：第一，天主所造的没有一样是不洁的；第二，即使它变成不洁，你也有补救的方法：划十字圣号、感谢天主、光荣天主，所有不洁就消失了。那么我们能够这样洁净祭献给偶像的食物吗？如果你不知道它是祭过的，就可以。但如果你知道，却仍吃它，你便是不洁的；不是因为它是祭献给偶像的，而是因为违背了明确的命令，从而与魔鬼相通。所以它并非本性不洁，而是因为你的故意悖逆而成为不洁的。那么，猪肉是不洁的吗？绝不是，当它被以感恩的心、并划十字圣号而领受时，没有任何东西是不洁的。你对天主不知感恩的态度才是不洁的。\par

第6节. \BibleQuote{你若将这些事提醒弟兄们，便是耶稣基督的好执事，在信仰的话语和你向来所跟随的好道理上得了栽培。}\BibleRef{弟前 4:6}\par

这里所指的\Quote{这些事}是什么？就是他先前提到的：\Quote{伟大的奥秘}，并说禁戒食物是魔鬼的道理，以及食物\Quote{因天主的话和祈祷而成洁}。\par

第7节. \BibleQuote{但要弃绝世俗和老妇的无稽之谈，在虔敬上操练自己。}\BibleRef{弟前 4:7}\par

他说：\Quote{提醒他们}；这里你见到的是完全的俯就，而非权威：他没有说\Quote{命令}或\Quote{吩咐}，而是\Quote{提醒他们}。就是说，把这些事作为建议提出，并以此与他们谈论信德，\Quote{被栽培}，他说，意味着持续专注于这些事。\par

因为我们每天为自己预备身体的食粮，同样，他意指我们要持续不断地接受关于信德的言论，并常以它们为滋养。什么是\Quote{被栽培}？就是反复思想它们；常常专注于同样的事，常常实行同样的事，因为它们所供应的不是普通的营养。\par

\BibleQuote{但拒绝凡俗的、老妇的荒诞传说。}这里指的是犹太人的传统，他称它们为\BibleQuote{荒诞传说}，或因其虚假，或因其不合时宜。因为合时宜的就有用，不合时宜的不仅无用而且有害。假设一个成年人由乳母哺乳，岂不是可笑，因为不合时宜？他称它们为\BibleQuote{凡俗的、老妇的荒诞传说}，部分是由于它们的陈旧，部分是由于它们妨碍信德。因为使那些已超脱这些事的灵魂再受恐惧的束缚，是亵渎的命令。\newline{}\BibleQuote{要操练自己达到虔敬。}也就是说，达到纯洁的信德和道德的生活；这就是虔敬。所以我们确实需要\Quote{操练}。\par

第8节：\BibleQuote{因为身体的操练益处甚少。}有人将此联系于禁食；但不可有这种想法！因为禁食不是身体上的操练，而是属神的操练。若是身体上的，它就会滋养身体，但禁食反而消耗身体使之消瘦，所以不是身体上的。因此他不是在谈论身体的纪律。我们所需要的，乃是灵魂的操练。因为身体的操练没有益处，或许对身体有点好处，但虔敬的操练却结出果实和益处，既在今世，也在来世。\newline{}\BibleRef{弟前 4:8}\par

\BibleQuote{这话是确实的}，就是说，虔敬在今世和来世都有益处是真实的。请注意他处处引入这一点，他不需要论证，只是简单地宣告，因为他是对弟茂德说话。\par

那么，即使在今世，我们也有美好的希望吗？因为那自知无罪、并在善事上结果实的人，甚至今世就喜乐；而恶人则在今世和来世都受惩罚。他活在永久的恐惧中，不能坦然面对任何人，脸色苍白，颤抖不安，充满焦虑。欺诈者和盗贼不也是这样吗？他们甚至对自己拥有的东西也不满足。谋杀者和奸淫者的生活不是最悲惨的吗？他们连看太阳都带着怀疑。这能称为生命吗？不，这简直是可怕死亡！\par

第10节：\BibleQuote{因此我们劳苦并受辱骂，是因为我们信赖永生的天主，祂是众人的救主，尤其是信徒的救主。}\newline{}\BibleRef{弟前 4:10}\par

这意思就是说，我们为何克苦自己，除非我们期待将来的福乐？我们忍受了那么多恶事，承受了那么多辱骂，遭受了那么多侮辱和毁谤，以及如此众多灾难，难道是徒然的吗？因为我们若不信赖永生的天主，我们为何承受这些事？但如果天主在此是未信者的救主，那么祂将来更是信徒的救主。他说的救恩是什么？是未来的拯救吗？他说：\BibleQuote{祂是众人的救主，尤其是信徒的救主。}目前他指的是今世的救恩。但是祂怎样是信徒的救主？若非如此，他们早已被消灭了，因为所有人都攻击他们。他在这里劝勉他忍受危险，以便有天主作他的救主，他就不致气馁，也不需要别人的帮助，而是甘愿并坚忍地承受一切。甚至那些急切追求世俗利益的人，因着获利的希望，也乐意承担劳苦的事业。\par

那么这是末期了。因为他说：\BibleQuote{在后期，有些人将背离信德。}\BibleQuote{禁止嫁娶。}难道我们禁止嫁娶吗？决不是！我们不禁止那些愿意结婚的人，而是劝勉那些不愿意结婚的人守贞。禁止是一回事，让人自由选择是另一回事。禁止的人是一劳永逸地禁止，而推荐贞洁为更高境界的人并不禁止婚姻，只是他更推崇贞洁。\par

\BibleQuote{禁止嫁娶}，\BibleQuote{并命令戒食那些天主所造、让信者并认识真理的人以感谢的心领受的食物。}说\BibleQuote{认识真理}很好。那么先前的事物是预像。因为没有任何东西在本性上是不洁的，它之所以不洁是由于取用者的良心。而禁止那么多肉食的目的是什么？为了约束过度的奢侈。但如果他说\Quote{为奢侈的缘故不要吃}，恐怕人们难以接受。因此他们被律法的必要性所束缚，从而因更大的畏惧而禁戒。鱼并没有被禁止，虽然它明显比猪肉更不洁。但他们可能从梅瑟的话中学到了奢侈是多么有害：\BibleQuote{耶叔仑养肥了，就踢跳。}\BibleRef{申 32:15}这些禁令的另一个原因可能是，他们因缺乏别的食物，不得不屠宰牛羊；因此他禁止其他东西，因为阿匹斯和牛犊是可憎的、忘恩负义的、污秽的、不洁的。\par

\BibleQuote{你要提醒他们这些事，要默想这些事}，因为\BibleQuote{在信德的话语和美好的教训中得到培养}这句话暗示他不仅应当向别人推荐这些事，自己也当实践。因为他他说：\BibleQuote{在信德的话语和美好的教训中得到培养，你已经达到了这些。但你要拒绝凡俗的、老妇的荒诞传说。}为什么他不说\Quote{远离}，而说\Quote{拒绝}？他以此暗示应该完全摒弃它们。\par

\Quote{但你要在虔诚上操练自己，}就是说，在于纯洁的生活和最德行端正的品行。操练自己的人，即使不在比赛季节，也总是如同在竞赛一般，实行斋戒，忍受一切劳苦，常怀忧虑，忍受许多辛劳。\Quote{操练自己，}他说，\Quote{在于虔诚；因为身体的操练益处甚少，但虔诚在一切事上都有益处，有今生和来生的应许。}有人问，为什么他提到身体的操练？为要通过比较显示另一者的优越，因为前者虽伴随许多劳苦，却没有实质的益处，而后者却有持久且丰厚的善。就如他吩咐妇女\Quote{应当打扮自己，不是用编发、金饰或昂贵的衣裳，而是用善行，这正合乎自称虔诚的妇女。}\BibleRef{弟前 2:9}\par

\Emph{道德训诲}。\Quote{这话是确实的，值得完全接纳。因为我们劳苦受辱，正是为此。}难道\Person{保禄}曾受辱，而你不耐烦吗？难道\Person{保禄}劳苦，而你想过安逸生活吗？倘若他过安逸生活，绝不可能获得如此大的福乐。因为即使属世的财物，既不确定又易朽坏，人们若不劳苦费力，尚且不能获得，何况属灵的？有人会说，但有些人继承了它们。然而即使继承了，若没有劳苦、关心和麻烦，也不会比那些获得它们的人更容易保存。我不必说，许多劳苦忍耐的人，在即将进港时，却被逆风吹翻了希望的船，在即将拥有时遭难。但我们这里没有这样的事。因为许下的是天主，而\Quote{这望德不会叫人蒙羞。}\BibleRef{罗 5:5}你们这些熟悉世俗事务的人，难道不知道有多少人，在无限辛劳之后，并未享受劳苦的果实，或因早逝，或遭不幸，或染疾病，或受诬告，或在人类各种变故中，因其他原因而被迫空手而去吗？\par

但你看那些幸运的人，有人说，他们用很少的劳力就获得了世上的好东西？什么好东西？金钱、房屋、那么多亩地、成群的仆人、成堆的金银？你能把这些称为好东西，而不羞愧得无地自容吗？一个被召追求天上智慧的人，却贪恋世俗事物，并称它们为\Quote{好东西}，那些毫无价值的东西！如果这些东西是好的，那么拥有它们的人就必须被称为好人。因为拥有好东西的人，难道不是好人吗？但当拥有这些东西的人犯欺诈和掠夺时，我们还会称他们为好人吗？如果财富是好的，但通过攫取而增加，那么它越增加，其拥有者就越被视为好人。那么，攫取的人会是好人吗？但如果财富是好的，且通过攫取而增加，一个人越攫取，他就越好。这岂不显然是矛盾的吗？但假设财富不是通过不正当手段获得的。这怎么可能呢？贪财是那么毁灭性的情欲，以至于不靠不义致富是不可能的。这正是\Person{基督}所说的：\Quote{你们要用不义的钱财结交朋友。}\BibleRef{路 16:19}但如果他继承了父亲的遗产呢？那么他接受的是通过不义聚敛的东西。因为他的祖先不是从亚当那里继承财富的，而是从许多前人之中，大概有人不义地夺取并享用了别人的财物。那么，他说，\Person{亚巴郎}拥有不义的财富吗？还有\Person{约伯}，那个无可指责、正义、忠诚的人，他\Quote{敬畏天主，远离邪恶}？他们的财富不在于金银，也不在于房屋，而在于牲畜。此外，他是天主所祝福的。那卷书的作者在叙述那位有福之人的遭遇时，提到他丢失了骆驼、母马和驴，但没有说到有金银财宝被夺走。\Person{亚巴郎}的财富也是他的仆从。那么，他不是买来的吗？不，因为\WorkTitle{圣经}恰恰说到，那三百一十八人是在他家里出生的。\BibleRef{创 19:14}他也有羊和牛。那么，他送给黎贝加的金子是从哪里来的呢？\BibleRef{创 24:22}是从他在埃及收到的礼物，没有暴力或不当手段。\par

那么，告诉我，你是从哪里得到财富的？你从谁那里得到它，而传给你的人又从谁那里得到？从他的父亲和他的祖父。但是，你能上溯许多代，证明这获得是正义的吗？不可能。它的根源和起源必定是不义。为什么？因为天主在起初并没有使一人富、一人穷。后来也没有取出金银财宝给一个人看，而禁止另一个人寻找；而是把土地同等地留给所有人。那么，如果它是公共的，你为什么有那么多亩地，而你的邻居却没有一点？是从我父亲传给我的。那么给他的是谁？是他的祖先。但你必须往回找到最初的所有者。\Person{雅各伯}有财富，但那是他劳动所得的报酬。\par

但我不打算过于紧逼这个论点。让你的财富以公正的方式获得，且没有抢劫。因为你对你父亲的贪婪行为不负责任。你的财富可能源于抢劫，但你并非掠夺者。或者即使他并非通过抢劫获得，他的金子是从某处地下挖出来的。那又怎样？那么财富因而就是善的吗？绝非如此。同时，他说，如果拥有者不贪婪，它就不坏；如果分施给穷人，它就不坏；否则它就是坏的，是引人犯罪的。他说：\Quote{但如果他不行恶，虽然也不做好事，这不算坏。}确实如此。但惟独你拥有主的财产，惟独你享受公共之物，这难道不是一种恶吗？难道不是\BibleQuote{大地属于天主，及其中所充满的}吗？那么，如果我们的财产属于同一位共同的主，它们也属于我们的同伴仆人。一位主的财产都是共有的。我们在大家族中难道没有看到这是既定规则吗？每个人都得到相同份额的供给，因为供给来自他们主的宝库。主人的房子向所有人敞开。国王的财产都是共有的，如城市、市场和公共大道。我们都平等地分享它们。\par

请注意天主智慧的安排。祂为使人类羞愧，使某些事物成为公有，如太阳、空气、大地和水，天、大海、光、星辰；其利益平等地分施给所有人，如同兄弟。我们都被造成同样的眼睛、同样的身体、同样的灵魂、各方面同样的结构，万物都出于大地，众人皆出于一人，且都住在同一居所。但这些不足以使我们羞愧。于是祂又使其他事物（如我们所说）成为公有，如浴场、城市、市场、大道。并且请注意，关于公有之物没有纷争，一切都和平。但当一个人试图占有某物，使之成为他自己的，这时纷争就引入了，仿佛大自然本身感到愤慨，因为当天主以各种方式使我们聚集在一起时，我们却渴望通过占有事物、使用那些冷酷的字眼\Quote{我的和你的}来分裂和分离我们自己。于是就有了纷争和不安宁。但在没有这种情况的地方，就不会产生争吵或纷争。因此，这种状态才是我们的遗产，更合乎本性。为什么人们在市场上从不争论？难道不是因为它属于众人吗？但关于一栋房子、关于财产，人们总是争论不休。必需品被公共地摆在我们面前；但即使在最小的事情上，我们也不遵守共有。然而那些更大的东西，祂已自由地向所有人开放，好让我们由此受教，将这些较低级的事物共有。然而尽管如此，我们仍未受教。\par

但如我所说，富有的人怎么能是善人？当他分施财富时，他是善的，所以他是善的，在他不再拥有它的时候，当他把它给别人时；但当他保留它时，他并不善。那么，那些保留它使人邪恶、舍弃它使人变善的东西，怎么能是善的呢？因此，不是拥有财富，而是不拥有财富，才使人显得善。所以财富不是一种善。但如果你在能够接受时却不接受，你又是善的。\par

那么，如果我们拥有它时，我们分施给别人；或者当它被提供给我们时，我们拒绝它；而如果我们接受或获得它时，我们并不善，这样东西本身怎么能是善呢？所以不要称它为善。你不拥有它，是因为你认为它是善，因为你渴望拥有它。洁净你的心灵，纠正你的判断，然后你就会是善的。学习什么是真正的善。它们是什么？美德与仁爱。这些而不是那些，才是真正的善。根据这一规则，你越有爱德，你就越被认为善。但如果你富有，你就不再是善的。因此，让我们成为这样的善，好让我们真正地善，并能在耶稣基督内获得未来的善，与祂一同，等等。\par

\SourceRecord{Translated by Philip Schaff.；Nicene and Post-Nicene Fathers, First Series；Vol. 13.；Edited by Philip Schaff.；Buffalo, NY: Christian Literature Publishing Co.,；1889.；<http://www.newadvent.org/fathers/230612.htm>.}

% structure: p0001 source=h1 target=section reason=page-title

\section{论弟茂德前书第十三篇讲道}

\Emph{\BibleRef{弟前 4:11}}\par

\BibleQuote{你要命令和教导这些事。不要让人小看你年轻；但要在言语、行为、爱德、精神、信德和纯洁上，成为信徒的榜样。直到我来到以前，你要专心于宣读、劝勉和教导。不要疏忽你内在的神恩，这神恩曾藉预言并在长老团的覆手下赐给了你。}\par

在某些情况下需要命令，在另一些情况下需要教导；因此，如果你在需要教导的地方命令，就会变得可笑。同样，如果你应当在命令的地方教导，你也会遭受同样的责备。例如，不该教导一个人不要作恶，而应当命令；以全部权柄禁止它。不该信奉犹太主义，这应当是一个命令，但当你引导人们舍弃财产、宣认守贞，或谈论信德时，则需要教导。因此保禄提及两者：\BibleQuote{你要命令和教导。}当一个人使用护身符或做类似的事，明知它是错的，他只需一个命令；但若他无知地去做，就应当教导他的错误。\BibleQuote{不要让人小看你年轻。}\par

请注意，司铎应当命令并权威地说话，而非总是教导。但由于一种常见的偏见，年轻人容易被轻视，因此他说：\BibleQuote{不要让人小看你年轻。}因为教师不应受轻蔑。但若他不应被轻蔑，那么温良与节制又何在呢？事实上，他个人所受的轻蔑应当承受；因为教导因忍耐而受称赞。但涉及他人时则不然；因为这并非温良，而是冷漠。若有人报复侮辱、恶言及对自身的伤害，你正当地责备他。但当涉及他人的得救时，要命令并以权威干预。这不是讲节制的时候，而是讲权威的时候，以免公共利益受损。他根据情况要求其中之一。不要因你的年轻而被人轻视。因为只要你的生活足以弥补，就不会因年轻而被轻视，反而更受钦佩。因此他接着说：\par

\BibleQuote{但你要在言语、行为、爱德、信德和纯洁上，成为信徒的榜样。}在一切事上显示自己为善行的榜样：也就是说，你自身要成为基督徒生活的模范，作为摆在他人面前的榜样，作为活的法律，作为善生活的规则与标准，因为教师应当如此。\Quote{在言语上，}使他能流利地说话，\Quote{在行为上、在爱德上、在信德上、在真正的纯洁上、在节制上。}\par

\BibleQuote{直到我来以前，你要专心于宣读、劝勉和教导。}\par

就连弟茂德也被命令专心阅读。那么我们当受教导，不可忽视对神圣著作的研习。再者，请注意他说：\BibleQuote{直到我来以前。}留意他如何安慰他，因为当他与他分离时，他仿佛成了孤儿，自然需要这样的安慰。他说：\BibleQuote{直到我来以前，你要专心阅读神圣著作，彼此劝勉，教导众人。}\par

\BibleQuote{不要疏忽你内在的神恩，这神恩曾藉预言赐给了你。}\par

此处他称教导为预言。\par

\BibleQuote{藉着长老团的覆手。}此处他说的不是长老，而是主教。因为不能认为长老祝圣了主教。\par

\BibleRef{弟前 4:15} \BibleQuote{你要思考这些事，全心投入其中。}\par

请注意他多么频繁地就同一件事给他劝告，从而表明教师尤其应当注意这些要点。\par

\BibleRef{弟前 4:16} \BibleQuote{你要注意自己和自己的教训；在这些事上恒心。}也就是说，注意你自己，也教导他人。\par

\BibleQuote{因为这样行，既能救你自己，也能救听你的人。}\par

他说得好：\BibleQuote{你也能救你自己。}因为那\Quote{在健全道理的言语上得到滋养}的，首先自己从中获益。在劝诫他人时，他自己也受到内心的刺痛。这些话不仅是对弟茂德说的，也是对所有人说的。如果这样的劝告是对那复活死人的人说的，那对我们又该说什么呢？基督也说明了教师的职责，祂说：\BibleQuote{天国好像一个家主，从他的宝库里取出新和旧的东西。} \BibleRef{玛 13:52}有福的保禄也给出同样的劝告：\BibleQuote{我们因圣经所生的忍耐和安慰，可以怀有希望。} \BibleRef{罗 15:4}他在这方面超越众人，在加玛里耳脚前受教于祖传的法律，此后自然乐于阅读；因为劝诫他人的人，自己必先遵从他给出的劝告。因此我们发现他不断诉诸先知们的见证，探究他们的著作。保禄于是致力于阅读，因为从圣经中获益匪浅。但我们却懒惰，漫不经心、漠不关心地听。我们该受何等惩罚呢！\par

\BibleQuote{使你的进步显明在众人面前。}\par

这样，他愿他在这一方面也显得伟大而可钦佩，表明这对他来说仍是必要的，因为他希望他的\Quote{进步}不仅显现在生活中，也显现在教导的道理上。\par

\BibleRef{弟前 5:1} \BibleQuote{不要斥责老年人。}\par

他现在是在说圣秩吗？我想不是，而是指任何年长的人。那么如果他需要纠正呢？不要斥责他，而要像对待犯错的父亲一样对他讲话。\par

\BibleRef{弟前 5:1} \BibleQuote{待年老的妇人如同母亲，待年轻的男子如同弟兄，待年轻的女子如同姐妹，凡事都当纯洁。}\par

斥责本身带有冒犯性，尤其是针对长者，又出自年轻人之口，就有三重冒失的表现。因此他要通过方式和温和来缓和它。因为如果能注意这一点，不冒犯地纠正是可能的；这需要极大的谨慎，但可以做到。\par

\Quote{以弟兄待青年男子。}他为何在此也推荐这点呢？考虑到青年男子天生的高昂心志，因此也应适度缓和规劝他们。\par

\Quote{以姊妹待青年妇女}；他又加上了\Quote{及一切纯洁。}他的意思是，不要只说避免与她们有罪恶的交往。连丝毫的嫌疑都不应有。因为与青年妇女亲密总是可疑的，然而一位主教又不能总是避免，他加上这些话表明，在这种亲密中需要\Quote{一切纯洁}。但保禄给弟茂德这条建议吗？是的，他说，因为我通过他是在对整个世界说话。但若弟茂德尚被这样劝诫，别人就该想想自己该有什么样的品行，不让那些想毁谤的人有任何怀疑的根据，任何借口的影子。\par

\Quote{要尊敬那些真为寡妇的寡妇。} \BibleRef{弟前 5:3}\par

为何他对守贞只字不提，也没有命令我们尊敬贞女？或许那时还没有宣认此身份的人，或是他们已从其中堕落了。他说：\Quote{因为有些已转去随从撒殚。} \BibleRef{弟前 5:15} 因为一个女人可能失去丈夫，却并非真正是寡妇。正如要成为贞女，单是不结婚并不够，还需要许多其他条件，如无可指摘和坚忍；同样，失去丈夫并不构成寡妇，而是需要忍耐、贞洁和与所有男人隔离。这样的寡妇他合理地命令我们尊敬，不如说是供养。因为他们需要支持，孤苦无依，没有丈夫为他们站出来。他们的处境在众人眼中显得卑贱和不祥。因此他希望她们从司铎那里得到更大的尊敬，更因她们是配得的。\par

\Quote{但若寡妇有子女或孙子女，让她们先学着在家表现孝爱，并报答父母。} \BibleRef{弟前 5:4}\par

请留意保禄的明智；他经常从人的考虑来劝勉。他并未在此设立什么伟大崇高的动机，而是易于明白的：\Quote{报答父母。}如何报答？因为抚养和教育他们。犹如他在说，你们从他们那里领受了极大的关怀。他们已经离世。你们无法回报他们。因为你们并未生养他们。在子孙身上报答他们，透过子女偿还债务。\Quote{让她们先学着在家表现孝爱。}这里他更直接地劝勉她们行善；然后为了进一步激励她们，他加上：\par

\Quote{因为这在天主面前是美好而可悦纳的。}既然他提到了那些\Quote{真为寡妇的人}，他就宣告谁是真正的寡妇。\par

\Quote{但那真为寡妇、且孤苦无依的，她寄望于天主，并昼夜恒常于恳求和祈祷。但那生活在享乐中的，虽活着也是死的。} \BibleRef{弟前 5:5}\par

那作为寡妇而并未选择世俗生活的，才是真寡妇；她按应当的寄望于天主，并昼夜恒切于祈祷，才是真寡妇。并非有子女的就不是真寡妇。因为他称赞那按应当养育子女的人。但若有人没有子女，他说，她是孤苦无依的，他安慰她，说她是真正的寡妇，她不仅失去了丈夫的安慰，也失去了来自子女的安慰，然而她有天主代替一切。她不会因没有子女而更差，反而祂以安慰充满她的缺乏，就是她没有子女。他所说的大意是：当人说寡妇应养育子女时，不要悲伤，好像因为你们没有子女，你们的价值就因此较低。你们是真寡妇，而那生活在享乐中的，虽活着也是死的。\par

既然许多有子女的人选择守寡的状态，并非为了切断世俗生活的机会，反而为了增加它们，好使他们能更加放纵地做他们想做的事，更自由地沉溺于世俗的欲望：因此他说：\BibleQuote{她好宴乐的，虽活着也是死的。} 那么寡妇难道不应该享乐吗？当然不。如果当自然和年龄衰弱时，享乐的生活都是不允许的，反而导致死亡，即永死；那么那些过享乐生活的人们有什么可说的呢？但他有理由说：\BibleQuote{她好宴乐的，虽活着也是死的。} 但为了让你明白这点，现在让我们看看死者的状态是什么，活者的状态是什么，我们将把这样的人放在哪一类？活者实行生命的工作，那未来生命的工作，那才是真正的生命。基督已经宣告了那未来生命的工作是什么，我们应该时常忙于这些工作。\BibleQuote{你们来承受那从创世以来为你们预备了的国吧，因为我饿了，你们给了我吃的；我渴了，你们给了我喝的。} \BibleRef{玛 25:34} 活者与死者的区别，不仅在于他们看见太阳、呼吸空气，而在于他们行善。因为如果这缺失了，活者并不比死者更好。为使你懂得这点，请听死者如何也能活着。因为经上说：\BibleQuote{天主不是死人的天主，而是活人的天主。} \BibleRef{玛 22:32} 但这你又说是个谜。让我们把两者都解开。一个享乐的人，活着的时候是死的。因为他只为他肚腹而活。在他的其他感官上，他不活着。他看不见他应当看见的，听不见他应当听的，说不出他应当说的。他也不实行活者的行为。但如同一个人摊在床上，闭着眼睛，眼皮紧闭，对一切事一无所知；这个人也是如此，或者更糟。因为前者对善恶都一样无感，但后者只对恶有感觉，对善却像前者一样无感。因此他是死的。因为关于未来生命的一切都不能打动或影响他。因为无节制将他揽入怀中，好像进入一个黑暗阴森、充满一切不洁的洞穴，使他完全住在黑暗中，如同死人。因为当他的所有时间都在宴饮和醉酒中度过时，他岂不是死的，埋葬在黑暗中吗？即使在早晨当他似乎清醒时，他实际上并不清醒，因为他还没有摆脱和洗净昨日的过量，仍在渴望重复，而他的夜晚和中午在狂欢中度过，整夜和大部分早晨则在沉睡中。\par

那么他还能被算在活人中吗？谁能描述那来自奢侈的风暴，袭击他的灵魂和身体？因为如同一片持续阴沉的天空不允许阳光照射穿过，同样，奢侈和酒气的烟雾笼罩着他的大脑，仿佛那是一块岩石，并在其上投下浓雾，不让理性发挥作用，反而使醉酒者被深沉的黑暗覆盖。对于这样受影响的人，他内在的风暴有多大，喧嚣有多猛烈！正如当洪水上涨，漫过了作坊的入口，我们看到里面的所有人都混乱不堪，使用桶、壶、海绵和许多其他工具来舀水，以免水既破坏建筑又毁坏里面的一切：同样，当奢侈淹没灵魂时，它内在的理性被扰乱。已经积存的无法排出，而引进更多时，就掀起猛烈的风暴。因为不要看那快活欢欣的面容，而要审视内心，你会看到它充满深深的沮丧。如果可能将灵魂显现，用我们的肉眼看见它，奢侈者的灵魂会显得沮丧、悲伤、可怜，因消瘦而憔悴；因为身体越光滑肥胖，灵魂就越消瘦虚弱；一个人越受纵容，灵魂就越受阻碍。正如当眼睛的瞳孔有太厚的外膜覆盖时，它就不能发挥视觉能力向外看，因为光线被厚覆盖层阻挡，常常导致黑暗；同样，当身体持续饱食时，灵魂必定被粗厚包裹。但你们说，死人腐烂朽坏，从不健康的湿气中滴出液体。同样，在她 \Quote{好宴乐的} 身上，可以看到黏液、痰、鼻卡他、打嗝、呕吐、嗳气等等，由于太不雅，我姑且不提。因为奢侈的统治如此强大，它使人忍受我们甚至认为不该提及的事情。\par

但你仍然问道，当身体还在吃东西和喝东西时，它如何被分解呢？这当然不是人类生命的标志，因为无理性的生物也吃也喝。哪里灵魂死了，吃喝有什么用？死尸，即使用花衣装裹，也不会因此受益；而当容光焕发的身体包裹一个死去的灵魂时，灵魂也不会受益。因为当它全部谈论都是关于厨师、承办人和甜点师，而不发出任何虔诚的话语时，它不是死的吗？让我们思考什么人是？外教人说，人是理性的动物，会死的，有理智和知识的能力。但我们不要从他们那里接受定义，而是从哪里？从圣经。圣经在哪里给了人的定义？请听它的话。\BibleQuote{有一个人，完美正直，敬畏天主，远离邪恶。} \BibleRef{约 1:2} 这真是一个人！此外，另一位说：\BibleQuote{人有大德，仁慈的人宝贵。} \BibleRef{箴 20:6} 那些不符合这描述的人，虽然他们分有理智，且有如此大的知识能力，圣经却拒绝承认他们是人，而称他们为狗、马、蛇、狐狸、狼，以及任何更可鄙的动物。如果这就是人，那么享乐的人就不是人；因为他怎能算是人呢，他从不思想他所应当思想的事？奢侈与节制不能共存：它们互相毁灭。连外教人也说，\par

\Quote{大腹便便生不出敏锐的心智。}\par

圣经称这样的人为没有灵魂的人。\BibleQuote{我的圣神（据说）不会永远在这些人体内，因为他们是血肉。}\BibleRef{创 6:3} 但他们有灵魂，但因为灵魂在他们内是死的，祂称他们为血肉。因为正如对有德者，虽然他们有身体，我们说：\Quote{他是全灵魂，全精神，} 对相反的人也同样如此。保禄也对那些没有完成肉性之工的人说：\BibleQuote{你们不属于肉性。}\BibleRef{罗 8:9} 因此那些生活在享乐中的人既不在灵魂内，也不在精神内。\par

\Emph{伦理}. \Quote{那生活在享乐中的寡妇活着也是死的。} 你们这些妇女，听啊！你们在狂欢和放纵中消磨时光，忽视穷人因饥饿而消瘦和死亡，同时你们自己因不断的奢侈而毁坏自己。这样，你们成了两种死亡的原因：那些死于匮乏的人，以及你们自己，都是由于无度。但如果你们用自己的丰裕弥补他们的匮乏，你们就能拯救两条性命。为什么你们这样用过多来填满自己的身体，却用缺乏来消耗穷人的身体？为什么过量地喂养这一个，而过分地限制那一个？想想食物变成什么，变成了什么。你们难道不厌恶提及它吗？为什么还要热衷于此种积聚？奢侈的增加不过是粪堆的增多！因为自然有其限度，超出这个限度的不是滋养，而是伤害，是粪便的增加。滋养身体，但不要摧毁它。食物被称为滋养，表明它的目的不是伤害身体，而是滋养身体。或许正因如此，食物变成了粪便，使我们不致成为享乐者。因为如果不是这样，如果它不无用且有害于身体，我们就不会停止互相吞噬。如果肚子能随意接收、消化并输送到身体，我们将看到无数的战争和战斗。即使现在，当我们的食物一部分变成粪便，一部分变成血液，一部分变成无用而有害的痰，我们仍然如此沉迷于奢侈，以至于一顿饭可能花掉整个庄园。如果奢侈不是这个结局，我们还会做出什么？我们生活得越奢侈，我们体内充满的气味就越难闻。身体就像一个膨胀的瓶子，到处漏气。打嗝让旁边的人头痛。由于内部的发酵热，蒸汽像熔炉一样散发出来，如果旁观者感到痛苦，那么你认为大脑内部不断受到这些气味的攻击会如何？更不用说加热和阻塞的血管、肝脏和脾脏这些储库，以及排泄粪便的通道了。我们注意保持街道上的排水沟畅通无阻。我们用杆子和拖把清理下水道，以免堵塞或溢出，但我们却不保持身体通道的畅通，反而堵塞和扼杀它们。当污垢上升到国王的宝座，即大脑时，我们却不在意，没有把它当作一位可敬的君王，而是当作一头不洁的野兽。天主特地将那些不洁的成员放在远处，免得我们受到冒犯。但我们不允许这样，反而因过度而毁坏一切。还有其他邪恶可以提及。堵塞下水道会引起瘟疫；但如果外部的恶臭是瘟疫，那么体内封闭且无法排出的恶臭，会对身体和灵魂造成怎样的紊乱？有些人奇怪地抱怨，不明白为什么天主命定我们要带着一堆粪便。但他们自己却增加了这个负担。天主这样设计是为了使我们脱离奢侈，说服我们不要执着于世俗之物。但你们却不愿因此停止贪食，尽管满足只到喉咙，只持续吃饭的时刻，甚至不那么久，快感就消失了，你们仍然沉溺于放纵。难道不是这样吗？一旦它经过口腔和喉咙，快感就停止了？因为它的感觉在于味觉，而味觉满足之后，接着就是恶心，胃无法消化食物，或者很难消化。因此，说得对：\Quote{那生活在享乐中的寡妇活着也是死的。} 因为享乐的灵魂无法听到或看到任何东西。它变得软弱、卑劣、缺乏男子气、狭隘、懦弱，充满厚颜无耻、奴性、无知、愤怒、暴力以及各种邪恶，缺乏相反的美德。因此他说：\par

第7节. \BibleQuote{你要嘱咐这些事，使她们无可指责。}\BibleRef{弟前 5:7}\par

他没有将其留给他们选择。他说，命令她们不要奢侈，假设这公认为一种邪恶，因为他认为奢侈者参与神圣奥迹是不合法且不允许的。\BibleQuote{你要嘱咐这些事，}他说，\BibleQuote{使她们无可指责。} 因此你看，这被视为有罪。因为如果是可选择的，即使不做，我们仍可能是无可指责的。所以，听从保禄，让我们命令奢侈的寡妇不能在寡妇名单中有一席之地。因为如果一个士兵频繁出入浴池、剧场和繁忙的生活场景，就被判断为擅离职守，更何况寡妇呢？那么，让我们不要在此寻求安息，好让我们将来能得到安息。让我们不要在此地生活享乐，好让我们将来能享受真正的快乐、真正的喜悦，它不带来任何邪恶，而是无限的美善。愿天主恩赐我们都参与其中，在耶稣基督内，偕同祂，等等。\par

\SourceRecord{Translated by Philip Schaff.；Nicene and Post-Nicene Fathers, First Series；Vol. 13.；Edited by Philip Schaff.；Buffalo, NY: Christian Literature Publishing Co.,；1889.；<http://www.newadvent.org/fathers/230613.htm>.}

% structure: p0001 source=h1 target=section reason=page-title

\section{弟茂德前书第十四篇讲道}

\Emph{\BibleRef{弟前 5:8}}\par

\Emph{\BibleQuote{但若人不照顾自己的亲属，尤其自己的家人，便是背弃信德，比不信的人更坏。}}\par

许多人认为自己的德行足以使他们获得救恩，如果妥善管理自己的生活，就不需要更多的东西来拯救他们。但他们在这一点上大大错误，这由埋藏了他唯一塔冷通的人的例子所证明，因为他带回来的不是减少的，而是完整的，正如交付给他时一样。这也由真福保禄在此所说的显示出来：\BibleQuote{若有人不照顾自己的亲属。}他所说的照顾是普遍的，涉及灵魂和身体，因为两者都当照顾。\par

\BibleQuote{若有人不照顾自己的亲属，尤其自己的家人，}即那些与他有近亲关系的人，\BibleQuote{他比不信的人更坏。}众先知之首依撒意亚也这样说：\BibleQuote{你不可忽略你自己的骨肉。} \BibleRef{依 58:7} 因为若有人遗弃那些因亲属和姻亲关系而结合的人，他怎能对他人有爱心呢？这岂不显得是虚荣，当他施恩于他人时，却轻视自己的亲属，不照顾他们？如果教导他人，却忽略自己的人，尽管他有更大的能力，也有更高的责任去关爱他们，那又将如何？岂不会有人说：这些基督徒真是有爱心，却忽视自己的亲属？\BibleQuote{他比不信的人更坏。} 为什么？因为后者若不施恩于外人，也不至于忽略自己的近亲。意思是：天主和自然的法律被那不顾自己家人的人违犯了。但如果那不顾他们的人已背弃了信德，比不信的人更坏，那伤害亲戚的人又该被列在哪里？他将被置于何处？但他如何背弃了信德？正如经上所说：\BibleQuote{他们说他们认识天主，但行为上却否认祂。} \BibleRef{铎 1:16} 他们所信的天主曾命令什么？\BibleQuote{不要遮掩自己骨肉。} \BibleRef{依 58:7} 那这样否认天主的人又怎能相信？让那些吝惜财富、忽视亲属的人思考这一点。天主的安排，透过亲属关系将我们连结，给予我们许多彼此行善的机会。因此，当你忽略了不信者都履行的责任时，你岂不背弃了信德？因为信德不仅仅是宣认信仰，而是要做与信德相称的行为。而且在每一件具体事上都有可能信或不信。因为他既已提到奢侈和放纵，就说这样的妇人受罚不仅因为她奢侈，还因为她的奢侈迫使她忽略自己的家庭。他说得有理：因为那为肚腹生活的人，也因此而丧亡，正如\BibleQuote{背弃了信德}一样。但她怎么比不信者更坏？因为忽略亲人与忽略陌生人不同。怎能相同呢？但这里的过错更大：遗弃一个认识的人比遗弃一个不认识的人更严重，遗弃朋友比遗弃非朋友更严重。\par

第9、10节：\BibleQuote{不要登记六十岁以下的寡妇，她应当只做过一个丈夫的妻子，并有行善的声望。}\par

他曾说：\BibleQuote{她们首先应当学习孝敬自己的家人，报答自己的父母。}他也曾说：\BibleQuote{那好宴乐的寡妇，虽活着也是死的。}他曾说：\BibleQuote{若她不照顾自己的亲属，就比不信的人更坏。}在提到哪些缺乏会使一个妇人不配列入寡妇名单后，现在他提到她还应具备什么。那么怎样？我们要因她的年纪而接纳她吗？那有什么功德？她六十岁并不是她自己所能决定的。因此他并不是仅仅谈论她的年龄，而是说，即使她已达到那些岁数，但若没有善行，仍不可列在数目中。那么他为何特别注重年龄？他后来给出了一个原因，并非出于他自己，而是出于寡妇本身。同时让我们听听接下来怎么说。\BibleQuote{并有行善的声望：如果她教养过儿女。}的确，教养儿女并非无足轻重的工作；但教养不仅是照顾他们，还必须教养得好；正如他先前所说：\BibleQuote{若她们持守信德、爱德和圣德，} \BibleRef{弟前 2:15} 请观察他何等经常将善待亲属置于善待陌生人之前。首先他说：\BibleQuote{如果她教养过儿女，}然后，\BibleQuote{如果她款待过旅客，给圣徒洗过脚，救济过患难的人，尽力行过各种善工。}但若她贫穷呢？即使如此，她也不被禁止教养儿女、款待旅客、救济患难者。她并不比那捐了两文钱的寡妇更穷。她虽贫穷，却有一个家，并不露宿街头。他说：\BibleQuote{如果她给圣徒洗过脚，}这不是一件花费很大的工作。\BibleQuote{如果她尽力行过各种善工。}他在这里给予什么训诫？他劝勉他们贡献身体的服务，因为妇女特别适合这样的服侍：为病人铺床，安抚他们休息。\par

奇怪！他对寡妇要求何等严格，几乎与对主教本人相同。因为他说：\Quote{若她殷勤追随一切善工}。这似乎是说，即使她不能亲自履行，也当分担与协作。当他禁止奢华时，他愿她有远见，善于持家，同时恒心祈祷。亚纳便是如此。他对寡妇要求这样的严格，甚至超过对童贞女的要求，尽管他对童贞女也要求严格，并具卓越德行。因为当他说到\Quote{那合宜的}和\Quote{叫她专心事奉主，没有分心}（\BibleRef{格前 7:35}）时，他实际上概括了所有德行。可见守寡不仅在于不再结婚，还需许多其他条件。但他为何劝阻再婚？难道再婚是受谴责的吗？绝非如此；那是异端。他只是愿她今后专注于属灵之事，将全部关怀转向德行。因为婚姻不是不洁的状态，而是充满忙碌。他提到她们有闲暇，而非因守独身而更纯洁。婚姻确实包含许多世俗事务。你若因要为天主服务而守独身，却并未善用那闲暇去款待旅客和圣人，那对你就无益——若你未善用闲暇，则你守独身并非为善的目的，而是如同谴责婚姻状态。因此，那未真正死于世界的贞女，因拒绝婚姻，似乎视之为被诅咒和不洁的。\par

请注意，此处所说的款待不仅是一种友好的接待，而是以热诚和欣悦、以甘心乐意去行，如同接待基督本人。寡妇应亲自做这些服务，而非交给婢女。因为主说：\Quote{我，你们的主和师傅，洗了你们的脚，你们也该彼此洗脚}（\BibleRef{若 13:14}）。即使一个妇人非常富有，出身最高贵，以她的门第和显赫家族为傲，但她与他人的差距，不及天主与门徒之间的差距。你若以基督的身份接待客人，就不要羞耻，反而要光荣。但若不以基督的身份接待，就根本不要接待。他说：\Quote{谁接待你们，就是接待我}（\BibleRef{玛 10:40}）。若不如此接待，便没有赏报。亚巴郎曾接待过路行人，他以为他们是旅客，却未曾将准备的劳役交给仆人，而是亲自承担了大部分服务，并命妻子和面，尽管他家里有三百一十八个生来的仆人，其中必有许多婢女；但他希望自己和妻子得到劳苦的赏报，不仅是花费的赏报。同样，我们也当亲自操劳款待，好使我们成圣，双手蒙福。你若施舍穷人，不要不屑于亲手施与，因为那所给的不是给穷人，而是给基督。谁还会如此可怜，不屑于向基督伸出自己的手呢？\par

这就是款待，这才是真正为天主而做的。但若你傲慢地发号施令，虽请他坐上座，那也不是款待，不是为天主而做的。客人需要很多的关照和鼓励，即便如此，他仍不易不感拘束；因为他的处境如此微妙，在蒙受恩惠时，他感到羞愧。我们应以最殷勤的服务消除那种羞愧，并以言语和行为表明，我们并不觉得自己在施恩，而是在领受恩惠，我们受惠多于施惠。如此善意倍增恩惠。因为，那视自己为受损、以为在施恩的人，会毁灭一切功绩；相反，那视自己为受惠的人，则增加赏报。\Quote{因为天主爱乐捐的人}（\BibleRef{格后 9:7}）。因此，你更应感激穷人接受你的善行。因为若没有穷人，你大部分的罪便不得消除。他们是医治你伤口的医生，他们的手为你带来疗愈。医生伸手敷药，其医术之效，不及穷人伸手接受你施舍所带来之治愈。你交出钱财，罪过也随之离去。古时的司祭正是如此，经上说：\Quote{他们以我子民的罪为食物}（\BibleRef{欧 4:8}）。因此，你所得多于所施，受益多于造福。你借给的是天主，不是人。你增加财富，而非减少。若你施舍而财富未减，那才真是减少！\par

\Quote{若她款待旅客，若她洗圣人的脚}。但这些人是谁？是受苦的圣人，而非任何圣人。因为可能有圣人备受众人服侍。不要探访那些身处富足的，而要探访那些身处患难、不为人知或鲜为人知的。他说：\Quote{为这些最小中的一个做的，就是为我做的}（\BibleRef{玛 25:40}）。\par

\Emph{道德训诲}。不要把你的施舍交给那些在教会里负责分发的人。你要亲手施舍，这样你不仅得到施舍的赏报，还得到善尽服务的赏报。亲自用你的双手施舍。亲自把种子撒入犁沟。在这里，不需要扶犁、套牛、等待季节、翻土或与严寒抗争。这里没有这样的劳苦，因为你播种在天上，那里没有霜冻、冬天或任何这类东西。你播撒在灵魂中，没有人会夺走所播下的，反而被悉心谨慎地保存。亲自撒种吧，为何要剥夺自己的赏报呢？即使管理别人之物，也有极大的赏报。不仅给予有赏报，妥善管理所给予的也有赏报。你为何不愿得到这赏报？因为对此有赏报，听我们读到宗徒们设立\Person{斯德望}负责寡妇的侍奉。\BibleRef{宗 6:5}\par

你要成为自己礼物的分施者。你的仁爱与对天主的敬畏指定你担任那职分。这样，虚荣心就被排除。这滋养灵魂，这圣化双手，这压伏骄傲。这教导你智慧，这激发你的热忱，这使你获得祝福。你离去时，你的头承受着寡妇们所有的祝福。\par

要更恳切地祈祷。用心寻找圣洁的人——真正如此的人，他们隐居在旷野中，不能乞讨，却完全献身于天主。长途跋涉去拜访他们，并亲手施舍。因为若你亲自施舍，你本人可能大有益处。你看见他们的帐篷、他们的住处吗？你看见旷野吗？你看见孤寂吗？当你去施舍金钱时，往往也献出了你的整个灵魂。你被留住，成了他的同囚，同样远离了世界。\par

即使见到穷人也有很大的益处。他说：\Quote{往居丧的家去，比往宴会的家去更好。}\BibleRef{训 7:2}因为后者使灵魂燃烧。你若能模仿奢华，就会被鼓励放纵；若不能，就会忧伤。在居丧的家里，没有这类事。你若无能力奢华，不会痛苦；若有能力，也会受节制。隐修院确实是居丧的家。那里有麻衣和灰烬，有孤寂，没有笑声，没有世俗事务的压力。有斋戒和躺卧地上；没有美食不洁的气味，没有流血，没有喧嚣、纷扰或拥挤。那是一处宁静的港湾。它们如同明灯从高处照耀远方的航海者。它们停泊在港口，吸引众人到自己的平静中，拯救那些注视它们的人免于覆舟，不让仰望它们的人在黑暗中行走。你去与他们结交，与他们为友，拥抱他们神圣的脚——触摸它们比触摸别人的头更光荣。若有人抱住雕像的脚，因为雕像带有君王的模样，难道你不愿抱住那怀有基督的人的脚，从而得救吗？圣人的脚是圣洁的，即使他们是穷人，而不敬虔之人的头也不可尊荣。圣人的脚有如此功效，以致当他们从脚下抖落尘土时，就施加惩罚。当一位圣人在我们中间时，我们不要为他所拥有的任何东西感到羞耻。所有圣人都是在正确的信仰下活出圣善生活的人，即使他们不行奇迹也不驱魔，仍然是圣人。\par

那么，去他们的帐幕吧。去圣人的隐修院，就如同从地上升到了天上。你在那里看不到在私人住宅里所见的东西。那团体免于一切不洁。那里有静默和深深的安宁。\Quote{我的}和\Quote{你的}这些话在他们中间并不使用。你若在那里待上一整天甚至两天，就会感到更大的快乐。在那里，天一亮，或者更确切地说，天未亮时，雄鸡啼叫，你不会像在房屋里看到的那样——仆人打鼾，门户紧闭，所有人都像死人般沉睡，而外面的骡夫摇着铃铛。那里没有这一切。所有人立即摆脱睡眠，当他们的院长召唤时，恭敬地起身，组成神圣的歌咏团，站立，同时举手歌唱圣诗。因为他们不像我们，需要许多小时才能从沉重的头脑中驱散睡意。我们确实一醒来，就坐一会儿伸伸四肢，去上厕所，然后洗脸洗手；之后穿鞋穿衣，花费大量时间。\par

那里并非如此。没有人呼唤仆人，因为每个人都自己服侍自己；也不需要许多衣服，更不需要驱散睡意。因为一睁开眼睛，他就如同已经清醒良久一般泰然自若。因为当内心没有被过多的食物所窒息时，它很快会恢复，立刻清醒。双手总是洁净的；因为他的睡眠有规律而安稳。在他们中间，没有人打鼾、喘粗气、在床上辗转反侧或身体暴露；他们睡眠中的姿态如同清醒时一样端庄，这一切都是他们灵魂有秩序的明证。这些人真是圣人，是人群中的天使。当你听到这些时，不必惊奇。因为他们对天主的极大敬畏，不容他们陷入沉睡深处、淹没理智，而是睡眠轻轻临到他们，仅仅使他们得到休息。正如他们的睡眠，他们的梦境也是如此，没有疯狂的幻想和怪异的景象。\par

但是，如我所说，在鸡鸣时，他们的院长来了，轻轻用脚触碰睡着的人，唤醒他们所有人。因为没有人赤裸而睡。然后他们一站起来，就站立着，用极大的和谐和优美的曲调歌唱先知式的诗歌。无论是竖琴、笛子还是其他乐器，都不能发出如此甜美的旋律，就像你听到这些圣人在他们深邃宁静的隐修院所唱的。这些诗歌本身也合宜，充满对天主的爱。他们说：\BibleQuote{在夜间，请向天主举起你们的手。我的灵魂在夜间切望你，我以我内在的精神清晨寻求你，}\BibleRef{依 26:9}还有\BibleBook{达味}{圣咏集}，使人泪如泉涌。因为当他唱道：\BibleQuote{我因呻吟而疲惫，每夜使我床榻漂浮，以泪水浸透我的床铺}\BibleRef{咏 6:6}：以及又：\BibleQuote{我吃灰烬如同食粮。}\BibleRef{咏 101:9}\BibleQuote{人算什么，你竟怀念他？}\BibleRef{咏 8:4}\BibleQuote{人好似虚幻，他的年日如影消逝。}\BibleRef{咏 143:4}\BibleQuote{不要害怕别人富裕，他家中的荣耀加增}\BibleRef{咏 48:16}；以及：\BibleQuote{使众人同心合意住在家里}\BibleRef{咏 67:6}：以及：\BibleQuote{我一日七次赞美你，因你正义的判断}\BibleRef{咏 118:164}：以及：\BibleQuote{半夜我起来感谢你，因你正义的判断}\BibleRef{咏 118:62}：以及：\BibleQuote{天主必拯救我的灵魂脱离阴间的权势}\BibleRef{咏 48:15}：以及：\BibleQuote{我虽走过死荫的幽谷，不惧怕灾祸，因你与我同在}\BibleRef{咏 22:4}：以及：\BibleQuote{我不怕黑夜的惊骇，白日飞的箭，黑暗中行的瘟疫，正午毁灭人的毒病}\BibleRef{咏 91:5}：以及：\BibleQuote{我们被看作待宰的群羊}\BibleRef{咏 43:22}：他表达了对天主炽热的爱。又当他们与天使一同歌唱时（因为天使那时也在歌唱）：\BibleQuote{请从天上赞美主。}\BibleRef{咏 148:1}而与此同时，我们却在打鼾，或搔头，或仰卧沉思无尽的欺骗。想想看，他们整夜从事这样的工作，这意味着什么。\par

当白天来临时，他们又休息；因为当我们开始工作时，他们有一段时间休息。但我们每个人，在白天，呼唤邻居，计算开支，然后去广场；颤抖着出现在法官面前，害怕算账。另一个人去戏院，另一个人去做自己的事。但这些圣人，在完成了晨祷和诗歌后，继续诵读圣经。有些人学会了抄写书籍，每人分配有自己的房间，在那里永享安静；没有人闲荡，没有一个人说话。然后在第三、第六、第九时辰和傍晚，他们进行祈祷，将一天分为四部分，每个部分结束时，他们用圣咏和诗歌赞美天主，而别人在吃饭、大笑、嬉戏、暴饮暴食时，他们专注于诗歌。因为他们没有时间用于餐桌或这些感官之事。饭后他们继续同样的行程，此前给自己一段时间睡眠。世俗的人在白天睡觉，但这些人在夜间警醒。他们真是光明之子！而当那些人睡了大半天，带着沉重出去时，这些人仍然收敛心神，直到傍晚都不进食，专注于诗歌。其他人，当傍晚来临时，急忙去洗澡和各种娱乐，但这些人在劳作之后，就去他们的餐桌，不召集许多仆人，也不使房屋喧闹混乱，也不摆上美味佳肴和香气蒸腾的食物，而是有人只吃面包和盐，其他人还加橄榄油，体弱者还有蔬菜和豆类。然后坐了片刻，或者不如说以诗歌结束一切后，他们各自上床休息，床只为休息而非奢侈。没有对官吏的恐惧，没有主人的傲慢，没有对奴隶的惊吓，没有妇女儿童的吵闹，没有成堆的箱子或多余的衣服堆放，没有金银，没有卫兵和哨兵，没有仓库。这一切都没有，但那里充满了祈祷、诗歌和属灵的芬芳。没有任何属肉体的东西。他们不惧怕盗贼的攻击，因为没有什么可被剥夺，没有财富，只有灵魂和身体，如果这些被夺去，那不是损失而是收益。因为经上说：\BibleQuote{对我来说，活着就是基督，死了就是获益。}\BibleRef{斐 1:21}他们已经摆脱了一切束缚。确实，\BibleQuote{喜乐的声音在义人的帐幕中。}\BibleRef{咏 117:15}\par

在那里听不到哀号和悲泣。他们的房屋没有那种忧郁和哭喊。虽然他们的身体并非不死，死亡确实在那里发生，但他们不知道死亡，如同死亡一般。去世的人被唱着赞美诗送到坟墓。他们称此为送殡，而非埋葬。当有人去世时，他们便充满喜乐，大大欢喜；或者说，没有人能忍心说某人死了，而是说他圆满了。于是便有感谢、极大的光荣和喜乐，每个人都祈求自己的结局如此，以便自己的战斗结束，从劳苦和挣扎中安息，并见到基督。若有病人，他们不哭泣哀叹，而是诉诸祈祷。常常不是医生的照料，而是信德独自使病人痊愈。如果需要医生，那时也有最大的坚定和泰然。没有妻子扯自己的头发，没有孩子提前哀悼自己成为孤儿，没有奴隶恳求临终者保证他们会被托付给好人家。从这一切中解脱出来，灵魂在最后的一息只着眼于一件事：得以在天主面前蒙受恩宠而离世。若有疾病发生，其原因与其说是可耻，不如说是光荣，不同于其他情况。因为疾病并非来自贪食或头脑胀满，而是来自警醒和斋戒等；因此它很容易消除，因为只需减轻这些操练的强度即可。\par

那么请告诉我，你会说，是否有人能在教会中为圣人洗脚？这样的人在我们中间能找到吗？是的：无疑他们是这样的。然而，当我们描述了这些圣人的生活时，不要轻视那些在教会中的人。我们中间常常有许多这样的人，尽管他们隐藏着。也不要因为他们走家串户、进入市场或公开露面而轻视他们。天主甚至命令了这样的服侍，说：\BibleQuote{审判孤儿，为寡妇辩护。}\BibleRef{依 1:17}行善的方法有许多，就像宝石有许多种类，尽管都称为宝石：一种四面明亮圆润，另一种有别的美丽。这如何可能？正如珊瑚，通过某种技艺，使线条延伸，棱角修整，另一种颜色比白色更悦目，绿宝石胜过一切绿色，另一种有丰富的血红色，另一种有胜过海洋的天蓝色，另一种比紫色更灿烂，如此在多样性上媲美一切花朵或太阳的颜色。然而，它们都称为宝石。圣人也是如此。有的管束自己，有的管束教会。因此，保禄说得好：\BibleQuote{若她洗过圣人的脚，若她周济过受苦的人。}因为他这样说，为的是激发我们众人效法。那么，让我们赶紧去行这样的事，好使我们将来能夸耀说，我们洗过圣人的脚。因为如果我们该洗他们的脚，我们更该亲手把钱给他们，同时力求隐藏自己。他说：\BibleQuote{你右手所做的，不要让左手知道。}\BibleRef{玛 6:3}\par

你为何找这么多见证人？不要让你的仆人知道，如果可能，连你的妻子也不要。那欺骗者有许多阻碍。那从不干涉的妻子，往往会因虚荣或其他动机而阻碍这些工作。就连那有贤妻的亚巴郎，当他准备献上自己儿子时，也对她隐瞒了这事，尽管他不知道将要发生什么，却深信自己必须宰杀儿子。那么，一个普通人会说什么呢？岂不是\Quote{是谁做出这样的事？}岂不指责他残忍和野蛮？他的妻子甚至不被允许见她的儿子，听他的遗言，目睹他垂死的挣扎。但他像俘虏一样把他带走了。那位义人，虽不知何事，却被热忱所陶醉，只想着如何完成所命令的。没有仆人，没有妻子在场，他自己也不知道结果如何。但他一心要献上纯洁的祭品，不愿以眼泪或任何反对来玷污它。请注意依撒格是多么温良地问：\BibleQuote{请看，火和柴都有了，但燔祭的羔羊在哪里呢？我的儿子，天主自会预备燔祭的羔羊。}\BibleRef{创 22:7}他这样说，预言了天主将在祂的圣子内预备燔祭，这话当时也应验了。但他为何向那要被献祭的人隐瞒呢？因为他怕他惊慌，怕他不配。他在整件事中是如此谨慎精明！因此，圣经说得好：\BibleQuote{你右手所做的，不要让左手知道。}如果我们有一个人如自己的肢体一样亲爱，就不要急于向他显示我们的施舍，除非必要。因为许多恶事可能由此而生：人会被引向虚荣，而且常常产生阻碍。为此，我们要隐藏这事，如果可能，连我们自己也要隐藏，好能获得所应许的福分，赖我们的主耶稣基督的恩宠和慈爱。\par

\SourceRecord{Translated by Philip Schaff.；Nicene and Post-Nicene Fathers, First Series；Vol. 13.；Edited by Philip Schaff.；Buffalo, NY: Christian Literature Publishing Co.,；1889.；<http://www.newadvent.org/fathers/230614.htm>.}

% structure: p0001 source=h1 target=section reason=page-title

\section{论弟茂德前书讲道词第十五篇}

\BibleRef{弟前 5:11-15}\par

\BibleQuote{至于年轻的寡妇，你要拒绝录用，因为她们当不起基督的约束，就想再嫁；她们被判罪，是因为抛弃了起初的信德。同时她们又学会懒惰，挨家闲游；不但懒惰，而且饶舌，好管闲事，说些不当说的话。所以我愿意年轻的寡妇再嫁，生儿育女，料理家务，不给敌人以诽谤的把柄。因为有些人已经随从了撒殚。}\par

保禄谈论了许多关于寡妇的事，规定了她们被接纳的年龄，说：\BibleQuote{不要选录六十岁以下的寡妇}，又描述了寡妇应有的资格：\BibleQuote{如果她养育过子女，如果她款待过旅客，如果她洗过圣徒的脚}，然后接着说：\BibleQuote{至于年轻的寡妇，你要拒绝录用}。但关于童贞女，虽然她们跌倒的情况更为严重，他却未提及这类事，这很合理。因为她们是怀着更高的志向注册的，她们的行为源于更高超的心神。因此，他用\BibleQuote{专心侍主，没有挂虑} \BibleRef{格前 7:34} 和\BibleQuote{没有出嫁的，挂虑主的事} \BibleRef{格前 7:34} 来代替款待旅客和洗圣徒的脚。而且他未为她们限定年龄，很可能是因为这件事已由他对寡妇所作的论述解决了。但正如我所说，童贞的选择源于更高的目的。此外，在此情况下已有跌倒发生，所以她们给了他制定规则的口实，而童贞女中未发生类似的事。因为显然有人已经堕落了，正如他所说：\BibleQuote{当她们当不起基督的约束，就想再嫁}；又说：\BibleQuote{因为有些人已经随从了撒殚。}\par

\BibleQuote{至于年轻的寡妇，你要拒绝录用，因为她们当不起基督的约束，就想再嫁}；也就是说，当她们变得傲慢和放纵时。就如我们可能对一个义人说：\Quote{让她走吧，因为她已成为别人的了。} 因此，他表明她们选择寡居并非出于判断。所以，一个寡妇因寡居的状态已许配给基督。因为祂说过：\BibleQuote{我是寡妇的保护人，孤儿的父亲} \BibleRef{咏 67:5}。他表明她们没有按应有的方式选择寡居，而是变得狂妄；不过他还是容忍她们。的确，他在别处说：\BibleQuote{我曾把你们许配给一个丈夫，把你们如同贞洁的童女献于基督} \BibleRef{格后 11:2}。在她们将自己的名字献给祂之后，\BibleQuote{她们想再嫁}，他说，\BibleQuote{她们被判罪，是因为抛弃了起初的信德。} 这里的信德指对盟约的忠信。就好像他说，她们对基督失信，羞辱了祂，违背了祂的盟约。\BibleQuote{同时她们又学会懒惰。}\par

因此，他不仅命令男人，也命令女人工作。因为懒惰是所有罪恶的老师。而且她们不仅遭受这种谴责，也陷入其他罪恶。如果对于一个已婚妇女来说，挨家闲游是不合适的，那么对于童贞女来说就更不合适了。不但懒惰，而且饶舌，好管闲事，说些不当说的话。\BibleQuote{所以我愿意年轻的寡妇再嫁，生儿育女，料理家务。}\par

那么，当对丈夫的牵挂被解除，而取悦天主的牵挂又不约束她们时，她们自然会成为懒惰者、饶舌者和好管闲事者。因为不管自己事务的人，就会干涉别人的事；而专心自己事务的人，则不会理会也不关心别人的事。没有比一个女子好管闲事更不合宜的了，对一个男子来说也同样不合宜。这是厚颜和放肆的极大标志。\par

\BibleQuote{所以我愿意}，他说（因为她们自己也愿意如此），\BibleQuote{年轻的寡妇再嫁，生儿育女，料理家务。}\par

这条道路至少比另一条好。她们本当关心天主的事，本当保存自己的信德。但既然她们不这样做，避开更坏的道路就好。天主并未因她们再婚而受辱，而且她们不会陷入那些已被指责的行为。这样的寡居不会带来任何益处，而这种婚姻却能带来好处。因此，妇女们能够纠正那种懒散和虚荣的心。\par

但是，既然有人已经堕落了，为什么他不说要对她们多加留意，以免她们陷入他提到的错误？为什么他命令她们结婚？因为婚姻并非被禁止，而且对她们是一种保护。为此他补充说，她们\BibleQuote{不给敌人以诽谤的把柄}，或口实，\BibleQuote{因为有些人已经随从了撒殚。} 因此，他主张拒绝这样的寡妇，意思并非不应有年轻的寡妇，而是不应有奸妇，不应有懒惰、好管闲事、说不当说的话的人，并且不应给敌人以可乘之机。若未发生这类事，他本来不会禁止她们。\par

\BibleRef{弟前 5:16} \BibleQuote{若有信男信女家中有寡妇，应照顾她们，不可累及教会，好使教会能照顾那些真正的寡妇。}\par

请留意，他再次称那些孤苦无依、从别处得不到帮助的人为\Quote{真正寡妇}。这样安排是好的。因为这样便达成了两个重要的目标：一方面给予行善的机会，另一方面使这些寡妇得到体面的供养，而教会也不受负担。他说\Quote{若有任何信者}这句话很恰当。因为信主的妇女不应由不信者供养，免得她们显得需要他们。注意他如何有说服力地说；他没有说\Quote{让她们奢侈地供养}，而是说\Quote{让她们救济她们}。他说：\Quote{为使教会能救济那些真正寡妇}。因此，她也获得这帮助的赏报，因为她帮助教会，不仅帮助教会，也帮助那些寡妇，教会因而得以更慷慨地供养她们。\Quote{所以我愿意年轻的寡妇}——做什么？过奢侈享乐的生活？绝不；而是——\Quote{结婚、生儿育女、治理家务}。为了避免让人以为他鼓励她们过奢侈生活，他补充说，免得给仇敌机会毁谤。她们本应超脱世俗之事，但既然不能，至少让她们正直地处在其中。\par

第17～18节：\Quote{那些善于管理的长老，应受加倍的敬重，尤其是那些在言语和道理上劳苦的。因为经上说：\BibleQuote{牛在场上踹谷的时候，不可笼住它的嘴。}又说：\BibleQuote{工人得工价是应当的。}}\par

他所说的\Quote{敬重}是指关注他们、供应他们所需，正如他接着引用的话所表明的：\BibleQuote{不可笼住牛在场上踹谷的嘴}\BibleRef{申 25:4}；以及\BibleQuote{工人得工价是应当的}\BibleRef{路 10:7}。因此，当他说\Quote{敬重寡妇}时，意思是供应她们一切所需。同样，他说\Quote{使教会能救济那些真正寡妇}；又说\Quote{敬重那些真正是寡妇的}，即那些贫穷的，因为越贫穷就越算真正寡妇。他引述法律，又引述基督的话，两者在此一致。因为法律说：\BibleQuote{牛在场上踹谷的时候，不可笼住它的嘴。}请看，他多么希望教师劳苦！因为实在没有其他劳苦能比得上教师的劳苦。这是来自法律的。但他又如何引用基督的话呢？\BibleQuote{工人得工价是应当的。}因此，我们不要只看赏报，也要看诫命的条款。他说：\BibleQuote{工人得工价是应当的。}所以如果有人懒惰、奢侈，他就不配。除非他像踹谷的牛一样，负轭忍受炎热和荆棘，直到将谷收入仓内，他才配得。因此，应毫不吝啬地供应教师所需，使他们不致沮丧灰心，也不致因分心于次要之事而失去更重要的事；使他们能专注于属灵之事，不顾念世俗之事。肋未人就是这样：他们没有世俗的挂虑，因为平信徒负责供应他们，法律也规定了他们的收入，如什一税、金祭、初熟之果、誓愿等等。法律恰当地将这些指定给他们，作为对现世之事的寻求。但我只说一句话：那些管理的人应当有食物和衣服，使他们不致因这些事分心。但\Quote{加倍的敬重}是什么意思？是对寡妇的双倍？或对执事的双倍？或只是丰厚的供应？我们不要只想到加倍的供养，也要想到\Quote{善于管理}这个附加条件。什么是善于管理？让我们听\Person{基督}的话，他说：\BibleQuote{善牧为羊舍命}\BibleRef{若 10:11}。所以，善于管理就是因关心羊群而绝不吝惜自己。\par

\Quote{尤其是那些在言语和道理上劳苦的。}那么，那些说言语和道理不需要的人在哪里呢？而他对\Person{弟茂德}说：\BibleQuote{你要思念这些事，全神贯注于这些事}；又说：\BibleQuote{你要专心读经、劝勉和教导；因为这样做，你不但救自己，也救听你的人}\BibleRef{弟前 4:15}。这些人就是他最希望被尊敬的人，他补充了理由，因为他们肩负巨大的劳苦。因为如果一个人既不清醒，也不勤奋，只是轻松无忧地坐在他的座位上，而另一个人因焦虑和努力而耗尽自己，尤其是如果他不熟悉世俗文学，那么后者不是更应该被尊敬吗？后者比其他人更投入这样的劳苦！因为他面对无数人的议论：有人指责他，有人称赞他，有人嘲笑他，有人批评他的记忆和文笔，这需要很大的力量才能忍受一切。这一点很重要，并且对教会的建立大有贡献：教会的管理者应善于教导。如果缺乏这一点，教会中许多事情就会败坏。因此，除了好客、温和、无可指责的生活之外，他还列举了这一点，说：\BibleQuote{善于教导}。不然为什么他被称为教师呢？有人说，他可以通过生活的榜样教导哲学，因此其他一切都是多余的，不需要言语教导就能成长。但为什么\Person{保禄}说\Quote{尤其是那些在言语和道理上劳苦的}呢？因为当涉及道理时，怎样的生活能起作用呢？他所说的言语是什么？不是浮夸的言词，也不是用外在装饰点缀的演讲，而是拥有圣神大能、充满智慧和悟性的言语。它不需要华丽的辞藻，而是需要思想来表达；不需要写作技巧，而是需要心智的力量。\par

第19节：\BibleQuote{控告长老的事，除非有两三个见证，不要受理。}\BibleRef{弟前 5:19}\par

那么，我们是否可以接受对一位年轻人的控告，或对任何人的控告而没有证人呢？难道我们不是应该在所有情况下都以最大的精确性来进行判断吗？那么他是什么意思呢？他意思是：对任何人都不应如此，特别是对年长者。因为他所说的年长者不是指职务，而是指年龄，因为年轻人比年长者更容易陷入罪恶。由此明显可见，教会，甚至整个\Place{亚细亚}的民众，现在都已托付给了\Person{弟茂德}，这就是为什么他和他谈论有关年长者之事的原因。\par

第20节。\BibleQuote{凡犯罪者，当在众人前受斥责，使其余人也惧怕。}\BibleRef{弟前 5:20}\par

不要仓促地革除他们，而要仔细查问所有情况，等你完全了解之后，再严厉地追究犯罪者，使他人引以为戒。因为草率地定罪是错误的，不惩罚明显的罪行就是为他人开辟道路，并使他们胆敢犯罪。\par

\Quote{斥责}，他说，是为了表明这不是轻率地做，而是严厉地做。因为这样他人就会有所戒惧。那么基督怎么说呢？他说：\BibleQuote{如果你的弟兄得罪了你，去，趁着他和你单独在一起时，指出他的错来。}\BibleRef{玛 18:15} 但基督自己也允许他在教会中受责备。那么怎样呢？难道在众人面前受斥责不是更大的丑闻吗？怎么会呢？因为知道罪行而不知道惩罚，这才是更大的丑闻。因为当罪人不受惩罚时，许多人就会犯罪；当他们受惩罚时，许多人就会变好。天主自己也是这样做的。祂使法老出来，并公开惩罚他。还有\Person{拿步高}，以及许多其他的人，无论是城市还是个人，我们都看到他们受到了惩罚。因此，\Person{保禄}希望所有人都敬畏他们的主教，并把他置于众人之上。\par

因为许多判断是基于猜疑形成的，他说，应当有证人，有人按照古法来定罪犯罪者。\BibleQuote{凭两个或三个见证人的口，一切案件都要确定。}\BibleRef{申 19:15} \Quote{不可收控告年长者的状子。} 他不是说\Quote{不要定罪}，而是说\Quote{不可收控告}，完全不要带他来审判。但如果两个见证人是假的呢？这很少发生，在审判的查问中可以发现。因为罪行是暗中犯下的，我们应当满意于两个证人，这对于查证是足够的。\par

但若罪行是昭彰的，却没有证人，只有强烈的怀疑呢？上面已经说过，他应该\BibleQuote{在教外人士中也必须有良好的声誉。}\BibleRef{弟前 3:7}\par

让我们因此怀着敬畏爱天主。法律原不是为义人设立的；但因为大多数人出于强制而非出于选择而行善，所以恐惧的原则对于根除他们的欲望大有裨益。因此，让我们听从地狱之火的威胁，好使我们因这有益的恐惧而受益。因为如果天主，有意将罪人投入其中，却没有事前用这火警告他们，许多人就会陷进去。因为，如果在这样的恐惧震撼我们灵魂时，有些人仍轻易犯罪，仿佛没有这回事一样；那么，如果这火没有被宣布和警告，我们岂不会犯下更大的恶行吗？所以，正如我常说的，地狱的威胁显示天主对我们的眷顾，不亚于天堂的应许。因为威胁与应许协力，借由恐惧催逼人进入天国。让我们不要以为这是残忍的事，而是怜悯和仁慈；是天主对我们的关怀和爱。如果在\Person{约纳}的日子，\Place{尼尼微}的毁灭没有被警告，那毁灭就不会被避免。若不是那威胁，\Place{尼尼微}不会幸免；那威胁是：\BibleQuote{尼尼微将被倾覆。}\BibleRef{纳 3:4}如果地狱没有被警告，我们都会堕入地狱。如果那火没有被宣告，没有人能逃脱那火。天主宣告祂将做祂不愿做的事，为能做成祂愿做的事。祂不愿罪人死去，因此祂以死亡警告罪人，为不用施加死亡。祂不仅说了话，还亲自显明了那事，好使我们能逃脱。为避免有人认为这仅是恐吓，祂已借由祂在地上所做的事显明了它的真实。难道你们没有在洪水中看到地狱的象征，在毁灭万物的雨水中看到吞噬一切的火的形象吗？\BibleQuote{当\Person{诺厄}的日子，}祂说，\BibleQuote{人们嫁娶}\BibleRef{玛 24:38}，现在也是如此。当时早在事情发生前很久就预言了，现在也预言了四百年甚至更久之前：却无人留意。它被视为寓言，被当作笑柄；无人畏惧它，无人想到它而哭泣或捶胸。火河沸腾，火焰点燃，我们却在欢笑、享乐、无忧无虑地犯罪。没有人记得那一日。没有人想到眼前的事物正在过去，只是暂时的，尽管每天发生的事件都在呐喊，发出可怕的声音。早逝、生活中的变故、自身的软弱和疾病，都不能教训我们。这些变化不仅在我们自身可见，在元素本身也是如此。每天在我们不同的年龄，我们经历某种死亡；在任何情况下，我们所见到的事物之特征就是不稳定。冬天、夏天、春天、秋天，没有一个是持久的；一切都在奔跑、飞逝、流过。我何必提及凋谢的花朵、高位、今日为王而明日不再的君王、富人、宏伟的屋宇、黑夜和白昼、太阳和月亮呢？因为月亮会亏缺，太阳有时会亏蚀，常常被云遮暗。总之，在可见的事物中，有什么是永恒的呢？没有！不，在我们身上，除了灵魂，没有什么是永恒的，而灵魂我们却忽略了。对于易变的事物，我们过分关心，仿佛它们是永久的；但对于那将永远存留的，我们却忽略，仿佛它很快就消逝。有人得以行大事，但它们只持续到明天，然后他就逝去，正如我们看见那些曾有更大权柄、而今已不见踪影的人的例子。生命是一场梦、一场戏；正如在舞台上，当场景转换，各种景象消失；正如阳光升起时，梦境消散；同样，当结局来临时，无论是宇宙的结局还是每个人的结局，一切都在消融和消散。你栽种的树依旧，你建造的房屋也依然立着；但栽种者和建造者却离去、消亡。然而这些事发生了，我们却毫不留意；我们继续生活在奢华享乐中，不断为自己置办这些东西，仿佛我们是不朽的。\par

请听撒罗满所说，他亲身经历过现世：\BibleQuote{我为自己建造了房屋，栽植了葡萄园，开辟了花园和果园，并挖了水池。我又聚集了金银，为自己召来了歌男歌女，以及许多奴仆和牲畜。}\BibleRef{训 2:4} 没有一个人比他生活得更奢华、更荣耀。没有一个人比他更聪明、更有权势，没有人看见万事都如此称心如意。那么结果如何呢？他从这一切中并没有得到享受。最后他自己是怎么说的？\BibleQuote{虚而又虚，万事皆虚。}\BibleRef{训 12:8} 这不是普通的虚，而是极端的虚。让我们相信他，抓住那没有虚、只有真理的东西，那建立在坚固磐石上的，那里没有衰老，没有衰落，一切都在兴盛繁荣，没有朽坏，不会变旧，也不会接近消亡。我恳求你们，让我们以真挚的感情爱天主，不是出于对地狱的恐惧，而是出于对天国的渴望。因为有什么能与看见基督相比呢？确实没有什么！有什么能与享受那些美善相比呢？确实没有什么！难怪没有什么，因为\BibleQuote{天主为爱祂的人所准备的，是眼所未见，耳所未闻，人心所未曾想到的。}\BibleRef{格前 2:9} 让我们渴望得到那些事物，轻视这一切。我们不是一直在抱怨人生是虚无吗？你为什么为虚无的事忧虑？你为什么为虚无的事物如此劳苦？你看见华丽的房屋，难道被它们迷惑了吗？抬头望天。将你的视线从石柱转向那美丽的建筑，与之相比，其他的都如同蚂蚁和蝼蚁的作品。从那景象中学习哲学，上升到天上的事物，从那里观看我们华丽的建筑，发现它们都是虚无，不过是小孩子的玩具。你没有看到吗，你升得越高，空气就越纯净、越轻盈、越清澈？在那里，施舍的人有他们的住所和帐幕。在地上的这些，在复活时都要消解，或者说在复活之前就已经被时间摧毁了。不，常常在他们最繁荣的时候，地震会推翻它们，或者火灾会彻底毁灭它们。因为不仅人的身体，就连他们的建筑也容易遭受过早的死亡。不，有时被时间侵蚀的东西在地震中屹立不倒，而闪耀的、牢固的、新建的建筑物却被闪电击中而毁灭。我相信，这是天主的干预，免得我们以我们的建筑为傲。\par

你还要另一个快乐的理由吗？去公共场所吧，你在那里与别人平等共享。因为最华丽的私人住宅，毕竟不如公共建筑辉煌。在那里，你可以想呆多久就呆多久。它们属于你，与属于别人一样，因为它们是你与别人共有的。它们是公共的，不是私人的。但你说，它们不令你快乐。它们不令你快乐，部分是因为你熟悉它们，部分是因为你的贪婪。所以，愉悦不在于美，而在于占有！所以，快乐在于贪得，在于希望把别人的东西都变成自己的！我们要被这些东西钉住多久？我们要被钉在地上，像虫子一样在泥土中爬行多久？天主给了我们一个由泥土构成的身体，为的是让我们带着它升天，而不是让我们把灵魂也拖到地上。它虽是属土的，但如果我们愿意，它可以成为属天的。看看天主多么看重我们，将如此卓越的身躯托付给我们。祂说：我造了天地，并赐给你创造的能力。使你的地成为天。因为这是在你能力之内的。\BibleQuote{我是那创造并改变万物的。}\BibleRef{亚 5:8} 这是天主论及自己所说的。祂也给了人类相似的能力；如同一位慈父般的画家，将自己的技艺教给儿子。祂说：我塑造了你美丽的身体，但我给了你塑造更美好事物的能力。使你的灵魂美丽。我说：\BibleQuote{地上要生出青草和所有结果实的树木。}\BibleRef{创 1:11} 你也要说：让这地结出它应有的果实，你所愿产生的，就必产生。\BibleQuote{我造夏天和云彩，我造闪电和风。}\BibleRef{亚 4:13} \BibleQuote{我造了龙，即魔鬼，来戏弄它。}\BibleRef{咏 103:26} 我也不吝惜给你同样的能力。你若愿意，可以戏弄它，并像绑麻雀一样绑住它。我使太阳升起，光照恶人也光照善人；你也要效法我，将你所有的分给善人和恶人。当我受嘲笑时，我忍耐，并善待那些嘲笑我的人。你要效法我，尽你所能。我行善，不图回报；你也要效法我，行善不求回报。我在天上点亮了光体。你要点亮比这些更亮的光体，因为你能，借着光照那些在错误中的人。因为认识我比看见太阳有更大的益处。你不能造一个人，但你能使他成为正义和蒙天主悦纳的。我塑造了他的实体，你要美化他的意志。看，我多么爱你，把更大的事的能力交给了你。\par

亲爱的，请看我们是多么受尊荣！然而有些人如此不讲理和忘恩负义，竟说：\Quote{我们为什么被赋予自由意志？} 但如果我们没有自由意志，又怎能在我们所列举的这一切事上效法天主呢？\BibleQuote{祂说：我统治天使，你们也借着那初果来统治。}\BibleRef{格前 15:23} 我坐在王座上，你们也借着那初果与我同坐。正如经上所记：\BibleQuote{祂使我们一同复活，并使我们在基督耶稣内一同坐在天上。}\BibleRef{弗 2:6} 借着那初果，革鲁宾和色辣芬，以及所有天上的万军、宰制者、掌权者、宝座和统御者，都朝拜你们。不要轻视你的身体，因为如此崇高的荣耀属于它，以至于无形体的大能者都因它而战栗。\par

但我要说什么呢？我不仅以此方式向你显示了我的爱，还通过我所受的苦。为了你，我被吐唾沫，被鞭打。我空虚了自己，舍弃了荣耀；我离开了我的圣父，来到你面前——你仇恨我，躲避我，厌恶听到我的名。我追赶你，跑在你后面，为了追上你。我将你联合并连接到我自己，我说：\Quote{食我，饮我。}我在高天怀抱你，在地上拥抱你。我在天上拥有你的\OriginalTerm{初果}{First-fruits}，这还不够吗？这不能满足你的爱情吗？我下降到地上：我不仅与你混合，我还交织在你内。我被咀嚼，碎成细微的颗粒，使散布、混合和结合更加完美。联合之物仍保留在各自的界限内，但我与你交织在一起。我不愿我们之间再有分隔。我愿意我们二人成为一体。\par

因此，知道这些事，并纪念祂对我们丰厚的眷顾，让我们做一切事，以显明我们不辜负祂的伟大恩赐。愿天主恩赐我们众人获得这恩赐，藉着我们的主基督耶稣的恩宠和慈爱，偕同祂，等等。\par

\SourceRecord{Translated by Philip Schaff.；Nicene and Post-Nicene Fathers, First Series；Vol. 13.；Edited by Philip Schaff.；Buffalo, NY: Christian Literature Publishing Co.,；1889.；<http://www.newadvent.org/fathers/230615.htm>.}

% structure: p0001 source=h1 target=section reason=page-title

\section{弟茂德前书第十六篇讲道}

\BibleRef{弟前 5:21-23}\par

\BibleQuote{我在天主和主耶稣基督并蒙拣选的天使面前嘱咐你：要遵守这些事，不可存成见，行事也不可有偏心。不可轻易给人覆手，也不要分担保别人的罪；要保守自己纯洁。不要再只喝水，但要稍微用点酒，为你的胃病和你屡次的软弱。}\par

在谈到主教和执事、男人和女人、寡妇和长老以及所有其他人，并显示了主教的权威何等伟大之后，现在他正讲到审判，他补充说：\BibleQuote{我在天主和主耶稣基督并蒙拣选的天使面前嘱咐你：要遵守这些事，不可存成见，行事也不可有偏心。} 他就是这样令人敬畏地告诫他。因为虽然弟茂德是他亲爱的儿子，但他并不因此畏惧他。正如他自己不羞于说：\BibleQuote{恐怕我传福音给别人，自己反被弃绝了}\BibleRef{格前 9:27}；在弟茂德的事上，他更不会畏惧或羞愧。他呼求圣父和圣子作证。但为什么是蒙拣选的天使？出于极大的谦逊，正如梅瑟所说：\BibleQuote{我呼天唤地向你作证}\BibleRef{申 4:26}；又说：\BibleQuote{山岭啊，听上主的申诉！大地的基础啊，侧耳倾听！}\BibleRef{米 6:2} 他呼求圣父和圣子为他所说的作证，向祂们呼吁，针对那将来的日子，这样若有什么不当做的事，他就与此罪无关。\par

\Quote{你要遵守这些事，不可存成见，行事也不可有偏心。} 也就是说，你要公平平等地对待那些受审和被审判的人，免得有人先入为主，或事先将你拉拢到他那边。\par

但谁是蒙拣选的天使？因为有些不是蒙拣选的。正如雅各伯呼求天主和石堆作证\BibleRef{创 31:45}，我们也常同时呼求上等和下等的人作证；见证是何等重要的事。就好像他说：我呼求天主和祂的圣子及祂的仆人们作证，我已经告诫了你；所以我在他们面前告诫你。他使弟茂德心中恐惧；之后他又加上对教会维持最为重要、最为关键的，就是关于圣秩授受的事。他说：\Quote{不可轻易覆手，也不要分担保别人的罪。} 什么是\Quote{轻易}？不是在第一、第二、第三次考验之后，而是在多次严格审查和谨慎之后。因为这不是一件寻常危险的事。因为你要对他所犯的罪负责，包括他过去的和未来的罪，因为你将这份权力托付给了他。如果你不当宽恕过去的罪，你也要为未来的罪负责，因为你使他处于那个职位，成了这些罪的起因，而且你也没有让他哀叹自责、悔改。因为你既分享他的善行，也分享他的罪。\par

\Quote{保守自己纯洁。} 这是指着贞洁说的。\par

\BibleQuote{不要再只喝水，但要稍微用点酒，为你的胃病和你屡次的软弱。} 如果一个如此斋戒、只喝水以至于患上\Quote{软弱}和\Quote{屡次的软弱}的人，这样被命令要节制，并且不拒绝劝告，那么我们受到别人的劝告时更不该生气。但是为什么保禄不恢复他胃的力量？不是因为他不能——因为他的衣服能复活死人，显然他也能做到——而是因为他有重要的计划不给予这种帮助。那么他的目的是什么？使我们现在若看见伟大的有德之人患病，也不致生气，因为这是有益的管教。如果保禄本人也被派了一个\BibleQuote{撒殚的使者}打他，免得他\BibleQuote{过于高傲}\BibleRef{格后 12:11}，那么弟茂德更可能如此。因为他所行的奇迹足以使他骄傲。因此他受制于医药的规律，好使他谦卑，别人也不致生气，而能知道行这些卓越事迹的人是与他们同性情的人。在其他方面，弟茂德似乎也常患病，从\Quote{你屡次的软弱}这句话可以看出来，不仅是胃，其他部分也是如此。但他并不允许他纵酒，只允许为健康而非为享乐的量。\par

第24节\BibleRef{弟前 5:24}：\BibleQuote{有些人的罪是明显的，先被带到审判前；也有些是随后跟来的。}\par

在讲到圣秩时，他曾说：\Quote{不要分担保别人的罪。} 但他或许会说，如果我不知情呢？为什么，\BibleQuote{有些人的罪是明显的，先被带到审判前；也有些是随后跟来的。} 他意思是，有些罪是明显的，因为它们先出现；而另一些是隐藏的，因为它们随后出现。\par

第25节\BibleRef{弟前 5:25}：\BibleQuote{同样，善行也是明显的；那些不明显的，也不能隐藏。}\par

第六章第1节\BibleRef{弟前 6:1}：\BibleQuote{凡在轭下作奴隶的，应认为自己的主人配受各种尊敬，以免天主的名号和祂的道理被人亵渎。}\par

他叫他们认为主人\Quote{配受各种尊敬}；不要以为因为你是信徒，你就是自由人；因为你的自由是更忠心地服事。如果外教人看见奴隶因信而傲慢，就会亵渎，好像这道理造成不顺服。但当他看见他们顺从时，就会更愿意相信，并更听从我们的话。但天主和我们所传的福音会因他们的不顺服而被亵渎。但如果他们自己的主人是外教人呢？即使在那种情况下，他们也该为天主的名而服从。\par

\Quote{凡有信主人的仆人，不可因他们是弟兄而轻视他们，反而要更加服事他们，因为他们是信德和蒙爱的、共享恩惠的人。} \BibleRef{弟前 6:2}\par

好像他说：如果你蒙受了如此大的恩惠，竟有主人作你的弟兄，你更当因此顺服。\par

\Quote{提前审判。}这句话的意思是：恶行的结局，有的在此世被遮掩，有的未被遮掩；但在那里，善与恶都不能隐藏。所谓\Quote{提前审判}是什么意思呢？就是当人所犯的罪已经定了他的罪，或当他怙恶不悛，有人想纠正他却不能成功时。那么，这有什么用处呢？因为即使有人在此世逃避了惩罚，将来也逃不掉。那里一切都将被揭露；这对行善者是最大的安慰。\par

然后，因为他曾说：\Quote{不可偏私行事，} 好像必须加以解释，他补充说：\Quote{凡在轭下的仆人。} 但你要说：这与主教有何相干？当然有关，因为他的职责也包括劝诫和教导这些人。在此他为他们制定了完美的规则。我们看到他在各处更命令仆人，而非他们的主人，向他们指明顺服之道，并以极大尊重对待他们。因此他劝勉他们要极其温良地顺服。但他劝勉主人要戒除恐吓。他说：\Quote{不可恐吓。} \BibleRef{弗 6:9} 他为何如此命令？在不信者的情况下，自然是因为向不听从他人说话不合情理；但关于信者，他为何如此？因为主人在仆人身上所施的恩惠，多于仆人对主人的贡献。因为主人出钱购买足够的食物和衣物，并在其他方面对他们付出许多照顾，所以主人实际上服事了仆人更多，这也是此处所暗示的，当他说：\Quote{他们是信德的、蒙爱的、共享恩惠的。} 他们为你们的安息付出许多辛劳和麻烦，难道不应该从仆人那里得到许多尊荣吗？\par

\Emph{道德教训。} 如果他劝勉仆人如此绝对服从，试想我们应当如何对待我们的主——祂从无中创造了我们，又养活并穿戴我们？如果不能以其他方式，至少让我们像仆人服事我们那样服事祂。难道他们不是安排一生为给主人安歇，他们的工作和生命岂不是照顾主人的事务？他们终日忙于主人的工作，只有很少一部分时间为自己忙碌？但我们却相反，常忙于自己的事，几乎不为我们的主做事，尽管祂不需要我们的服事，就如主人需要仆人的服事一样，那些服事反使我们自己受益。就他们而言，主人因仆人的服务获益；但就我们而言，仆人的服务并非使主获益，反而有益于仆人。正如圣咏作者所说的：\Quote{我的善行对祢毫无益处。} \BibleRef{咏 15:2} 你说，我正义对天主有何益处？我不义对祂又有何害？祂的本性是不朽的、不能受损的、超越一切苦难的。仆人一无所有，一切都是主人的，无论他们多么富有。但我们却有许多属于自己的东西。\par

并且我们不仅享受来自宇宙之王的如此大的尊荣。有哪个主人曾为自己的仆人献出自己的儿子？没有一个人，所有人都宁愿为儿子献出仆人。此处却相反：\Quote{祂没有怜惜自己的儿子，反而为我们众人把祂交出，} 为祂的仇敌，那些恨祂的人。仆人虽然被要求做极苦的劳役，却不急躁——至少那些好仆人不会。但我们却多少次发出怨言？主人应许仆人的，没有一件像天主应许我们的那样；他们应许的不过是现世的自由，这往往比奴役更糟，因为自由常因饥荒比奴役本身更痛苦。但这却是他们最大的恩惠。但在天主那里，没有一样是暂时的、会朽坏的；而是什么？你愿意知道吗？请听：\Quote{以后我不再称你们为仆人，你们是我的朋友。} \BibleRef{若 15:13}\par

亲爱的，让我们羞愧，让我们敬畏。让我们至少像我们的仆人服事我们那样服事我们的主。甚至我们连最小的服事也没有做到！迫使他们成为哲学家的，是不得已。他们只有食物和住所；我们却拥有许多，又期望更多，却以奢侈侮辱我们的恩主。如果别无他法，至少从他们身上学习哲学的准则。圣经惯常不断派遣人甚至向无理性的受造物学习，比如当它命令我们效法蜜蜂和蚂蚁时。但我劝你们只效法仆人：他们因害怕主人所做的，让我们因敬畏天主也照样做；因为我发现你们连这点也没有做到。他们因害怕我们而忍受许多侮辱，并以哲学家的忍耐默默承受。无论有理无理，他们遭受我们的暴力，却不反抗，反而恳求我们，虽然他们常常没有做错。他们满足于得着不多，常常还不敷所需；睡在草铺上，只吃面包，并不抱怨或嘟囔艰苦的生活，但因怕我们而克制了急躁。当他们受托钱财时，他们全部归还。我不是说那卑劣的，而是指中等好人。如果我们恐吓他们，他们立即畏惧。\par

这岂非哲学么？不要说他们受必然之束缚，因你自己亦受地狱恐惧之必然。然而你并未学到智慧，也没有将你所从仆人那里得到的同等尊荣归给天主。你的每个仆人都按你的规矩分得住所，他不侵犯邻舍的，也不因贪多而加害于人。这种忍耐因主人之恐惧而强加于仆役之间，你很少见到一个仆人抢劫或伤害同伴。但在自由人中却恰恰相反。我们彼此咬噬、吞食。我们不畏惧我们的主：我们抢劫掠夺我们的同伴，就在他眼前击打他们。仆人不会这样做；他若击打，必不在主人看见之时；他若辱骂，必不在主人听见之时。但我们却敢做任何事，尽管天主看见并听见一切。\par

主人的恐惧常在他们眼前，我们主的恐惧却从不在我们眼前。因此一切秩序颠倒，因此一切混乱毁坏。我们从不考虑自己所犯的过犯，但若仆人们做错，我们便严苛地追究每一项，就连最小的过失也不放过。我说这些并非要使仆人懈怠，而是要责备我们的麻木，将我们从懒散中唤醒，好让我们以仆人服事主人的同样热忱服事我们的天主；以受造物服事我们的同样忠诚服事我们的造物主，尽管他们从未从我们这里领受过这样的恩赐。因为他们本性也是自由的。对他们也说：\Quote{让他们管理鱼类。}\BibleRef{创 1:26}因为这种奴役并非出于本性，而是出于某种特殊原因或环境的结果。然而尽管如此，他们仍大大尊敬我们；我们极严格地要求他们的服事，而向天主我们几乎不奉献最小的部分，尽管这益处终将归于我们自己。因为我们越热忱地服事天主，我们得的益处就越大。所以我们不要剥夺自己如此重大的利益。因为天主自给自足，一无所缺；赏报和利益完全归于我们。因此，我恳求你们，要这样被感动，如同不是服事天主，而是服事我们自己，且要怀着恐惧和战栗服事他，好使我们获得所应许的福分，借着我们的主耶稣基督，愿光荣归于他，直到永远。阿们。\par

\SourceRecord{Translated by Philip Schaff.；Nicene and Post-Nicene Fathers, First Series；Vol. 13.；Edited by Philip Schaff.；Buffalo, NY: Christian Literature Publishing Co.,；1889.；<http://www.newadvent.org/fathers/230616.htm>.}

% structure: p0001 source=h1 target=section reason=page-title

\section{论弟茂德前书 第十七篇讲道}

\Emph{\BibleRef{弟前 6:2-7}}\par

\Emph{\Quote{这些事你要教导和劝勉。若有人讲异端，不听从我们主耶稣基督的健全言词，与那合乎虔敬的道理，他是骄傲的，一无所知，反而为争辩和言词上的角逐而痴迷，由此产生嫉妒、争吵、辱骂、邪恶的猜疑，以及心术不正之徒的谬论；这些人丧失了真理，以为虔敬便是获利。你应从这样的人中退避。但是虔敬与知足，才是大利。因为我们没有带什么到世上来，我们确实也不能带什么去。}}\par

教师不仅需要权威，也需要温和；不仅需要温和，也需要权威。蒙福的保禄教导这一切，他有时说：\Quote{这些事你要命令和教导}\BibleRef{弟前4:11}；有时又说：\Quote{这些事你要教导和劝勉。}如果医生恳求病人，不是为了自己的健康，而是为了减轻他们的病痛，恢复他们衰弱的体力，那么我们在教导人时更应当采用这种恳求的方法。因为蒙福的保禄不拒绝做他们的仆人。他说：\Quote{我们不是宣讲我们自己，而是宣讲基督耶稣为主，并且我们因耶稣的缘故做你们的仆人}\BibleRef{格后4:5}；又说：\Quote{一切都属于你们，无论是保禄，或是阿颇罗}\BibleRef{格前3:12}。他在这项服事中愉快地服务，因为这不是奴役，而是超越自由。因为主说：\Quote{凡犯罪的，就是罪的奴隶}\BibleRef{若 8:34}。\par

\Quote{若有人讲异端，不听从我们主耶稣基督的健全言词，与那合乎虔敬的道理，他是骄傲的，一无所知。}因此，骄傲不是来自知识，而是来自\Quote{一无所知}。因为知道虔敬道理的人也是最倾向节制的。知道健全言词的人，就不歪曲。因为身体的炎症犹如灵魂的骄傲。正如我们不会说发炎的部位是健康的，同样我们也不认为傲慢的人是健康的。因此，可能是有知识而仍然一无所知。因为不知道应当知道的事，就是一无所知。骄傲来自一无所知，这从以下事实可以证明。基督\Quote{使自己空虚}\BibleRef{斐2:7}，因此知道这事的人就不会自高自大。人除了从天主领受的以外，一无所有，因此他不会自高自大。\Quote{你有什么不是领受的呢？}\BibleRef{格前4:7}。他洗了门徒的脚，知道这事的人怎能自高自大呢？因此他说：\Quote{你们做完了一切，要说是无用的仆人}\BibleRef{路17:10}。税吏只因他的谦卑而被接纳，法利塞人却因自夸而灭亡。骄傲的人不知道这些事。再者，基督自己说：\Quote{我若说得不对，你指证不对的地方；若说得对，你为什么打我呢？}\BibleRef{若18:23}。\par

\Quote{痴迷于争辩}。因此提问就是痴迷。\Quote{言词上的角逐}，这说得对。因为当灵魂因推理而发热、骚动时，它才提问；但灵魂健康时，它不提问，而是接受信仰。但从提问和言词上的角逐中，什么也发现不了。因为当信仰只允许的事物被好问的精神所接受时，它既不能证明它们，也不能让我们理解它们。如果一个人闭上眼睛，他就无法找到他所寻找的任何东西；或者，如果他睁开眼睛，却把自己埋起来，遮住太阳，他就无法找到任何东西。同样，没有信仰，什么也辨别不出，反而必然产生争论。\Quote{由此产生辱骂、邪恶的猜疑}，也就是说，由提问而产生的错误意见和道理。因为当我们开始提问时，就会对天主臆断一些不当的事。\par

\Quote{谬论}，就是闲谈或交谈，或者他可能指互相交流，并且如同病羊通过接触将疾病传染给健康的羊，这些恶人也一样。\par

\Quote{丧失了真理，以为虔敬便是获利。}看，言词上的角逐产生了多少罪恶！贪财、无知和骄傲；因为骄傲是由无知产生的。\par

\Quote{你应从这样的人中退避。}他并没有说，和他们争论或对抗，而是\Quote{退避}，远离他们；正如他在别处说：\Quote{对异端的人，警告一两次后，就该弃绝他}\BibleRef{铎3:10}。这说明他们不是由于无知而犯错，而是将自己的无知归因于懈怠。你永远无法说服那些为金钱而争论的人。只有在你给与的时候，他们才可能被说服，即便如此，你也永远无法满足他们的欲望。因为经上说：\Quote{贪婪人的眼目总是不满足于一分}\BibleRef{德14:9}。因此，对于这样无可救药的人，远离他们是正确的。如果那位有义务为真理作战的人，也被劝告不与这样的人争论，那么作为门徒们，我们更应该避免。\par

他（保禄）说\Quote{他们认为虔敬是获利的途径}之后，又补充说：\Quote{但是虔敬与知足，才是大利}，不是指拥有财富，而是指没有财富。为了让人不因贫穷而沮丧，他鼓励并振奋他们的精神。他说，他们认为虔敬是获利的途径，事实也确实如此；不过不是他们所认为的方式，而是更高尚的方式。在推翻了他们的观点之后，他赞扬了另一种。因为世俗之利是虚无的，这很明显，因为它会被留下，在我们离开时不会跟着我们或与我们同行。这从何而知？因为我们来到这个世界时一无所有，所以我们离开时也一无所有。因为本性赤身来到世界，也必赤身离去。因此我们不需要多余之物；如果我们没有带来什么，也不能带走什么。\par

第八节：\Quote{只要有衣有食，就当知足。}\BibleRef{弟前 6:8}\par

我们应当只吃足以滋养我们的食物，只穿足以遮盖我们、包裹我们赤裸的衣服，此外无需更多；普通衣物便能满足此目的。接着，他从今世之事劝勉说：\par

第九节：\Quote{但那些想要发财的人}——并非那些富有的人，而是那些想要发财的人。因为一个人可以拥有钱财并善用它，不将其估价过高，而是施舍给穷人。因此，他所责备的并非这种人，而是贪财之人。\BibleRef{弟前 6:9}\par

\Quote{想要发财的人，就陷于诱惑、罗网，及许多无知而有害的私欲，这使人沉溺于败坏和灭亡之中。}\par

他公正地使用了\Quote{沉溺}一词，因为人无法从那深渊中浮起。\Quote{在败坏和灭亡之中。}\par

第十节：\Quote{因为贪财是万恶之根；有些人因贪求钱财，就离弃了信德，并使自己受了许多刺心的痛苦。}\BibleRef{弟前 6:10}\par

他提到两件事，而将他们可能视为更沉重的那件——\Quote{许多刺心的痛苦}——放在最后。要知此言何等真实，只有与富人同住，看看他们有多少忧伤、多少苦涩的抱怨，方能领会。\par

第十一节：\Quote{但你，属天主的人啊。}\BibleRef{弟前 6:11}\par

这是一个极尊贵的称号。因为所有人都是属天主的人，但义人尤其如此，不单因受造的权利，更因归属的权利。若你真是\Quote{属天主的人}，就不要寻求那些多余之物，它们不引领你归向天主，反而……\par

\Quote{你要避免这些事，而追求正义。}两个措辞都强调：他并没有说离开一件而靠近另一件，而是说\Quote{避免这些事，追求正义}，以免你成为贪婪之人。\par

\Quote{虔敬}，即道理纯正。\par

\Quote{信德}，与辩难相对。\par

\Quote{爱德}，忍耐，温和。\par

第十二节：\Quote{你要为信德打那美好的仗，持定永生。}看，这就是你的赏报，\Quote{你也是为此被召，并在许多见证人面前宣认了那美好的宣认}，为着永生的希望。\BibleRef{弟前 6:12}\par

这就是说，不要使那信德蒙羞。你为何徒劳无功？但他说那些想要发财的人所陷入的\Quote{诱惑和罗网}是什么？这使人离弃信德，陷入危险，变得胆怯。\Quote{无知的欲望}，他说。这岂非无知，当人们喜欢豢养白痴和侏儒，并非出于仁爱，而是为了取乐；当他们在厅堂中设鱼池，饲养野兽，将时间花在狗身上，装扮马匹，视之如子？所有这些事都是无知而多余的，全无必要，全无用处。\par

\Quote{无知而有害的私欲！}什么是有害的私欲？当人生活不法，贪求邻人之物，竭力奢侈，渴望酗酒，图谋杀人害命。由这些私欲，许多人图谋暴政而灭亡。无疑，以此等念头劳碌，既是无知，也是有害。他恰当地说：\Quote{他们离弃了信德。}因为贪财吸引他们的眼目，逐渐盗走他们的心智，使他们看不见道路。正如一个人走在直路上，却心不在焉，他固然在行路，却常常不知不觉地路过他原本要前往的城邑，双脚漫无目的地乱走：贪财正是如此。\Quote{他们使自己受了许多刺心的痛苦。}你明白他用\Quote{刺心}一词的含义吗？他以此比喻要表达的是：私欲如同荆棘，当人触碰荆棘，就会刺伤手，留下伤口；同样，陷入这些私欲的人必为所伤，用痛苦刺透自己的灵魂。那些被如此刺透的人所承受的焦虑和烦恼，无法言表。因此他说：\Quote{你要避免这些事，而追求正义、虔敬、信德、爱德、忍耐、温和。}因为温和源于爱德。\par

第十二节：\Quote{你要打那美好的仗。}\BibleRef{弟前 6:12}\par

这里他称赞弟茂德的勇敢和刚毅，因他在众人面前坦然\Quote{作了宣认}，并提醒他早年的教训。\par

\Quote{持定永生。}\par

不仅需要宣认，还需要忍耐以持守那宣认，以及激烈的争战和无数辛劳，免得你被击败。因为绊脚石和阻碍很多，所以\Quote{路是窄而狭的}\BibleRef{玛 7:14}。因此必须收敛心神，四面束紧。周围处处有享乐吸引灵魂的眼目：美貌、财富、奢侈、安逸、荣耀、报复、权势、统治的享乐，这些外表美丽可爱，足以俘获不坚定、不爱真理的人。因为真理的面容严苛而不诱人。为何如此？因为她所应许的福乐都是将来的，而别的享乐却提供现世的尊荣、愉悦和安逸——尽管都是虚假和伪造的。因此，粗俗、柔弱、怯懦、不愿为德行辛劳的心灵依附后者。正如在外教人的竞技中，不热切渴望冠冕的人从一开始就能沉溺于宴饮和醉酒，实际上懦弱怯战的斗士正是如此；而那些坚定注视冠冕的人，虽遭受无数打击，却被未来奖赏的希望所支撑和激励。\par

\Emph{道德训诲。}因此，让我们逃离这万恶的根源，我们便会摆脱一切罪恶。他说：\Quote{贪爱钱财}是\Quote{万恶的根源}；这是\Person{保禄}说的，更确切地说，是\Person{基督}藉\Person{保禄}说的。让我们看看为何如此。世界的实际经验也为此作证。因为有什么祸患不是由财富引起的呢？或者更准确地说，不是由财富本身，而是由那些不知如何善用财富之人的邪恶意志引起的？因为行善可以使用财富，甚至藉此继承天国。然而，那赐予我们用于救济穷人、弥补我们过去的罪过、赢得美名、并悦乐天主的财富，如今却被我们用来反对穷人和可怜人，甚至反对我们的灵魂，并极其触怒天主。因为对于他人，一个人夺去了他的财富，使他陷入贫困，却使自己陷入死亡；使他人在这里忍受贫乏之苦，却使自己在永罚中受煎熬。你以为他们所受的痛苦相等吗？\par

那么，它没有引发什么祸患呢？什么欺诈行为，什么抢劫！什么痛苦、仇恨、争斗、战争！它岂非伸手触及死者，甚至父亲和兄弟吗？那些被这种激情所控制的人，岂非违犯自然律和天主的诫命吗？简言之，一切不都是如此吗？岂不是这需要我们的法庭吗？因此，消除贪爱钱财，你就终止了战争、战斗、仇恨、纷争和冲突。这样的人理应像豺狼和害虫一样被逐出世界。因为正如逆风强风，扫过平静的大海，从深处搅动它，将深海的沙与上面的波浪混在一起，同样，贪财之人搅乱和动摇一切。贪财之人从不认识朋友：朋友，我说了吗？他甚至不认识天主自己，因为他被贪婪的激情逼得疯狂。你岂不见那些执剑出征的巨人吗？这是疯狂的写照。但贪财之人并不伪装，他们是真疯狂，神志不清；如果你能剖开他们的灵魂，你会发现他们不是佩着一两把剑，而是成千上万，不承认任何人，向所有人发泄怒气；飞扑向所有人，向所有人咆哮，宰杀的不是狗，而是人的灵魂，并向上天本身发出亵渎。被他们对财富的狂热所驱使，这些人颠覆和毁灭一切。\par

因为我不知道该控告谁！这是一种瘟疫，它抓住所有人，有的多些，有的少些，但所有人在某种程度上都罹患此病。如同烈火燃着林木，荒废并摧毁周围的一切，这种激情已经蹂躏了世界。君王、官吏、平民、穷人、女人、男人、儿童，都同样受影响。好像一层浓密的黑暗笼罩大地，没有人头脑清醒。然而，我们在公共场合和私下都听到许多反对贪婪的演说，但没有人因此改善。\par

那么该怎么办？我们如何熄灭这火焰？虽然它已经上升到天上，但它是可以被熄灭的。我们只要愿意，就能控制这大火。因为正如因我们的意志它得以抬头，同样因我们的意志它也可以被压制。难道不是我们自己的选择引起了它，难道同样的选择不能熄灭它吗？只要让我们愿意。但那愿望如何产生呢？如果我们思考财富的虚空和无益，它不能从这里与我们同去，就是在这里也遗弃我们，而它留下时，却给我们留下与我们同去的创伤。如果我们看到那里有财富，与此世的财富相比，此世的财富比粪土更可鄙。如果我们思考它伴随着无数的危险，暂时的快乐，快乐中掺杂着悲伤。如果我们正确默观永生的真财富，我们就能轻视世俗的财富。如果我们记住它既无益于荣耀，也无益于健康，或任何其他事情；反而使人沉沦于毁灭和丧亡。如果你思考你在这里富有，有许多人服从你；但当你离开这里时，你将赤身而去，孤身一人。如果我们常常向自己呈现这些事，并从他人那里听这些事，或许会有回归健全心智的时候，并从这可怕的惩罚中获得解救。\par

珍珠难道美丽吗？然而试想，它不过是海水，曾被抛入深渊的怀抱。金银难道美丽吗？它们不过是尘土灰烬。丝绸衣裳难道美丽吗？它们不过是虫子的纺丝。这种美丽仅存于意见，存于人类的偏见，而非事物的本性。因为那在本性上拥有美丽的，无需任何人指明。你若看到一枚镀金的铜币，起初会欣赏它，以为它是金子；但一旦通晓者向你揭示其伪，你的惊叹便随欺骗消失。因此，美丽并不在事物的本性中。银子也是如此；你可能会把锡当作银子欣赏，就如把铜当作金子欣赏，你需要有人指点你该欣赏什么。如此，我们的眼睛不足以辨别差异。对于花卉，它们远为美丽，情况却非如此。你若看到一朵玫瑰，无需任何人指点，你自能区别银莲花、紫罗兰或百合，以及其他各种花卉。因此，这不过是偏见。为显明这毁灭性的情感不过是偏见，请告诉我：倘若皇帝降旨，规定银子比金子更贵重，难道你不会将你的喜爱和羡慕转向银子吗？可见，我们处处受贪婪和意见的支配。为证明这点，以及事物因其稀有而非本性被珍视，可见如下事实：在我们这里被视为低贱的果品，在卡帕多西亚人那里备受推崇，在赛里斯人那里甚至比我们这里最珍贵的果品更有价值——这些衣裳便来自那个国度；在出产香料和宝石的阿拉伯和印度，还可举出许多类似例子。因此，这种偏好不过是偏见和人的意见。我们不依判断行事，而是随意而行，为偶然所决定。但让我们从这种醉态中苏醒，让我们定睛于那真正美丽的、本性美丽的，即虔敬和义德；好使我们获得所应许的福分，赖我们的主耶稣基督的恩宠和慈爱，愿光荣归于祂，等等。\par

\SourceRecord{Translated by Philip Schaff.；Nicene and Post-Nicene Fathers, First Series；Vol. 13.；Edited by Philip Schaff.；Buffalo, NY: Christian Literature Publishing Co.,；1889.；<http://www.newadvent.org/fathers/230617.htm>.}

% structure: p0001 source=h1 target=section reason=page-title

\section{关于弟茂德前书第十八篇讲道}

\Emph{弟前 6:13-16}\BibleRef{弟前 6:13-16}\par

\BibleQuote{我在那赋予万有生命的天主面前，并在那在本丢比拉多前宣过美好宣认的基督耶稣面前嘱咐你：务要持守这命令，毫无玷污，无可指摘，直到我们的主耶稣基督显现；到了时机，那当受称颂、独一无二的全能者、万王之王、万主之主，必会将它显明出来；祂独享不朽，住在不可接近的光中，没有人见过祂，也不能看见祂；愿尊贵和永远的权能归于祂。阿们。}\par

他再次呼求天主作证，如同稍早所作的一样，这既为增加弟茂德的敬畏，也为保障他的安全，并表明这些并非人的诫命，使他既从主亲自领受这命令，又时刻铭记他所听见的见证者，因而将这话更深刻地印在心上。\par

\Quote{我嘱咐你，}他说，\Quote{在那赋予万有生命的天主面前。}\par

这里既给予他面对危险时的安慰，又唤醒他记起复活。\par

\Quote{并在耶稣基督面前，祂在本丢比拉多前宣过美好宣认。}\par

劝诫又来自于祂师傅的榜样，他的意思是：祂怎样做了，你也应该照样做，因为祂为此而作了\Quote{宣认}\BibleRef{伯前 2:21}，好使我们追随祂的足迹。\par

\Quote{美好的宣认。}他在致希伯来人书信中所做的——\Quote{仰望信德的创始者和完成者耶稣，祂因那摆在面前的喜乐，轻视了羞辱，忍受了十字架，如今坐在天主宝座的右边；你们要仔细想想那位忍受了罪人这样反对的，免得你们灰心丧志}\BibleRef{希 12:2}——如今他同样对待他的弟子弟茂德。仿佛在说：不要畏惧死亡，因为你所事奉的天主是能赐予万有生命的。\par

但他所指的是哪一次\Quote{美好的宣认}？是当本丢比拉多问\Quote{你是君王吗？}祂说\Quote{我为此而生}\InnerQuote{，又说}\Quote{我来是为给真理作证；看，这些人听见了我。}\BibleRef{若 18:37}或者当被问\Quote{你是天主子吗？}祂回答\Quote{你们说我是（天主子）。}\BibleRef{路 22:70}他或许指此，或许指其他许多祂所作的证词和宣认。\par

第14节。\BibleQuote{务要持守这命令，毫无玷污，无可指摘，直到我们的主耶稣基督显现。}\BibleRef{弟前 6:14}\par

也就是说，直到你的结局，即你离开人世，虽然他并不这么表达，而是为了更多地激发他，说\Quote{直到祂显现。}但何为\Quote{持守这命令毫无玷污}？就是不沾染任何污秽，无论是教义还是生活上的污秽。\par

第15节。\BibleQuote{到了时机，那当受称颂、独一无二的全能者、万王之王、万主之主，独享不朽，住在不可接近的光中，必会将它显明出来。}\BibleRef{弟前 6:15}\par

这些话是指谁说的？指圣父还是指圣子？无疑是指圣子；这样说为安慰弟茂德，使他不要害怕或畏惧地上的君王。\par

\Quote{到了时机}，即适当且相宜的时候，使他不要因尚未到来而不耐烦。从哪里能显明祂会显出这个时机呢？因为祂是\Quote{独一无二的全能者}。那么，祂必会显出它，祂是\Quote{当受称颂的}，不，祂就是福乐本身；这样说为了表明在祂的显现中没有任何痛苦或困扰。\par

他说\Quote{独一无二的}，或是与世人相对而言，或是因祂无始，或是如同我们称赞某个人时所常说的。\par

\Quote{独享不朽。}那么，圣子没有不朽吗？祂不就是不朽本身吗？与圣父同一本体的圣子，怎么会没有不朽呢？\par

\Quote{住在不可接近的光中。}那么，祂自己是光，还有另一个光祂住在其中？祂岂是被地方所限制？切勿这样想。这句话描述了天主性体的不可理解性。他这样谈论天主，是在他所能的最好方式下说出的。请注意，舌头当要说出什么伟大的事时，它便力量不足。\par

\Quote{没有人见过祂，也不能看见祂。}诚然，没有人见过圣子，也不能看见祂。\par

\Quote{愿尊贵和永远的权能归于祂。阿们。}他如此恰当且有意义地谈论了天主。因为他曾呼求祂作证，于是多讲述这位见证者，使他的门徒更加敬畏。他用这些话将光荣归于祂，我们所能做或所说的，也就仅此而已。我们不该过分好奇地探究祂是谁。如果永远的权能属于祂，就不要害怕。即使现在尚未显现，尊贵属于祂，权能永远属于祂。\par

第17节。\BibleQuote{你要嘱咐今世富裕的人不要心高气傲。}\BibleRef{弟前 6:17}\par

他说\Quote{今世富裕}说得好，因为另有人将来在世富裕。他提出这一劝诫，知道没有比财富更普遍地产生骄傲与傲慢的了。因此，为缓解这个，他立即加上\Quote{也不要依赖无定的财富}，因为那是骄傲的根源；凡寄望于天主的，就不会自高。你为何将希望寄托在转眼即逝的东西上？财富便是如此！你为何寄望于你所不能确信的？但你说，他们如何才能避免心高气傲？要思量财富的无定与不稳，并意识到寄望于天主远为宝贵，因为天主本身就是财富的创造者。\par

第17节。\BibleQuote{但要寄望于永生的天主，}他说，\Quote{祂把万物丰盛地赐给我们享用。}\BibleRef{弟前 6:17}\par

这\Quote{万物丰盛}说得对，是指四季的变迁、空气、光、水以及其他恩赐。这一切是何等丰盛、多么慷慨地赐予！你若寻求财富，当寻求那些稳固持久、出于善工的财富。他以下文表明这正是他的意思。\par

第 18 节。\Quote{好行善事，}他说，\Quote{富有善工，乐于施舍，愿意与人分享。}\BibleRef{弟前 6:18}\par

第一句指财富，第二句指爱德。因为愿意与人分享，意味着他们是合群而善良的。\par

第 19 节。\Quote{为自己积蓄美好的根基，以备将来。}\BibleRef{弟前 6:19}\par

那里没有不确定的事，因为根基稳固，没有动摇；一切都是坚固、固定、不可移动、稳妥、持久的。\par

第 19 节。\Quote{使他们把握，}他说，\Quote{永生。}\BibleRef{弟前 6:19}\par

因为行善工能确保享受永生。\par

第 20 节。\Quote{哦，弟茂德，要保守所受的托付。}\BibleRef{弟前 6:20}\par

不可让它减少。这不是你的。你受托管理他人的财产，不可将其减损。\par

第 20 节。\Quote{躲避亵渎的空谈，以及伪称为知识的对抗。}\BibleRef{弟前 6:20}\par

他这样称呼得很好。因为没有信德，就没有真知；凡是出于我们推理的，都不是知识。或者他说这话，是因为那时有些人自称\OriginalTerm{诺斯替派信徒}{Gnostics}，以为比别人知道得更多。\par

第 21 节。\Quote{有些人自称有这知识，却在信德上偏离了。}\BibleRef{弟前 6:21}\par

你看，他再次命令弟茂德甚至不要与他们来往。\Quote{躲避对抗。}因此，有些对抗我们不值得回答，因为它们使人离开信德，不容人坚定稳固于其中。我们不要追求这学问，而要持守信德，那不可动摇的磐石。因为无论是洪水还是狂风都不能伤害我们，因为我们立于磐石上坚定不移。同样，在此生中，如果我们选择那真正是根基的那一位，我们就站稳了，没有伤害能袭击我们。因为选择了来生的财富、尊荣、光荣、快乐的人，有什么能伤害他呢？这一切都是稳固的，其中没有变化；这里的一切都易逆转，且永远变化。你想要什么呢？光荣？圣咏作者说：\BibleQuote{他的荣耀不会随他下去。}\BibleRef{咏 48:17}而且，它往往在他活着时也不停留。但德行却非如此，所有属于德行的事都是持久的。在这里，从职位获得光荣的人，一旦别人接替他的职位，他就成为平民而无荣。富人因强盗袭击或谄媚者和恶棍的陷害而变穷。基督徒却不是这样。节制的人，若谨慎自守，就不会被剥夺其德行。自制的人，不会沦为平常人和臣服者。\par

而且会看出这统治高于任何其他。因为，请告诉我，统治我们的人类同胞却做自己情欲的奴隶，有什么益处呢？或者，如果我们超脱了情欲的暴政，即使无人受我们统治，我们又有什么损失呢？那才是自由，才是统治，才是王权和主权。相反，即使一个人戴上无数王冠，却是奴隶。因为当众多主人从内部支配他——贪财、享乐、愤怒及其他情欲——他的王冠又有何用？那些情欲的暴政更加严厉，因为即使他的王冠也不能救他脱离其奴役。正如一个人曾是国王，后被蛮族贬为奴隶，蛮族为了更绝对地显示其权力，不但不剥夺他的紫袍和王冠，反而强迫他穿着这些做各种卑贱的事，挑水、烧饭，使他的耻辱和他们的光荣更加明显：我们的情欲正是以这样的方式比任何蛮族更野蛮地统治我们。因为蔑视情欲的人也能蔑视蛮族；但屈服于情欲的人，将遭受比蛮族更严厉的痛苦。蛮族当权时，可能折磨身体，但这些情欲折磨灵魂，将其全面撕裂。蛮族得胜时，使人暂时死亡，但这些情欲导致将来的死亡。所以，只有那内心自由的人才是自由人；屈服于这些非理性情欲的人，是奴隶。\par

没有哪个主人，无论多么残忍，会下达如此严厉且不人道的命令。他们实际上对他说：\Quote{无端无限地玷辱你的灵魂——得罪你的天主——对自然的要求充耳不闻；即使是你父亲或母亲，也不要羞于与他们对抗。}贪财的命令就是这样。\Quote{她说：向我祭献，不是牛犊，而是人。}先知确实说：\BibleQuote{献人吧，因为牛犊已经缺乏。}\BibleRef{欧 13:2}但贪财说：\Quote{即使还有牛犊，也要献人。献上那些从未伤害过你的人，不仅如此，还要杀害他们，即使他们是你的恩人。}或者说：\Quote{要争战，要成为一切人、自然本身和天主的公敌。积聚金子，不是为了享受，而是为了守住它，给自己带来更大的折磨。}因为贪财者不可能享受它，他害怕金子减少，库存耗尽。\Quote{要警惕，}它说，\Quote{要怀疑每一个人，甚至家人和朋友。垂涎他人的财物。虽然你看见穷人饿死，也不要给他任何东西；如果可能，还要剥他的皮。违背誓言，撒谎，起假誓。做控告者，做假证人。如果必要，不要拒绝赴火，经受千百次死亡，饿死，与疾病搏斗。}难道不是贪财强加这些律法吗？\Quote{要粗鲁无礼，厚颜无耻，凶恶邪恶，忘恩负义，冷酷无情，不友善，不忠诚，没有爱心，弑亲，像野兽而不像人。在狠毒上超过蛇，在贪婪上超过狼。在残忍上甚至超过野兽，是的，即使必须达到魔鬼的恶毒，也不拒绝。在你的恩人面前做陌生人。}\par

贪婪不是说这一切，难道就没有人听吗？相反，天主却说：要作众人的朋友，要温和，为众人所爱，不无故得罪任何人。\BibleQuote{孝敬你的父亲和你的母亲}。赢得可敬的名声。不要作人，而要作天使。不说任何不正经的话，不说任何假话，甚至连想都不要想。救济穷人。不要因掠夺别人而给自己惹麻烦。不要胆大妄为，不要傲慢无礼。天主这样说，却没有人听从。那么，地狱、不灭的火和不死的虫，岂不公正地受到威胁吗？我们要这样把自己推下悬崖多久？我们要在荆棘上行走，用钉子刺伤自己，还为此感恩多久？我们把自己置于残忍的暴君之下，却拒绝温和的主，祂不施加任何痛苦、野蛮、沉重或无益的事，而是一切都有用、有价值、有益处。那么，让我们振作起来，收敛心神，集合力量。让我们按我们所应当的，爱天主，好使我们获得祂为爱祂的人所预许的福分，借着我主耶稣基督的恩宠和慈爱，愿光荣归于祂及圣父，及圣神，永生永王。\par

\SourceRecord{Translated by Philip Schaff.；Nicene and Post-Nicene Fathers, First Series；Vol. 13.；Edited by Philip Schaff.；Buffalo, NY: Christian Literature Publishing Co.,；1889.；<http://www.newadvent.org/fathers/230618.htm>.}
